% !TEX program = lualatex
% ---------------------------------------------------------------------------------------------
%  Readback note 34: Exponents of elements of Kisin modules, the cost of deepening them, and cokernel killers
%
%  Build:  python3 tools/rb.py note <slug>
%
%  One unit of a Lean artifact, rewritten as ordinary mathematics. The prose is hand-written. Lean names,
%  line ranges and counts come from data.tex; quoted Lean statements are copied from the artifact by the build.
%
%  Rules of the text: standard terminology; every symbol is explained where it first appears and listed again
%  in the notation at the end; no letter has two meanings; a proof is complete or says what it takes for granted;
%  Lean names appear only in the LEAN lines and in the concordance.
%
%  The signature rule (the build enforces it): a function, relation or predicate that the artifact defines is
%  written only through its notation macro, which the build makes from the spec (`notation`) with one slot for
%  each explicit argument of the Lean definition. No argument is fixed and dropped, also not "for brevity" and
%  not for one section. Every argument and every hypothesis of a stated declaration stands inside its
%  statement, or in a Setting that the statement cites by number; only implicit and instance arguments (the
%  ambient objects) may stand in the conventions.
% ---------------------------------------------------------------------------------------------
\documentclass[11pt,a4paper]{article}
\ifdefined\pdfvariable\pdfvariable suppressoptionalinfo \numexpr 1+2+4+8\relax\fi
\usepackage[a4paper,hmargin=35mm,top=22mm,bottom=27mm,footskip=12mm,nohead]{geometry}

\usepackage{fontspec}
\directlua{
  luaotfload.add_fallback("gfallback", {"STIXTwoMath-Regular.otf:mode=node;", "DejaVuSans.ttf:mode=node;"})
  luaotfload.add_fallback("gmonofallback", {"DejaVuSansMono.ttf:mode=node;", "STIXTwoMath-Regular.otf:mode=node;"})
}
\setmainfont{LibertinusSerif}[Extension=.otf,UprightFont=*-Regular,ItalicFont=*-Italic,
  BoldFont=*-Semibold,BoldItalicFont=*-SemiboldItalic,RawFeature={fallback=gfallback}]
\setsansfont{Inter}[Extension=.otf,UprightFont=*-Regular,ItalicFont=*-Italic,
  BoldFont=*-SemiBold,BoldItalicFont=*-SemiBoldItalic,RawFeature={fallback=gfallback}]
\setmonofont{JetBrainsMono}[Extension=.otf,UprightFont=*-Regular,ItalicFont=*-Italic,
  BoldFont=*-SemiBold,BoldItalicFont=*-SemiBoldItalic,Scale=MatchLowercase,
  RawFeature={-calt;-liga;fallback=gmonofallback}]
\usepackage{amsmath}
\usepackage{amsthm}
\usepackage[math-style=TeX]{unicode-math}
\setmathfont{LibertinusMath-Regular.otf}

\usepackage{xcolor}
\usepackage{tikz}
\usetikzlibrary{calc,arrows.meta,positioning,backgrounds}
\usepackage[protrusion=true,expansion=true]{microtype}
\usepackage{enumitem}
\usepackage{needspace}
\usepackage{booktabs}
\usepackage{array}
\usepackage{longtable}
\usepackage{fancyhdr}
\usepackage{lastpage}
\usepackage{multicol}
\usepackage[hidelinks,pdfusetitle]{hyperref}

\definecolor{ink}{HTML}{0F172A}
\definecolor{inksoft}{HTML}{475569}
\definecolor{inkfaint}{HTML}{94A3B8}
\definecolor{rulegrey}{HTML}{CBD5E1}
\definecolor{arith}{HTML}{C2410C}
\definecolor{est}{HTML}{7C3AED}
\definecolor{ana}{HTML}{0891B2}
\definecolor{ok}{HTML}{047857}
\definecolor{leancol}{HTML}{3730A3}

\title{Exponents of elements of Kisin modules, the cost of deepening them, and cokernel killers}
\author{grh-compact}
\date{}

% ------------------------------------------------------------------------------------- macros
\newcommand{\leanfont}{\ttfamily\fontsize{8.6}{10.5}\selectfont\color{leancol}}
\newcommand{\lean}[1]{\csname LN@#1\endcsname}
\newcommand{\fillin}[1]{{\color{arith}\itshape[#1]}}  % a place of the template that is not written yet; the build reports every one that is left
\newcommand{\leanquote}[1]{\csname LQ@#1\endcsname}
% generated by rb.py note; do not edit
\newcommand{\ConcDecls}{77}
\newcommand{\Pin}{72fead10}
\newcommand{\PinText}{that of commit {\leanfont 72fead10}}
\newcommand{\ArtifactName}{grh-compact}
\expandafter\newcommand\csname LN@IsInJumpWindow\endcsname{\href{https://github.com/lsernau/grh-lean/blob/72fead102f663d5c7a407964a13dc0cf9d977207/grh-compact/src/New/Lift/CokernelKillers.lean\#L17-L18}{{\leanfont Lift.\allowbreak{}IsInJumpWindow}}}
\expandafter\newcommand\csname LN@raisedAtE\endcsname{\href{https://github.com/lsernau/grh-lean/blob/72fead102f663d5c7a407964a13dc0cf9d977207/grh-compact/src/New/Lift/CokernelKillers.lean\#L23-L23}{{\leanfont Lift.\allowbreak{}raisedAtE}}}
\expandafter\newcommand\csname LN@twoWindows\endcsname{\href{https://github.com/lsernau/grh-lean/blob/72fead102f663d5c7a407964a13dc0cf9d977207/grh-compact/src/New/Lift/CokernelKillers.lean\#L27-L27}{{\leanfont Lift.\allowbreak{}twoWindows}}}
\expandafter\newcommand\csname LN@NoMultipleNear\endcsname{\href{https://github.com/lsernau/grh-lean/blob/72fead102f663d5c7a407964a13dc0cf9d977207/grh-compact/src/New/Lift/CokernelKillers.lean\#L30-L30}{{\leanfont Lift.\allowbreak{}NoMultipleNear}}}
\expandafter\newcommand\csname LN@sumtwoWindowsle\endcsname{\href{https://github.com/lsernau/grh-lean/blob/72fead102f663d5c7a407964a13dc0cf9d977207/grh-compact/src/New/Lift/CokernelKillers.lean\#L32-L38}{{\leanfont Lift.\allowbreak{}sum\_twoWindows\_le}}}
\expandafter\newcommand\csname LN@taupowsmulmemainfLatticeofchi\endcsname{\href{https://github.com/lsernau/grh-lean/blob/72fead102f663d5c7a407964a13dc0cf9d977207/grh-compact/src/New/Lift/CokernelKillers.lean\#L46-L51}{{\leanfont Lift.\allowbreak{}KummerChart.\allowbreak{}JumpLattice.\allowbreak{}tau\_pow\_smul\_mem\_ainfLattice\_of\_chi}}}
\expandafter\newcommand\csname LN@constantAt\endcsname{\href{https://github.com/lsernau/grh-lean/blob/72fead102f663d5c7a407964a13dc0cf9d977207/grh-compact/src/New/Lift/CokernelKillers.lean\#L57-L58}{{\leanfont Lift.\allowbreak{}KummerChart.\allowbreak{}PairingSide.\allowbreak{}constantAt}}}
\expandafter\newcommand\csname LN@liftIdeal\endcsname{\href{https://github.com/lsernau/grh-lean/blob/72fead102f663d5c7a407964a13dc0cf9d977207/grh-compact/src/New/Lift/CokernelKillers.lean\#L62-L64}{{\leanfont Lift.\allowbreak{}KummerChart.\allowbreak{}PairingSide.\allowbreak{}liftIdeal}}}
\expandafter\newcommand\csname LN@KillsCokernel\endcsname{\href{https://github.com/lsernau/grh-lean/blob/72fead102f663d5c7a407964a13dc0cf9d977207/grh-compact/src/New/Lift/CokernelKillers.lean\#L68-L68}{{\leanfont Lift.\allowbreak{}KummerChart.\allowbreak{}PairingSide.\allowbreak{}KillsCokernel}}}
\expandafter\newcommand\csname LN@memliftIdealiff\endcsname{\href{https://github.com/lsernau/grh-lean/blob/72fead102f663d5c7a407964a13dc0cf9d977207/grh-compact/src/New/Lift/CokernelKillers.lean\#L73-L76}{{\leanfont Lift.\allowbreak{}KummerChart.\allowbreak{}PairingSide.\allowbreak{}mem\_liftIdeal\_iff}}}
\expandafter\newcommand\csname LN@ppowmemliftIdealself\endcsname{\href{https://github.com/lsernau/grh-lean/blob/72fead102f663d5c7a407964a13dc0cf9d977207/grh-compact/src/New/Lift/CokernelKillers.lean\#L78-L79}{{\leanfont Lift.\allowbreak{}KummerChart.\allowbreak{}PairingSide.\allowbreak{}p\_pow\_mem\_liftIdeal\_self}}}
\expandafter\newcommand\csname LN@phipowsmulconstantAt\endcsname{\href{https://github.com/lsernau/grh-lean/blob/72fead102f663d5c7a407964a13dc0cf9d977207/grh-compact/src/New/Lift/CokernelKillers.lean\#L84-L85}{{\leanfont Lift.\allowbreak{}KummerChart.\allowbreak{}PairingSide.\allowbreak{}phi\_pow\_smul\_constantAt}}}
\expandafter\newcommand\csname LN@mulEpowmemliftIdealone\endcsname{\href{https://github.com/lsernau/grh-lean/blob/72fead102f663d5c7a407964a13dc0cf9d977207/grh-compact/src/New/Lift/CokernelKillers.lean\#L87-L90}{{\leanfont Lift.\allowbreak{}KummerChart.\allowbreak{}PairingSide.\allowbreak{}mul\_E\_pow\_mem\_liftIdeal\_one}}}
\expandafter\newcommand\csname LN@liftIdealleofppowmem\endcsname{\href{https://github.com/lsernau/grh-lean/blob/72fead102f663d5c7a407964a13dc0cf9d977207/grh-compact/src/New/Lift/CokernelKillers.lean\#L95-L99}{{\leanfont Lift.\allowbreak{}KummerChart.\allowbreak{}PairingSide.\allowbreak{}liftIdeal\_le\_of\_p\_pow\_mem}}}
\expandafter\newcommand\csname LN@constantAtmemlocLetterPart\endcsname{\href{https://github.com/lsernau/grh-lean/blob/72fead102f663d5c7a407964a13dc0cf9d977207/grh-compact/src/New/Lift/CokernelKillers.lean\#L101-L104}{{\leanfont Lift.\allowbreak{}KummerChart.\allowbreak{}PairingSide.\allowbreak{}constantAt\_mem\_locLetterPart}}}
\expandafter\newcommand\csname LN@phipowmemliftIdealiff\endcsname{\href{https://github.com/lsernau/grh-lean/blob/72fead102f663d5c7a407964a13dc0cf9d977207/grh-compact/src/New/Lift/CokernelKillers.lean\#L106-L111}{{\leanfont Lift.\allowbreak{}KummerChart.\allowbreak{}PairingSide.\allowbreak{}phi\_pow\_mem\_liftIdeal\_iff}}}
\expandafter\newcommand\csname LN@ppowmulzpowmemliftIdealiff\endcsname{\href{https://github.com/lsernau/grh-lean/blob/72fead102f663d5c7a407964a13dc0cf9d977207/grh-compact/src/New/Lift/CokernelKillers.lean\#L113-L115}{{\leanfont Lift.\allowbreak{}KummerChart.\allowbreak{}PairingSide.\allowbreak{}p\_pow\_mul\_z\_pow\_mem\_liftIdeal\_iff}}}
\expandafter\newcommand\csname LN@existsliftofmemliftIdeal\endcsname{\href{https://github.com/lsernau/grh-lean/blob/72fead102f663d5c7a407964a13dc0cf9d977207/grh-compact/src/New/Lift/CokernelKillers.lean\#L120-L127}{{\leanfont Lift.\allowbreak{}KummerChart.\allowbreak{}PairingSide.\allowbreak{}exists\_lift\_of\_mem\_liftIdeal}}}
\expandafter\newcommand\csname LN@memliftIdealofderivation\endcsname{\href{https://github.com/lsernau/grh-lean/blob/72fead102f663d5c7a407964a13dc0cf9d977207/grh-compact/src/New/Lift/CokernelKillers.lean\#L132-L144}{{\leanfont Lift.\allowbreak{}KummerChart.\allowbreak{}PairingSide.\allowbreak{}mem\_liftIdeal\_of\_derivation}}}
\expandafter\newcommand\csname LN@memliftIdealzeroofmemone\endcsname{\href{https://github.com/lsernau/grh-lean/blob/72fead102f663d5c7a407964a13dc0cf9d977207/grh-compact/src/New/Lift/CokernelKillers.lean\#L149-L157}{{\leanfont Lift.\allowbreak{}KummerChart.\allowbreak{}PairingSide.\allowbreak{}mem\_liftIdeal\_zero\_of\_mem\_one}}}
\expandafter\newcommand\csname LN@memliftIdealofmulEpow\endcsname{\href{https://github.com/lsernau/grh-lean/blob/72fead102f663d5c7a407964a13dc0cf9d977207/grh-compact/src/New/Lift/CokernelKillers.lean\#L162-L174}{{\leanfont Lift.\allowbreak{}KummerChart.\allowbreak{}PairingSide.\allowbreak{}mem\_liftIdeal\_of\_mul\_E\_pow}}}
\expandafter\newcommand\csname LN@existszpowmemliftIdeal\endcsname{\href{https://github.com/lsernau/grh-lean/blob/72fead102f663d5c7a407964a13dc0cf9d977207/grh-compact/src/New/Lift/CokernelKillers.lean\#L178-L188}{{\leanfont Lift.\allowbreak{}KummerChart.\allowbreak{}PairingSide.\allowbreak{}exists\_z\_pow\_mem\_liftIdeal}}}
\expandafter\newcommand\csname LN@killsCokernelofisUnramified\endcsname{\href{https://github.com/lsernau/grh-lean/blob/72fead102f663d5c7a407964a13dc0cf9d977207/grh-compact/src/New/Lift/CokernelKillers.lean\#L192-L202}{{\leanfont Lift.\allowbreak{}KummerChart.\allowbreak{}PairingSide.\allowbreak{}killsCokernel\_of\_isUnramified}}}
\expandafter\newcommand\csname LN@PairingSide\endcsname{\href{https://github.com/lsernau/grh-lean/blob/72fead102f663d5c7a407964a13dc0cf9d977207/grh-compact/src/New/Lift/CupProductBound.lean\#L17-L18}{{\leanfont Lift.\allowbreak{}KummerChart.\allowbreak{}PairingSide}}}
\expandafter\newcommand\csname LN@extLattice\endcsname{\href{https://github.com/lsernau/grh-lean/blob/72fead102f663d5c7a407964a13dc0cf9d977207/grh-compact/src/New/Lift/CupProductBound.lean\#L24-L27}{{\leanfont Lift.\allowbreak{}KummerChart.\allowbreak{}PairingSide.\allowbreak{}extLattice}}}
\expandafter\newcommand\csname LN@trivLattice\endcsname{\href{https://github.com/lsernau/grh-lean/blob/72fead102f663d5c7a407964a13dc0cf9d977207/grh-compact/src/New/Lift/CupProductBound.lean\#L31-L32}{{\leanfont Lift.\allowbreak{}KummerChart.\allowbreak{}PairingSide.\allowbreak{}trivLattice}}}
\expandafter\newcommand\csname LN@projTriv\endcsname{\href{https://github.com/lsernau/grh-lean/blob/72fead102f663d5c7a407964a13dc0cf9d977207/grh-compact/src/New/Lift/CupProductBound.lean\#L36-L36}{{\leanfont Lift.\allowbreak{}KummerChart.\allowbreak{}PairingSide.\allowbreak{}projTriv}}}
\expandafter\newcommand\csname LN@cost\endcsname{\href{https://github.com/lsernau/grh-lean/blob/72fead102f663d5c7a407964a13dc0cf9d977207/grh-compact/src/New/Lift/ExponentLattices.lean\#L24-L26}{{\leanfont Lift.\allowbreak{}KummerChart.\allowbreak{}cost}}}
\expandafter\newcommand\csname LN@existsdepthstep\endcsname{\href{https://github.com/lsernau/grh-lean/blob/72fead102f663d5c7a407964a13dc0cf9d977207/grh-compact/src/New/Lift/ExponentLattices.lean\#L32-L64}{{\leanfont Lift.\allowbreak{}KummerChart.\allowbreak{}exists\_depth\_step}}}
\expandafter\newcommand\csname LQ@existsdepthstep\endcsname{\leanquoteline{0}{lemma Lift.KummerChart.exists\_depth\_step (hn : n ≠ 0) : ∃ g : galMonoid, ∃ ω : \_ˣ, ∃ υ X,}\leanquoteline{2}{g • (k.frobE a).1 = ↑ω * k.frobE a ∧ (∀ s : k.S, (k.frobE a).1 ∣ g • s - s) ∧}\leanquoteline{2}{(k.frobE a).1 ∣ υ * X - p \textasciicircum{} cost a n ∧ zTau * X = (ω \textasciicircum{} n).1 - 1}}
\expandafter\newcommand\csname LN@existsdepthstepzero\endcsname{\href{https://github.com/lsernau/grh-lean/blob/72fead102f663d5c7a407964a13dc0cf9d977207/grh-compact/src/New/Lift/ExponentLattices.lean\#L68-L81}{{\leanfont Lift.\allowbreak{}KummerChart.\allowbreak{}exists\_depth\_step\_zero}}}
\expandafter\newcommand\csname LN@withExponents\endcsname{\href{https://github.com/lsernau/grh-lean/blob/72fead102f663d5c7a407964a13dc0cf9d977207/grh-compact/src/New/Lift/ExponentLattices.lean\#L88-L90}{{\leanfont Lift.\allowbreak{}KummerChart.\allowbreak{}JumpLattice.\allowbreak{}withExponents}}}
\expandafter\newcommand\csname LN@locLetterPart\endcsname{\href{https://github.com/lsernau/grh-lean/blob/72fead102f663d5c7a407964a13dc0cf9d977207/grh-compact/src/New/Lift/ExponentLattices.lean\#L94-L94}{{\leanfont Lift.\allowbreak{}KummerChart.\allowbreak{}JumpLattice.\allowbreak{}locLetterPart}}}
\expandafter\newcommand\csname LN@valueAtzero\endcsname{\href{https://github.com/lsernau/grh-lean/blob/72fead102f663d5c7a407964a13dc0cf9d977207/grh-compact/src/New/Lift/ExponentLattices.lean\#L99-L99}{{\leanfont Lift.\allowbreak{}KummerChart.\allowbreak{}JumpLattice.\allowbreak{}valueAt\_zero}}}
\expandafter\newcommand\csname LN@frobEsmulmemkisinLocAinfiff\endcsname{\href{https://github.com/lsernau/grh-lean/blob/72fead102f663d5c7a407964a13dc0cf9d977207/grh-compact/src/New/Lift/ExponentLattices.lean\#L101-L102}{{\leanfont Lift.\allowbreak{}KummerChart.\allowbreak{}JumpLattice.\allowbreak{}frobE\_smul\_mem\_kisinLocAinf\_iff}}}
\expandafter\newcommand\csname LN@frobEUnitzpowsmulmemkisinLoc\endcsname{\href{https://github.com/lsernau/grh-lean/blob/72fead102f663d5c7a407964a13dc0cf9d977207/grh-compact/src/New/Lift/ExponentLattices.lean\#L104-L108}{{\leanfont Lift.\allowbreak{}KummerChart.\allowbreak{}JumpLattice.\allowbreak{}frobEUnit\_zpow\_smul\_mem\_kisinLoc}}}
\expandafter\newcommand\csname LN@wittMapmemkisinLocAinf\endcsname{\href{https://github.com/lsernau/grh-lean/blob/72fead102f663d5c7a407964a13dc0cf9d977207/grh-compact/src/New/Lift/ExponentLattices.lean\#L110-L111}{{\leanfont Lift.\allowbreak{}KummerChart.\allowbreak{}JumpLattice.\allowbreak{}Hom.\allowbreak{}wittMap\_mem\_kisinLocAinf}}}
\expandafter\newcommand\csname LN@wittMapmemkisinLoc\endcsname{\href{https://github.com/lsernau/grh-lean/blob/72fead102f663d5c7a407964a13dc0cf9d977207/grh-compact/src/New/Lift/ExponentLattices.lean\#L113-L114}{{\leanfont Lift.\allowbreak{}KummerChart.\allowbreak{}JumpLattice.\allowbreak{}Hom.\allowbreak{}wittMap\_mem\_kisinLoc}}}
\expandafter\newcommand\csname LN@galsmulmemkisinLocAinf\endcsname{\href{https://github.com/lsernau/grh-lean/blob/72fead102f663d5c7a407964a13dc0cf9d977207/grh-compact/src/New/Lift/ExponentLattices.lean\#L116-L120}{{\leanfont Lift.\allowbreak{}KummerChart.\allowbreak{}JumpLattice.\allowbreak{}gal\_smul\_mem\_kisinLocAinf}}}
\expandafter\newcommand\csname LN@memwithExponentsofmemhodgeFil\endcsname{\href{https://github.com/lsernau/grh-lean/blob/72fead102f663d5c7a407964a13dc0cf9d977207/grh-compact/src/New/Lift/ExponentLattices.lean\#L122-L125}{{\leanfont Lift.\allowbreak{}KummerChart.\allowbreak{}JumpLattice.\allowbreak{}mem\_withExponents\_of\_mem\_hodgeFil}}}
\expandafter\newcommand\csname LN@memwithExponentsofle\endcsname{\href{https://github.com/lsernau/grh-lean/blob/72fead102f663d5c7a407964a13dc0cf9d977207/grh-compact/src/New/Lift/ExponentLattices.lean\#L127-L128}{{\leanfont Lift.\allowbreak{}KummerChart.\allowbreak{}JumpLattice.\allowbreak{}mem\_withExponents\_of\_le}}}
\expandafter\newcommand\csname LN@valueAtmemhodgeFilofmemainfLattice\endcsname{\href{https://github.com/lsernau/grh-lean/blob/72fead102f663d5c7a407964a13dc0cf9d977207/grh-compact/src/New/Lift/ExponentLattices.lean\#L130-L134}{{\leanfont Lift.\allowbreak{}KummerChart.\allowbreak{}JumpLattice.\allowbreak{}valueAt\_mem\_hodgeFil\_of\_mem\_ainfLattice}}}
\expandafter\newcommand\csname LN@mapmemwithExponents\endcsname{\href{https://github.com/lsernau/grh-lean/blob/72fead102f663d5c7a407964a13dc0cf9d977207/grh-compact/src/New/Lift/ExponentLattices.lean\#L139-L142}{{\leanfont Lift.\allowbreak{}KummerChart.\allowbreak{}JumpLattice.\allowbreak{}Hom.\allowbreak{}map\_mem\_withExponents}}}
\expandafter\newcommand\csname LN@memwithExponentsofsmul\endcsname{\href{https://github.com/lsernau/grh-lean/blob/72fead102f663d5c7a407964a13dc0cf9d977207/grh-compact/src/New/Lift/ExponentLattices.lean\#L147-L150}{{\leanfont Lift.\allowbreak{}KummerChart.\allowbreak{}JumpLattice.\allowbreak{}mem\_withExponents\_of\_smul}}}
\expandafter\newcommand\csname LN@frobEUnitzpowmemwithExponents\endcsname{\href{https://github.com/lsernau/grh-lean/blob/72fead102f663d5c7a407964a13dc0cf9d977207/grh-compact/src/New/Lift/ExponentLattices.lean\#L155-L161}{{\leanfont Lift.\allowbreak{}KummerChart.\allowbreak{}JumpLattice.\allowbreak{}frobEUnit\_zpow\_mem\_withExponents}}}
\expandafter\newcommand\csname LN@kisinLocAinflefracSpan\endcsname{\href{https://github.com/lsernau/grh-lean/blob/72fead102f663d5c7a407964a13dc0cf9d977207/grh-compact/src/New/Lift/ExponentLattices.lean\#L163-L164}{{\leanfont Lift.\allowbreak{}KummerChart.\allowbreak{}JumpLattice.\allowbreak{}kisinLocAinf\_le\_fracSpan}}}
\expandafter\newcommand\csname LN@valueAtsmulsub\endcsname{\href{https://github.com/lsernau/grh-lean/blob/72fead102f663d5c7a407964a13dc0cf9d977207/grh-compact/src/New/Lift/ExponentLattices.lean\#L166-L168}{{\leanfont Lift.\allowbreak{}KummerChart.\allowbreak{}JumpLattice.\allowbreak{}valueAt\_smul\_sub}}}
\expandafter\newcommand\csname LN@smulsubmemwithExponents\endcsname{\href{https://github.com/lsernau/grh-lean/blob/72fead102f663d5c7a407964a13dc0cf9d977207/grh-compact/src/New/Lift/ExponentLattices.lean\#L170-L172}{{\leanfont Lift.\allowbreak{}KummerChart.\allowbreak{}JumpLattice.\allowbreak{}smul\_sub\_mem\_withExponents}}}
\expandafter\newcommand\csname LN@smulsubmemzTaukisinLocAinf\endcsname{\href{https://github.com/lsernau/grh-lean/blob/72fead102f663d5c7a407964a13dc0cf9d977207/grh-compact/src/New/Lift/ExponentLattices.lean\#L176-L188}{{\leanfont Lift.\allowbreak{}KummerChart.\allowbreak{}JumpLattice.\allowbreak{}smul\_sub\_mem\_zTau\_kisinLocAinf}}}
\expandafter\newcommand\csname LN@phipowmemkisinLoc\endcsname{\href{https://github.com/lsernau/grh-lean/blob/72fead102f663d5c7a407964a13dc0cf9d977207/grh-compact/src/New/Lift/ExponentLattices.lean\#L190-L195}{{\leanfont Lift.\allowbreak{}KummerChart.\allowbreak{}JumpLattice.\allowbreak{}phi\_pow\_mem\_kisinLoc}}}
\expandafter\newcommand\csname LN@spankisinLoc\endcsname{\href{https://github.com/lsernau/grh-lean/blob/72fead102f663d5c7a407964a13dc0cf9d977207/grh-compact/src/New/Lift/ExponentLattices.lean\#L197-L201}{{\leanfont Lift.\allowbreak{}KummerChart.\allowbreak{}JumpLattice.\allowbreak{}span\_kisinLoc}}}
\expandafter\newcommand\csname LN@valueAtmemvalueField\endcsname{\href{https://github.com/lsernau/grh-lean/blob/72fead102f663d5c7a407964a13dc0cf9d977207/grh-compact/src/New/Lift/ExponentLattices.lean\#L203-L214}{{\leanfont Lift.\allowbreak{}KummerChart.\allowbreak{}JumpLattice.\allowbreak{}valueAt\_mem\_valueField}}}
\expandafter\newcommand\csname LN@existscoordRetract\endcsname{\href{https://github.com/lsernau/grh-lean/blob/72fead102f663d5c7a407964a13dc0cf9d977207/grh-compact/src/New/Lift/ExponentLattices.lean\#L216-L221}{{\leanfont Lift.\allowbreak{}KummerChart.\allowbreak{}JumpLattice.\allowbreak{}exists\_coordRetract}}}
\expandafter\newcommand\csname LN@coordRetract\endcsname{\href{https://github.com/lsernau/grh-lean/blob/72fead102f663d5c7a407964a13dc0cf9d977207/grh-compact/src/New/Lift/ExponentLattices.lean\#L225-L225}{{\leanfont Lift.\allowbreak{}KummerChart.\allowbreak{}JumpLattice.\allowbreak{}coordRetract}}}
\expandafter\newcommand\csname LN@coordRetractsum\endcsname{\href{https://github.com/lsernau/grh-lean/blob/72fead102f663d5c7a407964a13dc0cf9d977207/grh-compact/src/New/Lift/ExponentLattices.lean\#L227-L229}{{\leanfont Lift.\allowbreak{}KummerChart.\allowbreak{}JumpLattice.\allowbreak{}coordRetract\_sum}}}
\expandafter\newcommand\csname LN@coordRetractofmemkisinLoc\endcsname{\href{https://github.com/lsernau/grh-lean/blob/72fead102f663d5c7a407964a13dc0cf9d977207/grh-compact/src/New/Lift/ExponentLattices.lean\#L231-L232}{{\leanfont Lift.\allowbreak{}KummerChart.\allowbreak{}JumpLattice.\allowbreak{}coordRetract\_of\_mem\_kisinLoc}}}
\expandafter\newcommand\csname LN@mapcoordRetract\endcsname{\href{https://github.com/lsernau/grh-lean/blob/72fead102f663d5c7a407964a13dc0cf9d977207/grh-compact/src/New/Lift/ExponentLattices.lean\#L234-L238}{{\leanfont Lift.\allowbreak{}KummerChart.\allowbreak{}JumpLattice.\allowbreak{}Hom.\allowbreak{}map\_coordRetract}}}
\expandafter\newcommand\csname LN@valueAtmemofhodgeFille\endcsname{\href{https://github.com/lsernau/grh-lean/blob/72fead102f663d5c7a407964a13dc0cf9d977207/grh-compact/src/New/Lift/ExponentLattices.lean\#L243-L252}{{\leanfont Lift.\allowbreak{}KummerChart.\allowbreak{}JumpLattice.\allowbreak{}valueAt\_mem\_of\_hodgeFil\_le}}}
\expandafter\newcommand\csname LN@smulsubmemzTauUnitsmulofdeep\endcsname{\href{https://github.com/lsernau/grh-lean/blob/72fead102f663d5c7a407964a13dc0cf9d977207/grh-compact/src/New/Lift/ExponentLattices.lean\#L256-L258}{{\leanfont Lift.\allowbreak{}KummerChart.\allowbreak{}JumpLattice.\allowbreak{}smul\_sub\_mem\_zTauUnit\_smul\_of\_deep}}}
\expandafter\newcommand\csname LN@locKisinPart\endcsname{\href{https://github.com/lsernau/grh-lean/blob/72fead102f663d5c7a407964a13dc0cf9d977207/grh-compact/src/New/Lift/HodgeExpansion.lean\#L23-L23}{{\leanfont Lift.\allowbreak{}KummerChart.\allowbreak{}JumpLattice.\allowbreak{}locKisinPart}}}
\expandafter\newcommand\csname LN@invariantsMod\endcsname{\href{https://github.com/lsernau/grh-lean/blob/72fead102f663d5c7a407964a13dc0cf9d977207/grh-compact/src/New/Lift/HodgeExpansion.lean\#L27-L28}{{\leanfont Lift.\allowbreak{}KummerChart.\allowbreak{}JumpLattice.\allowbreak{}invariantsMod}}}
\expandafter\newcommand\csname LN@costleofisUnramified\endcsname{\href{https://github.com/lsernau/grh-lean/blob/72fead102f663d5c7a407964a13dc0cf9d977207/grh-compact/src/New/Lift/HodgeExpansion.lean\#L30-L31}{{\leanfont Lift.\allowbreak{}KummerChart.\allowbreak{}JumpLattice.\allowbreak{}cost\_le\_of\_isUnramified}}}
\expandafter\newcommand\csname LN@improvedExponent\endcsname{\href{https://github.com/lsernau/grh-lean/blob/72fead102f663d5c7a407964a13dc0cf9d977207/grh-compact/src/New/Lift/HodgeExpansion.lean\#L35-L36}{{\leanfont Lift.\allowbreak{}KummerChart.\allowbreak{}JumpLattice.\allowbreak{}improvedExponent}}}
\expandafter\newcommand\csname LN@memwithExponentsofdeep\endcsname{\href{https://github.com/lsernau/grh-lean/blob/72fead102f663d5c7a407964a13dc0cf9d977207/grh-compact/src/New/Lift/HodgeExpansion.lean\#L41-L46}{{\leanfont Lift.\allowbreak{}KummerChart.\allowbreak{}JumpLattice.\allowbreak{}mem\_withExponents\_of\_deep}}}
\expandafter\newcommand\csname LN@galMonoidsmullocalFieldOf\endcsname{\href{https://github.com/lsernau/grh-lean/blob/72fead102f663d5c7a407964a13dc0cf9d977207/grh-compact/src/New/Lift/HodgeExpansion.lean\#L48-L49}{{\leanfont Lift.\allowbreak{}KummerChart.\allowbreak{}JumpLattice.\allowbreak{}galMonoid\_smul\_localFieldOf}}}
\expandafter\newcommand\csname LN@existsreachesdepthsucc\endcsname{\href{https://github.com/lsernau/grh-lean/blob/72fead102f663d5c7a407964a13dc0cf9d977207/grh-compact/src/New/Lift/HodgeExpansion.lean\#L54-L74}{{\leanfont Lift.\allowbreak{}KummerChart.\allowbreak{}JumpLattice.\allowbreak{}exists\_reaches\_depth\_succ}}}
\expandafter\newcommand\csname LN@wittPairinga\endcsname{\href{https://github.com/lsernau/grh-lean/blob/72fead102f663d5c7a407964a13dc0cf9d977207/grh-compact/src/New/Lift/HodgeExpansion.lean\#L82-L83}{{\leanfont Lift.\allowbreak{}KummerChart.\allowbreak{}JumpLattice.\allowbreak{}wittPairing}}}
\expandafter\newcommand\csname LN@wittPairingletterProjleft\endcsname{\href{https://github.com/lsernau/grh-lean/blob/72fead102f663d5c7a407964a13dc0cf9d977207/grh-compact/src/New/Lift/HodgeExpansion.lean\#L87-L88}{{\leanfont Lift.\allowbreak{}KummerChart.\allowbreak{}JumpLattice.\allowbreak{}wittPairing\_letterProj\_left}}}
\expandafter\newcommand\csname LN@wittPairingmemsaturateBy\endcsname{\href{https://github.com/lsernau/grh-lean/blob/72fead102f663d5c7a407964a13dc0cf9d977207/grh-compact/src/New/Lift/HodgeExpansion.lean\#L90-L92}{{\leanfont Lift.\allowbreak{}KummerChart.\allowbreak{}JumpLattice.\allowbreak{}wittPairing\_mem\_saturateBy}}}
\expandafter\newcommand\csname LN@frobExponent\endcsname{\href{https://github.com/lsernau/grh-lean/blob/72fead102f663d5c7a407964a13dc0cf9d977207/grh-compact/src/New/Lift/HodgeExpansion.lean\#L96-L98}{{\leanfont Lift.\allowbreak{}KummerChart.\allowbreak{}JumpLattice.\allowbreak{}frobExponent}}}
\expandafter\newcommand\csname LN@memhodgeFilofhodgebasis\endcsname{\href{https://github.com/lsernau/grh-lean/blob/72fead102f663d5c7a407964a13dc0cf9d977207/grh-compact/src/New/Lift/HodgeExpansion.lean\#L102-L118}{{\leanfont Lift.\allowbreak{}KummerChart.\allowbreak{}JumpLattice.\allowbreak{}mem\_hodgeFil\_of\_hodge\_basis}}}
\expandafter\newcommand\csname LN@valueAtcoordRetractsubmemhodgeFil\endcsname{\href{https://github.com/lsernau/grh-lean/blob/72fead102f663d5c7a407964a13dc0cf9d977207/grh-compact/src/New/Lift/HodgeExpansion.lean\#L120-L125}{{\leanfont Lift.\allowbreak{}KummerChart.\allowbreak{}JumpLattice.\allowbreak{}valueAt\_coordRetract\_sub\_mem\_hodgeFil}}}
\expandafter\newcommand\csname LN@coordRetractmemwithExponents\endcsname{\href{https://github.com/lsernau/grh-lean/blob/72fead102f663d5c7a407964a13dc0cf9d977207/grh-compact/src/New/Lift/HodgeExpansion.lean\#L127-L132}{{\leanfont Lift.\allowbreak{}KummerChart.\allowbreak{}JumpLattice.\allowbreak{}coordRetract\_mem\_withExponents}}}
\expandafter\newcommand\csname LN@phipowmemwithExponents\endcsname{\href{https://github.com/lsernau/grh-lean/blob/72fead102f663d5c7a407964a13dc0cf9d977207/grh-compact/src/New/Lift/HodgeExpansion.lean\#L134-L143}{{\leanfont Lift.\allowbreak{}KummerChart.\allowbreak{}JumpLattice.\allowbreak{}phi\_pow\_mem\_withExponents}}}
\expandafter\newcommand\csname LN@memwithExponentsstdExponentiff\endcsname{\href{https://github.com/lsernau/grh-lean/blob/72fead102f663d5c7a407964a13dc0cf9d977207/grh-compact/src/New/Lift/HodgeExpansion.lean\#L148-L174}{{\leanfont Lift.\allowbreak{}KummerChart.\allowbreak{}JumpLattice.\allowbreak{}mem\_withExponents\_stdExponent\_iff}}}
\expandafter\newcommand\csname LQ@memwithExponentsstdExponentiff\endcsname{\leanquoteline{0}{lemma mem\_withExponents\_stdExponent\_iff : x ∈ withExponents i (stdExponent l) ↔}\leanquoteline{2}{x ∈ frobKisinAinf l ∧ x ∈ letterPart i}}
\expandafter\newcommand\csname LN@letterProjmemwithExponents\endcsname{\href{https://github.com/lsernau/grh-lean/blob/72fead102f663d5c7a407964a13dc0cf9d977207/grh-compact/src/New/Lift/HodgeExpansion.lean\#L176-L179}{{\leanfont Lift.\allowbreak{}KummerChart.\allowbreak{}JumpLattice.\allowbreak{}letterProj\_mem\_withExponents}}}
\expandafter\newcommand\csname LN@letterProjmemofvalueFieldPoints\endcsname{\href{https://github.com/lsernau/grh-lean/blob/72fead102f663d5c7a407964a13dc0cf9d977207/grh-compact/src/New/Lift/HodgeExpansion.lean\#L181-L183}{{\leanfont Lift.\allowbreak{}KummerChart.\allowbreak{}JumpLattice.\allowbreak{}letterProj\_mem\_of\_valueFieldPoints}}}
\expandafter\newcommand\csname LN@wittPairingmemwithExponents\endcsname{\href{https://github.com/lsernau/grh-lean/blob/72fead102f663d5c7a407964a13dc0cf9d977207/grh-compact/src/New/Lift/HodgeExpansion.lean\#L188-L207}{{\leanfont Lift.\allowbreak{}KummerChart.\allowbreak{}JumpLattice.\allowbreak{}wittPairing\_mem\_withExponents}}}
\newcommand{\ConcordanceRows}{\crow{IsInJumpWindow}{\href{https://github.com/lsernau/grh-lean/blob/72fead102f663d5c7a407964a13dc0cf9d977207/grh-compact/src/New/Lift/CokernelKillers.lean\#L17-L18}{{\leanfont Lift.\allowbreak{}IsInJumpWindow}}}{CokernelKillers.lean}{17--18}
\crow{raisedAtE}{\href{https://github.com/lsernau/grh-lean/blob/72fead102f663d5c7a407964a13dc0cf9d977207/grh-compact/src/New/Lift/CokernelKillers.lean\#L23-L23}{{\leanfont Lift.\allowbreak{}raisedAtE}}}{CokernelKillers.lean}{23--23}
\crow{twoWindows}{\href{https://github.com/lsernau/grh-lean/blob/72fead102f663d5c7a407964a13dc0cf9d977207/grh-compact/src/New/Lift/CokernelKillers.lean\#L27-L27}{{\leanfont Lift.\allowbreak{}twoWindows}}}{CokernelKillers.lean}{27--27}
\crow{NoMultipleNear}{\href{https://github.com/lsernau/grh-lean/blob/72fead102f663d5c7a407964a13dc0cf9d977207/grh-compact/src/New/Lift/CokernelKillers.lean\#L30-L30}{{\leanfont Lift.\allowbreak{}NoMultipleNear}}}{CokernelKillers.lean}{30--30}
\crow{sumtwoWindowsle}{\href{https://github.com/lsernau/grh-lean/blob/72fead102f663d5c7a407964a13dc0cf9d977207/grh-compact/src/New/Lift/CokernelKillers.lean\#L32-L38}{{\leanfont Lift.\allowbreak{}sum\_twoWindows\_le}}}{CokernelKillers.lean}{32--38}
\crow{taupowsmulmemainfLatticeofchi}{\href{https://github.com/lsernau/grh-lean/blob/72fead102f663d5c7a407964a13dc0cf9d977207/grh-compact/src/New/Lift/CokernelKillers.lean\#L46-L51}{{\leanfont Lift.\allowbreak{}KummerChart.\allowbreak{}JumpLattice.\allowbreak{}tau\_pow\_smul\_mem\_ainfLattice\_of\_chi}}}{CokernelKillers.lean}{46--51}
\crow{constantAt}{\href{https://github.com/lsernau/grh-lean/blob/72fead102f663d5c7a407964a13dc0cf9d977207/grh-compact/src/New/Lift/CokernelKillers.lean\#L57-L58}{{\leanfont Lift.\allowbreak{}KummerChart.\allowbreak{}PairingSide.\allowbreak{}constantAt}}}{CokernelKillers.lean}{57--58}
\crow{liftIdeal}{\href{https://github.com/lsernau/grh-lean/blob/72fead102f663d5c7a407964a13dc0cf9d977207/grh-compact/src/New/Lift/CokernelKillers.lean\#L62-L64}{{\leanfont Lift.\allowbreak{}KummerChart.\allowbreak{}PairingSide.\allowbreak{}liftIdeal}}}{CokernelKillers.lean}{62--64}
\crow{KillsCokernel}{\href{https://github.com/lsernau/grh-lean/blob/72fead102f663d5c7a407964a13dc0cf9d977207/grh-compact/src/New/Lift/CokernelKillers.lean\#L68-L68}{{\leanfont Lift.\allowbreak{}KummerChart.\allowbreak{}PairingSide.\allowbreak{}KillsCokernel}}}{CokernelKillers.lean}{68--68}
\crow{memliftIdealiff}{\href{https://github.com/lsernau/grh-lean/blob/72fead102f663d5c7a407964a13dc0cf9d977207/grh-compact/src/New/Lift/CokernelKillers.lean\#L73-L76}{{\leanfont Lift.\allowbreak{}KummerChart.\allowbreak{}PairingSide.\allowbreak{}mem\_liftIdeal\_iff}}}{CokernelKillers.lean}{73--76}
\crow{ppowmemliftIdealself}{\href{https://github.com/lsernau/grh-lean/blob/72fead102f663d5c7a407964a13dc0cf9d977207/grh-compact/src/New/Lift/CokernelKillers.lean\#L78-L79}{{\leanfont Lift.\allowbreak{}KummerChart.\allowbreak{}PairingSide.\allowbreak{}p\_pow\_mem\_liftIdeal\_self}}}{CokernelKillers.lean}{78--79}
\crow{phipowsmulconstantAt}{\href{https://github.com/lsernau/grh-lean/blob/72fead102f663d5c7a407964a13dc0cf9d977207/grh-compact/src/New/Lift/CokernelKillers.lean\#L84-L85}{{\leanfont Lift.\allowbreak{}KummerChart.\allowbreak{}PairingSide.\allowbreak{}phi\_pow\_smul\_constantAt}}}{CokernelKillers.lean}{84--85}
\crow{mulEpowmemliftIdealone}{\href{https://github.com/lsernau/grh-lean/blob/72fead102f663d5c7a407964a13dc0cf9d977207/grh-compact/src/New/Lift/CokernelKillers.lean\#L87-L90}{{\leanfont Lift.\allowbreak{}KummerChart.\allowbreak{}PairingSide.\allowbreak{}mul\_E\_pow\_mem\_liftIdeal\_one}}}{CokernelKillers.lean}{87--90}
\crow{liftIdealleofppowmem}{\href{https://github.com/lsernau/grh-lean/blob/72fead102f663d5c7a407964a13dc0cf9d977207/grh-compact/src/New/Lift/CokernelKillers.lean\#L95-L99}{{\leanfont Lift.\allowbreak{}KummerChart.\allowbreak{}PairingSide.\allowbreak{}liftIdeal\_le\_of\_p\_pow\_mem}}}{CokernelKillers.lean}{95--99}
\crow{constantAtmemlocLetterPart}{\href{https://github.com/lsernau/grh-lean/blob/72fead102f663d5c7a407964a13dc0cf9d977207/grh-compact/src/New/Lift/CokernelKillers.lean\#L101-L104}{{\leanfont Lift.\allowbreak{}KummerChart.\allowbreak{}PairingSide.\allowbreak{}constantAt\_mem\_locLetterPart}}}{CokernelKillers.lean}{101--104}
\crow{phipowmemliftIdealiff}{\href{https://github.com/lsernau/grh-lean/blob/72fead102f663d5c7a407964a13dc0cf9d977207/grh-compact/src/New/Lift/CokernelKillers.lean\#L106-L111}{{\leanfont Lift.\allowbreak{}KummerChart.\allowbreak{}PairingSide.\allowbreak{}phi\_pow\_mem\_liftIdeal\_iff}}}{CokernelKillers.lean}{106--111}
\crow{ppowmulzpowmemliftIdealiff}{\href{https://github.com/lsernau/grh-lean/blob/72fead102f663d5c7a407964a13dc0cf9d977207/grh-compact/src/New/Lift/CokernelKillers.lean\#L113-L115}{{\leanfont Lift.\allowbreak{}KummerChart.\allowbreak{}PairingSide.\allowbreak{}p\_pow\_mul\_z\_pow\_mem\_liftIdeal\_iff}}}{CokernelKillers.lean}{113--115}
\crow{existsliftofmemliftIdeal}{\href{https://github.com/lsernau/grh-lean/blob/72fead102f663d5c7a407964a13dc0cf9d977207/grh-compact/src/New/Lift/CokernelKillers.lean\#L120-L127}{{\leanfont Lift.\allowbreak{}KummerChart.\allowbreak{}PairingSide.\allowbreak{}exists\_lift\_of\_mem\_liftIdeal}}}{CokernelKillers.lean}{120--127}
\crow{memliftIdealofderivation}{\href{https://github.com/lsernau/grh-lean/blob/72fead102f663d5c7a407964a13dc0cf9d977207/grh-compact/src/New/Lift/CokernelKillers.lean\#L132-L144}{{\leanfont Lift.\allowbreak{}KummerChart.\allowbreak{}PairingSide.\allowbreak{}mem\_liftIdeal\_of\_derivation}}}{CokernelKillers.lean}{132--144}
\crow{memliftIdealzeroofmemone}{\href{https://github.com/lsernau/grh-lean/blob/72fead102f663d5c7a407964a13dc0cf9d977207/grh-compact/src/New/Lift/CokernelKillers.lean\#L149-L157}{{\leanfont Lift.\allowbreak{}KummerChart.\allowbreak{}PairingSide.\allowbreak{}mem\_liftIdeal\_zero\_of\_mem\_one}}}{CokernelKillers.lean}{149--157}
\crow{memliftIdealofmulEpow}{\href{https://github.com/lsernau/grh-lean/blob/72fead102f663d5c7a407964a13dc0cf9d977207/grh-compact/src/New/Lift/CokernelKillers.lean\#L162-L174}{{\leanfont Lift.\allowbreak{}KummerChart.\allowbreak{}PairingSide.\allowbreak{}mem\_liftIdeal\_of\_mul\_E\_pow}}}{CokernelKillers.lean}{162--174}
\crow{existszpowmemliftIdeal}{\href{https://github.com/lsernau/grh-lean/blob/72fead102f663d5c7a407964a13dc0cf9d977207/grh-compact/src/New/Lift/CokernelKillers.lean\#L178-L188}{{\leanfont Lift.\allowbreak{}KummerChart.\allowbreak{}PairingSide.\allowbreak{}exists\_z\_pow\_mem\_liftIdeal}}}{CokernelKillers.lean}{178--188}
\crow{killsCokernelofisUnramified}{\href{https://github.com/lsernau/grh-lean/blob/72fead102f663d5c7a407964a13dc0cf9d977207/grh-compact/src/New/Lift/CokernelKillers.lean\#L192-L202}{{\leanfont Lift.\allowbreak{}KummerChart.\allowbreak{}PairingSide.\allowbreak{}killsCokernel\_of\_isUnramified}}}{CokernelKillers.lean}{192--202}
\crow{PairingSide}{\href{https://github.com/lsernau/grh-lean/blob/72fead102f663d5c7a407964a13dc0cf9d977207/grh-compact/src/New/Lift/CupProductBound.lean\#L17-L18}{{\leanfont Lift.\allowbreak{}KummerChart.\allowbreak{}PairingSide}}}{CupProductBound.lean}{17--18}
\crow{extLattice}{\href{https://github.com/lsernau/grh-lean/blob/72fead102f663d5c7a407964a13dc0cf9d977207/grh-compact/src/New/Lift/CupProductBound.lean\#L24-L27}{{\leanfont Lift.\allowbreak{}KummerChart.\allowbreak{}PairingSide.\allowbreak{}extLattice}}}{CupProductBound.lean}{24--27}
\crow{trivLattice}{\href{https://github.com/lsernau/grh-lean/blob/72fead102f663d5c7a407964a13dc0cf9d977207/grh-compact/src/New/Lift/CupProductBound.lean\#L31-L32}{{\leanfont Lift.\allowbreak{}KummerChart.\allowbreak{}PairingSide.\allowbreak{}trivLattice}}}{CupProductBound.lean}{31--32}
\crow{projTriv}{\href{https://github.com/lsernau/grh-lean/blob/72fead102f663d5c7a407964a13dc0cf9d977207/grh-compact/src/New/Lift/CupProductBound.lean\#L36-L36}{{\leanfont Lift.\allowbreak{}KummerChart.\allowbreak{}PairingSide.\allowbreak{}projTriv}}}{CupProductBound.lean}{36--36}
\crow{cost}{\href{https://github.com/lsernau/grh-lean/blob/72fead102f663d5c7a407964a13dc0cf9d977207/grh-compact/src/New/Lift/ExponentLattices.lean\#L24-L26}{{\leanfont Lift.\allowbreak{}KummerChart.\allowbreak{}cost}}}{ExponentLattices.lean}{24--26}
\crow{existsdepthstep}{\href{https://github.com/lsernau/grh-lean/blob/72fead102f663d5c7a407964a13dc0cf9d977207/grh-compact/src/New/Lift/ExponentLattices.lean\#L32-L64}{{\leanfont Lift.\allowbreak{}KummerChart.\allowbreak{}exists\_depth\_step}}}{ExponentLattices.lean}{32--64}
\crow{existsdepthstepzero}{\href{https://github.com/lsernau/grh-lean/blob/72fead102f663d5c7a407964a13dc0cf9d977207/grh-compact/src/New/Lift/ExponentLattices.lean\#L68-L81}{{\leanfont Lift.\allowbreak{}KummerChart.\allowbreak{}exists\_depth\_step\_zero}}}{ExponentLattices.lean}{68--81}
\crow{withExponents}{\href{https://github.com/lsernau/grh-lean/blob/72fead102f663d5c7a407964a13dc0cf9d977207/grh-compact/src/New/Lift/ExponentLattices.lean\#L88-L90}{{\leanfont Lift.\allowbreak{}KummerChart.\allowbreak{}JumpLattice.\allowbreak{}withExponents}}}{ExponentLattices.lean}{88--90}
\crow{locLetterPart}{\href{https://github.com/lsernau/grh-lean/blob/72fead102f663d5c7a407964a13dc0cf9d977207/grh-compact/src/New/Lift/ExponentLattices.lean\#L94-L94}{{\leanfont Lift.\allowbreak{}KummerChart.\allowbreak{}JumpLattice.\allowbreak{}locLetterPart}}}{ExponentLattices.lean}{94--94}
\crow{valueAtzero}{\href{https://github.com/lsernau/grh-lean/blob/72fead102f663d5c7a407964a13dc0cf9d977207/grh-compact/src/New/Lift/ExponentLattices.lean\#L99-L99}{{\leanfont Lift.\allowbreak{}KummerChart.\allowbreak{}JumpLattice.\allowbreak{}valueAt\_zero}}}{ExponentLattices.lean}{99--99}
\crow{frobEsmulmemkisinLocAinfiff}{\href{https://github.com/lsernau/grh-lean/blob/72fead102f663d5c7a407964a13dc0cf9d977207/grh-compact/src/New/Lift/ExponentLattices.lean\#L101-L102}{{\leanfont Lift.\allowbreak{}KummerChart.\allowbreak{}JumpLattice.\allowbreak{}frobE\_smul\_mem\_kisinLocAinf\_iff}}}{ExponentLattices.lean}{101--102}
\crow{frobEUnitzpowsmulmemkisinLoc}{\href{https://github.com/lsernau/grh-lean/blob/72fead102f663d5c7a407964a13dc0cf9d977207/grh-compact/src/New/Lift/ExponentLattices.lean\#L104-L108}{{\leanfont Lift.\allowbreak{}KummerChart.\allowbreak{}JumpLattice.\allowbreak{}frobEUnit\_zpow\_smul\_mem\_kisinLoc}}}{ExponentLattices.lean}{104--108}
\crow{wittMapmemkisinLocAinf}{\href{https://github.com/lsernau/grh-lean/blob/72fead102f663d5c7a407964a13dc0cf9d977207/grh-compact/src/New/Lift/ExponentLattices.lean\#L110-L111}{{\leanfont Lift.\allowbreak{}KummerChart.\allowbreak{}JumpLattice.\allowbreak{}Hom.\allowbreak{}wittMap\_mem\_kisinLocAinf}}}{ExponentLattices.lean}{110--111}
\crow{wittMapmemkisinLoc}{\href{https://github.com/lsernau/grh-lean/blob/72fead102f663d5c7a407964a13dc0cf9d977207/grh-compact/src/New/Lift/ExponentLattices.lean\#L113-L114}{{\leanfont Lift.\allowbreak{}KummerChart.\allowbreak{}JumpLattice.\allowbreak{}Hom.\allowbreak{}wittMap\_mem\_kisinLoc}}}{ExponentLattices.lean}{113--114}
\crow{galsmulmemkisinLocAinf}{\href{https://github.com/lsernau/grh-lean/blob/72fead102f663d5c7a407964a13dc0cf9d977207/grh-compact/src/New/Lift/ExponentLattices.lean\#L116-L120}{{\leanfont Lift.\allowbreak{}KummerChart.\allowbreak{}JumpLattice.\allowbreak{}gal\_smul\_mem\_kisinLocAinf}}}{ExponentLattices.lean}{116--120}
\crow{memwithExponentsofmemhodgeFil}{\href{https://github.com/lsernau/grh-lean/blob/72fead102f663d5c7a407964a13dc0cf9d977207/grh-compact/src/New/Lift/ExponentLattices.lean\#L122-L125}{{\leanfont Lift.\allowbreak{}KummerChart.\allowbreak{}JumpLattice.\allowbreak{}mem\_withExponents\_of\_mem\_hodgeFil}}}{ExponentLattices.lean}{122--125}
\crow{memwithExponentsofle}{\href{https://github.com/lsernau/grh-lean/blob/72fead102f663d5c7a407964a13dc0cf9d977207/grh-compact/src/New/Lift/ExponentLattices.lean\#L127-L128}{{\leanfont Lift.\allowbreak{}KummerChart.\allowbreak{}JumpLattice.\allowbreak{}mem\_withExponents\_of\_le}}}{ExponentLattices.lean}{127--128}
\crow{valueAtmemhodgeFilofmemainfLattice}{\href{https://github.com/lsernau/grh-lean/blob/72fead102f663d5c7a407964a13dc0cf9d977207/grh-compact/src/New/Lift/ExponentLattices.lean\#L130-L134}{{\leanfont Lift.\allowbreak{}KummerChart.\allowbreak{}JumpLattice.\allowbreak{}valueAt\_mem\_hodgeFil\_of\_mem\_ainfLattice}}}{ExponentLattices.lean}{130--134}
\crow{mapmemwithExponents}{\href{https://github.com/lsernau/grh-lean/blob/72fead102f663d5c7a407964a13dc0cf9d977207/grh-compact/src/New/Lift/ExponentLattices.lean\#L139-L142}{{\leanfont Lift.\allowbreak{}KummerChart.\allowbreak{}JumpLattice.\allowbreak{}Hom.\allowbreak{}map\_mem\_withExponents}}}{ExponentLattices.lean}{139--142}
\crow{memwithExponentsofsmul}{\href{https://github.com/lsernau/grh-lean/blob/72fead102f663d5c7a407964a13dc0cf9d977207/grh-compact/src/New/Lift/ExponentLattices.lean\#L147-L150}{{\leanfont Lift.\allowbreak{}KummerChart.\allowbreak{}JumpLattice.\allowbreak{}mem\_withExponents\_of\_smul}}}{ExponentLattices.lean}{147--150}
\crow{frobEUnitzpowmemwithExponents}{\href{https://github.com/lsernau/grh-lean/blob/72fead102f663d5c7a407964a13dc0cf9d977207/grh-compact/src/New/Lift/ExponentLattices.lean\#L155-L161}{{\leanfont Lift.\allowbreak{}KummerChart.\allowbreak{}JumpLattice.\allowbreak{}frobEUnit\_zpow\_mem\_withExponents}}}{ExponentLattices.lean}{155--161}
\crow{kisinLocAinflefracSpan}{\href{https://github.com/lsernau/grh-lean/blob/72fead102f663d5c7a407964a13dc0cf9d977207/grh-compact/src/New/Lift/ExponentLattices.lean\#L163-L164}{{\leanfont Lift.\allowbreak{}KummerChart.\allowbreak{}JumpLattice.\allowbreak{}kisinLocAinf\_le\_fracSpan}}}{ExponentLattices.lean}{163--164}
\crow{valueAtsmulsub}{\href{https://github.com/lsernau/grh-lean/blob/72fead102f663d5c7a407964a13dc0cf9d977207/grh-compact/src/New/Lift/ExponentLattices.lean\#L166-L168}{{\leanfont Lift.\allowbreak{}KummerChart.\allowbreak{}JumpLattice.\allowbreak{}valueAt\_smul\_sub}}}{ExponentLattices.lean}{166--168}
\crow{smulsubmemwithExponents}{\href{https://github.com/lsernau/grh-lean/blob/72fead102f663d5c7a407964a13dc0cf9d977207/grh-compact/src/New/Lift/ExponentLattices.lean\#L170-L172}{{\leanfont Lift.\allowbreak{}KummerChart.\allowbreak{}JumpLattice.\allowbreak{}smul\_sub\_mem\_withExponents}}}{ExponentLattices.lean}{170--172}
\crow{smulsubmemzTaukisinLocAinf}{\href{https://github.com/lsernau/grh-lean/blob/72fead102f663d5c7a407964a13dc0cf9d977207/grh-compact/src/New/Lift/ExponentLattices.lean\#L176-L188}{{\leanfont Lift.\allowbreak{}KummerChart.\allowbreak{}JumpLattice.\allowbreak{}smul\_sub\_mem\_zTau\_kisinLocAinf}}}{ExponentLattices.lean}{176--188}
\crow{phipowmemkisinLoc}{\href{https://github.com/lsernau/grh-lean/blob/72fead102f663d5c7a407964a13dc0cf9d977207/grh-compact/src/New/Lift/ExponentLattices.lean\#L190-L195}{{\leanfont Lift.\allowbreak{}KummerChart.\allowbreak{}JumpLattice.\allowbreak{}phi\_pow\_mem\_kisinLoc}}}{ExponentLattices.lean}{190--195}
\crow{spankisinLoc}{\href{https://github.com/lsernau/grh-lean/blob/72fead102f663d5c7a407964a13dc0cf9d977207/grh-compact/src/New/Lift/ExponentLattices.lean\#L197-L201}{{\leanfont Lift.\allowbreak{}KummerChart.\allowbreak{}JumpLattice.\allowbreak{}span\_kisinLoc}}}{ExponentLattices.lean}{197--201}
\crow{valueAtmemvalueField}{\href{https://github.com/lsernau/grh-lean/blob/72fead102f663d5c7a407964a13dc0cf9d977207/grh-compact/src/New/Lift/ExponentLattices.lean\#L203-L214}{{\leanfont Lift.\allowbreak{}KummerChart.\allowbreak{}JumpLattice.\allowbreak{}valueAt\_mem\_valueField}}}{ExponentLattices.lean}{203--214}
\crow{existscoordRetract}{\href{https://github.com/lsernau/grh-lean/blob/72fead102f663d5c7a407964a13dc0cf9d977207/grh-compact/src/New/Lift/ExponentLattices.lean\#L216-L221}{{\leanfont Lift.\allowbreak{}KummerChart.\allowbreak{}JumpLattice.\allowbreak{}exists\_coordRetract}}}{ExponentLattices.lean}{216--221}
\crow{coordRetract}{\href{https://github.com/lsernau/grh-lean/blob/72fead102f663d5c7a407964a13dc0cf9d977207/grh-compact/src/New/Lift/ExponentLattices.lean\#L225-L225}{{\leanfont Lift.\allowbreak{}KummerChart.\allowbreak{}JumpLattice.\allowbreak{}coordRetract}}}{ExponentLattices.lean}{225--225}
\crow{coordRetractsum}{\href{https://github.com/lsernau/grh-lean/blob/72fead102f663d5c7a407964a13dc0cf9d977207/grh-compact/src/New/Lift/ExponentLattices.lean\#L227-L229}{{\leanfont Lift.\allowbreak{}KummerChart.\allowbreak{}JumpLattice.\allowbreak{}coordRetract\_sum}}}{ExponentLattices.lean}{227--229}
\crow{coordRetractofmemkisinLoc}{\href{https://github.com/lsernau/grh-lean/blob/72fead102f663d5c7a407964a13dc0cf9d977207/grh-compact/src/New/Lift/ExponentLattices.lean\#L231-L232}{{\leanfont Lift.\allowbreak{}KummerChart.\allowbreak{}JumpLattice.\allowbreak{}coordRetract\_of\_mem\_kisinLoc}}}{ExponentLattices.lean}{231--232}
\crow{mapcoordRetract}{\href{https://github.com/lsernau/grh-lean/blob/72fead102f663d5c7a407964a13dc0cf9d977207/grh-compact/src/New/Lift/ExponentLattices.lean\#L234-L238}{{\leanfont Lift.\allowbreak{}KummerChart.\allowbreak{}JumpLattice.\allowbreak{}Hom.\allowbreak{}map\_coordRetract}}}{ExponentLattices.lean}{234--238}
\crow{valueAtmemofhodgeFille}{\href{https://github.com/lsernau/grh-lean/blob/72fead102f663d5c7a407964a13dc0cf9d977207/grh-compact/src/New/Lift/ExponentLattices.lean\#L243-L252}{{\leanfont Lift.\allowbreak{}KummerChart.\allowbreak{}JumpLattice.\allowbreak{}valueAt\_mem\_of\_hodgeFil\_le}}}{ExponentLattices.lean}{243--252}
\crow{smulsubmemzTauUnitsmulofdeep}{\href{https://github.com/lsernau/grh-lean/blob/72fead102f663d5c7a407964a13dc0cf9d977207/grh-compact/src/New/Lift/ExponentLattices.lean\#L256-L258}{{\leanfont Lift.\allowbreak{}KummerChart.\allowbreak{}JumpLattice.\allowbreak{}smul\_sub\_mem\_zTauUnit\_smul\_of\_deep}}}{ExponentLattices.lean}{256--258}
\crow{locKisinPart}{\href{https://github.com/lsernau/grh-lean/blob/72fead102f663d5c7a407964a13dc0cf9d977207/grh-compact/src/New/Lift/HodgeExpansion.lean\#L23-L23}{{\leanfont Lift.\allowbreak{}KummerChart.\allowbreak{}JumpLattice.\allowbreak{}locKisinPart}}}{HodgeExpansion.lean}{23--23}
\crow{invariantsMod}{\href{https://github.com/lsernau/grh-lean/blob/72fead102f663d5c7a407964a13dc0cf9d977207/grh-compact/src/New/Lift/HodgeExpansion.lean\#L27-L28}{{\leanfont Lift.\allowbreak{}KummerChart.\allowbreak{}JumpLattice.\allowbreak{}invariantsMod}}}{HodgeExpansion.lean}{27--28}
\crow{costleofisUnramified}{\href{https://github.com/lsernau/grh-lean/blob/72fead102f663d5c7a407964a13dc0cf9d977207/grh-compact/src/New/Lift/HodgeExpansion.lean\#L30-L31}{{\leanfont Lift.\allowbreak{}KummerChart.\allowbreak{}JumpLattice.\allowbreak{}cost\_le\_of\_isUnramified}}}{HodgeExpansion.lean}{30--31}
\crow{improvedExponent}{\href{https://github.com/lsernau/grh-lean/blob/72fead102f663d5c7a407964a13dc0cf9d977207/grh-compact/src/New/Lift/HodgeExpansion.lean\#L35-L36}{{\leanfont Lift.\allowbreak{}KummerChart.\allowbreak{}JumpLattice.\allowbreak{}improvedExponent}}}{HodgeExpansion.lean}{35--36}
\crow{memwithExponentsofdeep}{\href{https://github.com/lsernau/grh-lean/blob/72fead102f663d5c7a407964a13dc0cf9d977207/grh-compact/src/New/Lift/HodgeExpansion.lean\#L41-L46}{{\leanfont Lift.\allowbreak{}KummerChart.\allowbreak{}JumpLattice.\allowbreak{}mem\_withExponents\_of\_deep}}}{HodgeExpansion.lean}{41--46}
\crow{galMonoidsmullocalFieldOf}{\href{https://github.com/lsernau/grh-lean/blob/72fead102f663d5c7a407964a13dc0cf9d977207/grh-compact/src/New/Lift/HodgeExpansion.lean\#L48-L49}{{\leanfont Lift.\allowbreak{}KummerChart.\allowbreak{}JumpLattice.\allowbreak{}galMonoid\_smul\_localFieldOf}}}{HodgeExpansion.lean}{48--49}
\crow{existsreachesdepthsucc}{\href{https://github.com/lsernau/grh-lean/blob/72fead102f663d5c7a407964a13dc0cf9d977207/grh-compact/src/New/Lift/HodgeExpansion.lean\#L54-L74}{{\leanfont Lift.\allowbreak{}KummerChart.\allowbreak{}JumpLattice.\allowbreak{}exists\_reaches\_depth\_succ}}}{HodgeExpansion.lean}{54--74}
\crow{wittPairinga}{\href{https://github.com/lsernau/grh-lean/blob/72fead102f663d5c7a407964a13dc0cf9d977207/grh-compact/src/New/Lift/HodgeExpansion.lean\#L82-L83}{{\leanfont Lift.\allowbreak{}KummerChart.\allowbreak{}JumpLattice.\allowbreak{}wittPairing}}}{HodgeExpansion.lean}{82--83}
\crow{wittPairingletterProjleft}{\href{https://github.com/lsernau/grh-lean/blob/72fead102f663d5c7a407964a13dc0cf9d977207/grh-compact/src/New/Lift/HodgeExpansion.lean\#L87-L88}{{\leanfont Lift.\allowbreak{}KummerChart.\allowbreak{}JumpLattice.\allowbreak{}wittPairing\_letterProj\_left}}}{HodgeExpansion.lean}{87--88}
\crow{wittPairingmemsaturateBy}{\href{https://github.com/lsernau/grh-lean/blob/72fead102f663d5c7a407964a13dc0cf9d977207/grh-compact/src/New/Lift/HodgeExpansion.lean\#L90-L92}{{\leanfont Lift.\allowbreak{}KummerChart.\allowbreak{}JumpLattice.\allowbreak{}wittPairing\_mem\_saturateBy}}}{HodgeExpansion.lean}{90--92}
\crow{frobExponent}{\href{https://github.com/lsernau/grh-lean/blob/72fead102f663d5c7a407964a13dc0cf9d977207/grh-compact/src/New/Lift/HodgeExpansion.lean\#L96-L98}{{\leanfont Lift.\allowbreak{}KummerChart.\allowbreak{}JumpLattice.\allowbreak{}frobExponent}}}{HodgeExpansion.lean}{96--98}
\crow{memhodgeFilofhodgebasis}{\href{https://github.com/lsernau/grh-lean/blob/72fead102f663d5c7a407964a13dc0cf9d977207/grh-compact/src/New/Lift/HodgeExpansion.lean\#L102-L118}{{\leanfont Lift.\allowbreak{}KummerChart.\allowbreak{}JumpLattice.\allowbreak{}mem\_hodgeFil\_of\_hodge\_basis}}}{HodgeExpansion.lean}{102--118}
\crow{valueAtcoordRetractsubmemhodgeFil}{\href{https://github.com/lsernau/grh-lean/blob/72fead102f663d5c7a407964a13dc0cf9d977207/grh-compact/src/New/Lift/HodgeExpansion.lean\#L120-L125}{{\leanfont Lift.\allowbreak{}KummerChart.\allowbreak{}JumpLattice.\allowbreak{}valueAt\_coordRetract\_sub\_mem\_hodgeFil}}}{HodgeExpansion.lean}{120--125}
\crow{coordRetractmemwithExponents}{\href{https://github.com/lsernau/grh-lean/blob/72fead102f663d5c7a407964a13dc0cf9d977207/grh-compact/src/New/Lift/HodgeExpansion.lean\#L127-L132}{{\leanfont Lift.\allowbreak{}KummerChart.\allowbreak{}JumpLattice.\allowbreak{}coordRetract\_mem\_withExponents}}}{HodgeExpansion.lean}{127--132}
\crow{phipowmemwithExponents}{\href{https://github.com/lsernau/grh-lean/blob/72fead102f663d5c7a407964a13dc0cf9d977207/grh-compact/src/New/Lift/HodgeExpansion.lean\#L134-L143}{{\leanfont Lift.\allowbreak{}KummerChart.\allowbreak{}JumpLattice.\allowbreak{}phi\_pow\_mem\_withExponents}}}{HodgeExpansion.lean}{134--143}
\crow{memwithExponentsstdExponentiff}{\href{https://github.com/lsernau/grh-lean/blob/72fead102f663d5c7a407964a13dc0cf9d977207/grh-compact/src/New/Lift/HodgeExpansion.lean\#L148-L174}{{\leanfont Lift.\allowbreak{}KummerChart.\allowbreak{}JumpLattice.\allowbreak{}mem\_withExponents\_stdExponent\_iff}}}{HodgeExpansion.lean}{148--174}
\crow{letterProjmemwithExponents}{\href{https://github.com/lsernau/grh-lean/blob/72fead102f663d5c7a407964a13dc0cf9d977207/grh-compact/src/New/Lift/HodgeExpansion.lean\#L176-L179}{{\leanfont Lift.\allowbreak{}KummerChart.\allowbreak{}JumpLattice.\allowbreak{}letterProj\_mem\_withExponents}}}{HodgeExpansion.lean}{176--179}
\crow{letterProjmemofvalueFieldPoints}{\href{https://github.com/lsernau/grh-lean/blob/72fead102f663d5c7a407964a13dc0cf9d977207/grh-compact/src/New/Lift/HodgeExpansion.lean\#L181-L183}{{\leanfont Lift.\allowbreak{}KummerChart.\allowbreak{}JumpLattice.\allowbreak{}letterProj\_mem\_of\_valueFieldPoints}}}{HodgeExpansion.lean}{181--183}
\crow{wittPairingmemwithExponents}{\href{https://github.com/lsernau/grh-lean/blob/72fead102f663d5c7a407964a13dc0cf9d977207/grh-compact/src/New/Lift/HodgeExpansion.lean\#L188-L207}{{\leanfont Lift.\allowbreak{}KummerChart.\allowbreak{}JumpLattice.\allowbreak{}wittPairing\_mem\_withExponents}}}{HodgeExpansion.lean}{188--207}}
% the notations of the defined functions: one slot for each explicit argument of the Lean declaration
\newcommand{\JumpWin}[5]{#5\subseteq[-#1,#1]\cup#2\cdot[#3,#3+#4]}
\newcommand{\raisedE}[3]{\operatorname{rs}_{#1}(#2,#3)}
\newcommand{\twoWin}[2]{\mathcal{W}_{#1}(#2)}
\newcommand{\NoMult}[3]{\operatorname{NM}(#1;#2,#3)}
\newcommand{\constAt}[3]{\operatorname{cst}_{#1,#2}(#3)}
\newcommand{\liftId}[4]{\mathfrak{a}_{#1}(#2;#3,#4)}
\newcommand{\Kills}[4]{\operatorname{Kill}_{#1}(#2;#3,#4)}
\newcommand{\extL}[1]{\tilde{Y}_{#1}}
\newcommand{\trivL}[1]{C^{\mathrm{tr}}_{#1}}
\newcommand{\projC}[1]{\pr^{C}_{#1}}
\newcommand{\cost}[2]{\operatorname{cost}(#1,#2)}
\newcommand{\locLP}[2]{\mathfrak{N}_{#1,#2}}
\newcommand{\locKP}[1]{(#1)_{\fS}}
\newcommand{\invMod}[1]{\operatorname{Inv}(#1)}
\newcommand{\impE}[4]{\operatorname{imp}(#1,#2)(#3,#4)}
\newcommand{\wPair}[3]{{#1}_{W}(#2,#3)}
\newcommand{\frobX}[5]{\operatorname{fr}_{#1,#2}(#3)(#4,#5)}
\newcommand{\UnitMembers}{141}
\newcommand{\UnitInNote}{77}
\newcommand{\UnitProofs}{69}
\newcommand{\UnitSteps}{375}
\newcommand{\UnitAllSteps}{430}


\linespread{1.08}                       % 11 pt on 14.7 pt
\setlength{\parindent}{0pt}
\setlength{\parskip}{5.5pt plus 1.5pt minus 0.5pt}
\setlength{\abovedisplayskip}{7.5pt plus 2pt minus 3pt}
\setlength{\belowdisplayskip}{7.5pt plus 2pt minus 3pt}
\setlength{\abovedisplayshortskip}{3pt plus 2pt}
\setlength{\belowdisplayshortskip}{6pt plus 2pt minus 2pt}
\emergencystretch=1.5em
\clubpenalty=10000 \widowpenalty=10000 \displaywidowpenalty=10000
\relpenalty=7000 \binoppenalty=9000
\raggedbottom
\pagestyle{fancy}\fancyhf{}\renewcommand{\headrulewidth}{0pt}
\fancyfoot[C]{\sffamily\fontsize{7.6}{9}\selectfont\color{inkfaint}\thepage\,/\,\pageref*{LastPage}}
\setlist{topsep=2pt,partopsep=0pt,parsep=0pt,itemsep=4pt,leftmargin=2.1em,labelsep=0.55em}

\newcommand{\Q}{\mathbb{Q}}
\newcommand{\Z}{\mathbb{Z}}
\newcommand{\F}{\mathbb{F}}
\newcommand{\len}{\operatorname{len}}
\newcommand{\cok}{\operatorname{cok}}
\newcommand{\End}{\operatorname{End}}
\newcommand{\Hom}{\operatorname{Hom}}
\newcommand{\Gal}{\operatorname{Gal}}
\newcommand{\diag}{\operatorname{diag}}
\newcommand{\kw}[1]{\emph{#1}}
\newcommand{\m}[1]{\mbox{$#1$}}        % a short formula that must not be broken across lines
\newcommand{\head}[1]{\par\addvspace{17pt plus 4pt minus 2pt}\Needspace{6\baselineskip}%
  \noindent{\fontsize{13.2}{16}\selectfont\bfseries\color{ink}#1}\par\nopagebreak\vspace{-1pt}}
\newcommand{\sect}[2]{\refstepcounter{section}\label{#2}\par\addvspace{19pt plus 4pt minus 2pt}\Needspace{7\baselineskip}%
  \noindent\llap{{\fontsize{13.2}{16}\selectfont\bfseries\color{inkfaint}\thesection}\hspace{9pt}}%
  {\fontsize{13.2}{16}\selectfont\bfseries\color{ink}#1}\par\nopagebreak\vspace{-1pt}}
\newcommand{\subhead}[1]{\par\addvspace{9pt plus 2pt}\Needspace{4\baselineskip}\noindent\textbf{#1}\enspace\ignorespaces}
% where a statement lives in the Lean archive
\newcommand{\leantag}{{\sffamily\bfseries\fontsize{6.3}{8}\selectfont\color{leancol}\addfontfeatures{LetterSpace=7}LEAN}}
\newcommand{\sep}{\unskip\nobreak\hspace{5pt}{\color{inkfaint}\textperiodcentered}\hspace{5pt plus 2pt}\ignorespaces}
\newcommand{\leanline}[1]{\par\nopagebreak\vspace{-2.5pt}{\fontsize{8.9}{12.4}\selectfont\color{inksoft}\raggedright
  \noindent\llap{\leantag\hspace{9pt}}#1\par}\vspace{1pt}}
% a displayed statement in a tinted box: #1 colour, #2 text
\newcommand{\statement}[2]{\par\addvspace{7pt}\begingroup\setlength{\fboxsep}{8.5pt}\setlength{\fboxrule}{0.6pt}%
  \noindent\fcolorbox{#1!30!white}{#1!3!white}{\parbox{\dimexpr\linewidth-2\fboxsep-2\fboxrule\relax}{#2}}%
  \endgroup\par\addvspace{7pt}}
\newcommand{\nt}[2]{\par\noindent\hangindent=1.1em\hangafter=1 #1\enspace{\color{inksoft}#2}\par\vspace{1pt}}
\newcommand{\legdot}[1]{\tikz[baseline=-0.62ex]\fill[#1] (0,0) circle (2.5pt);\hspace{3pt}}
\newcommand{\crow}[4]{\csname CS@#1\endcsname & #2 & \textcolor{inksoft}{#3} & \textcolor{inksoft}{#4}\\}
\newcommand{\leanquoteline}[2]{\par\noindent\hangindent=\dimexpr#1\fontdimen2\font+2.2em\relax\hangafter=1
  \hspace*{#1\fontdimen2\font}#2}
\newcommand{\quoted}[1]{\par\vspace{3pt}{\ttfamily\fontsize{7.5}{10.4}\selectfont\color{leancol}\raggedright
  \setlength{\parskip}{0pt}\leanquote{#1}\par}\vspace{4pt}}
\newcommand{\bx}[3]{{\sffamily\bfseries\fontsize{7.3}{9.2}\selectfont\color{#1!82!black}#2}\\[1.5pt]#3}
\newcommand{\cs}[2]{\expandafter\newcommand\csname CS@#1\endcsname{#2}}

% theorems
\newtheoremstyle{plainx}{7pt plus 2pt}{7pt plus 2pt}{\itshape}{}{\bfseries}{.}{0.5em}{}
\newtheoremstyle{uprightx}{7pt plus 2pt}{7pt plus 2pt}{}{}{\bfseries}{.}{0.5em}{}
\theoremstyle{plainx}
\newtheorem{theorem}{Theorem}[section]
\newtheorem{lemma}[theorem]{Lemma}
\newtheorem{proposition}[theorem]{Proposition}
\newtheorem{corollary}[theorem]{Corollary}
\theoremstyle{uprightx}
\newtheorem{definition}[theorem]{Definition}
\newtheorem{setting}[theorem]{Setting}          % hypotheses shared by several statements; each of them cites it by number
\newtheorem{remark}[theorem]{Remark}
\newtheorem{example}[theorem]{Example}
\renewcommand{\qedsymbol}{{\setlength{\fboxsep}{0pt}\setlength{\fboxrule}{0.5pt}\fbox{\rule{0pt}{0.46em}\hspace{0.46em}}}}
% what holds throughout the note: the ambient objects only (implicit and instance arguments), never a hypothesis
\newenvironment{conventions}{\par\addvspace{5pt}\noindent\textbf{Conventions.}\enspace\ignorespaces}{\par\addvspace{3pt}}
\makeatletter
\renewenvironment{proof}[1][Proof]{\par\pushQED{\qed}\normalfont\topsep0pt\partopsep0pt
  \trivlist\item[\hskip\labelsep\itshape #1.]\ignorespaces}{\popQED\endtrivlist\@endpefalse}
\makeatother

% where each declaration of the concordance is treated: \cs{key}{place}
\cs{existsdepthstep}{Theorem A}
\cs{memwithExponentsstdExponentiff}{Theorem B}
\cs{wittPairingmemwithExponents}{Theorem C}
\cs{killsCokernelofisUnramified}{Theorem D}
\cs{cost}{Definition~\ref{d:cost}}
\cs{withExponents}{Definition~\ref{d:exp}}
\cs{locLetterPart}{Definition~\ref{d:exp}}
\cs{locKisinPart}{Definition~\ref{d:ops}}
\cs{invariantsMod}{Definition~\ref{d:ops}}
\cs{improvedExponent}{Definition~\ref{d:ops}}
\cs{frobExponent}{Definition~\ref{d:ops}}
\cs{wittPairinga}{Definition~\ref{d:ops}}
\cs{PairingSide}{Definition~\ref{d:side}}
\cs{extLattice}{Definition~\ref{d:side}}
\cs{trivLattice}{Definition~\ref{d:side}}
\cs{projTriv}{Definition~\ref{d:side}}
\cs{IsInJumpWindow}{Definition~\ref{d:win}}
\cs{raisedAtE}{Definition~\ref{d:win}}
\cs{twoWindows}{Definition~\ref{d:win}}
\cs{NoMultipleNear}{Definition~\ref{d:win}}
\cs{constantAt}{Definition~\ref{d:lift}}
\cs{liftIdeal}{Definition~\ref{d:lift}}
\cs{KillsCokernel}{Definition~\ref{d:lift}}
\cs{existsdepthstepzero}{Lemma~\ref{l:zero}}
\cs{coordRetract}{Definition~\ref{d:ret}}
\cs{existscoordRetract}{Lemma~\ref{l:ret}}
\cs{smulsubmemzTaukisinLocAinf}{Lemma~\ref{l:ztau}}
\cs{smulsubmemzTauUnitsmulofdeep}{Lemma~\ref{l:ztau}}
\cs{valueAtmemofhodgeFille}{Lemma~\ref{l:any}}
\cs{costleofisUnramified}{Lemma~\ref{l:small}}
\cs{galMonoidsmullocalFieldOf}{Lemma~\ref{l:small}}
\cs{memwithExponentsofdeep}{Lemma~\ref{l:deep}}
\cs{existsreachesdepthsucc}{Proposition~\ref{p:succ}}
\cs{memhodgeFilofhodgebasis}{Lemma~\ref{l:basis}}
\cs{phipowmemwithExponents}{Lemma~\ref{l:shift}}
\cs{letterProjmemwithExponents}{Lemma~\ref{l:proj}}
\cs{letterProjmemofvalueFieldPoints}{Lemma~\ref{l:proj}}
\cs{wittPairingletterProjleft}{Lemma~\ref{l:pair}}
\cs{wittPairingmemsaturateBy}{Lemma~\ref{l:pair}}
\cs{sumtwoWindowsle}{Lemma~\ref{l:count}}
\cs{taupowsmulmemainfLatticeofchi}{Lemma~\ref{l:chi}}
\cs{existsliftofmemliftIdeal}{Lemma~\ref{l:liftS}}
\cs{memliftIdealofderivation}{Lemma~\ref{l:der}}
\cs{memliftIdealzeroofmemone}{Lemma~\ref{l:down}}
\cs{memliftIdealofmulEpow}{Lemma~\ref{l:down}}
\cs{existszpowmemliftIdeal}{Lemma~\ref{l:zpow}}
\cs{frobEsmulmemkisinLocAinfiff}{Lemma~\ref{l:loc}}
\cs{frobEUnitzpowsmulmemkisinLoc}{Lemma~\ref{l:loc}}
\cs{wittMapmemkisinLocAinf}{Lemma~\ref{l:loc}}
\cs{wittMapmemkisinLoc}{Lemma~\ref{l:loc}}
\cs{galsmulmemkisinLocAinf}{Lemma~\ref{l:loc}}
\cs{kisinLocAinflefracSpan}{Lemma~\ref{l:loc}}
\cs{valueAtsmulsub}{Lemma~\ref{l:loc}}
\cs{phipowmemkisinLoc}{Lemma~\ref{l:loc}}
\cs{spankisinLoc}{Lemma~\ref{l:loc}}
\cs{valueAtmemvalueField}{Lemma~\ref{l:loc}}
\cs{valueAtzero}{Lemma~\ref{l:loc}}
\cs{memwithExponentsofmemhodgeFil}{Lemma~\ref{l:exp}}
\cs{memwithExponentsofle}{Lemma~\ref{l:exp}}
\cs{valueAtmemhodgeFilofmemainfLattice}{Lemma~\ref{l:exp}}
\cs{mapmemwithExponents}{Lemma~\ref{l:exp}}
\cs{memwithExponentsofsmul}{Lemma~\ref{l:exp}}
\cs{frobEUnitzpowmemwithExponents}{Lemma~\ref{l:exp}}
\cs{smulsubmemwithExponents}{Lemma~\ref{l:exp}}
\cs{coordRetractsum}{Lemma~\ref{l:retprop}}
\cs{coordRetractofmemkisinLoc}{Lemma~\ref{l:retprop}}
\cs{mapcoordRetract}{Lemma~\ref{l:retprop}}
\cs{valueAtcoordRetractsubmemhodgeFil}{Lemma~\ref{l:retprop}}
\cs{coordRetractmemwithExponents}{Lemma~\ref{l:retprop}}
\cs{memliftIdealiff}{Lemma~\ref{l:ideal}}
\cs{ppowmemliftIdealself}{Lemma~\ref{l:ideal}}
\cs{phipowsmulconstantAt}{Lemma~\ref{l:ideal}}
\cs{mulEpowmemliftIdealone}{Lemma~\ref{l:ideal}}
\cs{liftIdealleofppowmem}{Lemma~\ref{l:ideal}}
\cs{constantAtmemlocLetterPart}{Lemma~\ref{l:ideal}}
\cs{phipowmemliftIdealiff}{Lemma~\ref{l:ideal}}
\cs{ppowmulzpowmemliftIdealiff}{Lemma~\ref{l:ideal}}
% ---------- common symbols of the notes on unit 136
\providecommand{\Zp}{\mathbb{Z}_p}\providecommand{\Qp}{\mathbb{Q}_p}\providecommand{\GQp}{G_{\mathbb{Q}_p}}
\providecommand{\Ainf}{A_{\inf}}\providecommand{\WC}{W(C^{\flat})}\providecommand{\BdR}{B_{\mathrm{dR}}}
\providecommand{\Fil}{\operatorname{Fil}}\providecommand{\fS}{\mathfrak{S}}\providecommand{\cE}{\mathcal{E}}
\providecommand{\cO}{\mathcal{O}}\providecommand{\fc}{\mathfrak{c}}\providecommand{\fd}{\mathfrak{d}}
\providecommand{\WW}{\mathbb{W}}\providecommand{\Per}{\operatorname{Per}}\providecommand{\pr}{\operatorname{pr}}
\providecommand{\Minf}{\mathfrak{M}^{\inf}}\providecommand{\Mloc}{\mathfrak{M}_{\mathrm{loc}}}
\providecommand{\Mlocinf}{\mathfrak{M}^{\inf}_{\mathrm{loc}}}\providecommand{\std}{\operatorname{std}}
\providecommand{\ret}{\operatorname{ret}}
\providecommand{\inputbox}[1]{\statement{inksoft}{#1}}
% the exponent module (the spec says why it has no generated macro): lattice datum, letter, exponent function
\providecommand{\wExp}[3]{\mathfrak{X}_{#1}(#2;#3)}

\begin{document}
\color{ink}

{\raggedright
{\sffamily\bfseries\fontsize{7.2}{9}\selectfont\color{leancol}\addfontfeatures{LetterSpace=8}READBACK NOTE · CHAPTER VIII · UNIT 136 · PART 1 OF 3\par}
\vspace{5pt}
{\fontsize{21.5}{25.5}\selectfont\bfseries Exponents of elements of Kisin modules, the cost of deepening them, and cokernel killers\par}
\vspace{5pt}
{\fontsize{12.4}{16}\selectfont\itshape\color{inksoft} One unit of the Lean artifact \ArtifactName, read back into
mathematics\par}}
\vspace{8pt}
{\color{rulegrey}\hrule height 0.6pt}
\vspace{7pt}
\thispagestyle{empty}

\textbf{Summary.} The last step of the choice of the twists bounds the loss that the conditions at~$\ell$ cause in
Selmer groups, by lifting crystalline extension classes in large weight. Its tool is local and $p$-adic: the
Breuil--Kisin module of a crystalline lattice, seen inside \m{\WC\otimes L}. This note measures an element of that
module, with all Frobenius translates \m{\cE_\nu=\varphi^\nu(\cE)} of the Eisenstein element inverted, by an
\emph{exponent} at each~$\cE_\nu$ and each Hodge--Tate weight. Three things are proved. A Galois operator
\m{(g-1)/(z\tau)} raises the exponent at~$\cE_\nu$ by one at the price of an explicit power of~$p$ (the
\emph{cost}). The standard exponents of level~$l$ describe exactly the $l$-th Frobenius pullback of the Kisin
module, and exponents add under equivariant pairings. Finally, for an extension of a trivial lattice by a
crystalline one over an unramified base, \m{p^{\kappa}\cE^{\rho}} kills the cokernel of the map of Kisin modules as
soon as all weights are at most \m{(p-3)(\kappa+\rho)}.

\statement{leancol}{\fontsize{9.6}{13.2}\selectfont\textbf{What this note is.} The artifact \ArtifactName{} is a
body of Lean code: the compact edition of a formal proof, due to L.~Alpöge with the AI system Claude, of the Riemann
hypothesis and of its generalisation to Dirichlet $L$-functions. That headline is a claim. The authors report that
every file compiles with Lean v4.35.0 against a pinned version of Mathlib and that Lean's kernel accepted the
proofs; an independent build is recorded in the handout's record of checks when it reports. No mathematician has
read the argument through, and nothing here says that the Riemann hypothesis is proved. This note is the first of
three on one unit of the artifact, the folder on lifting; it treats \UnitInNote{} of the unit's \UnitMembers{}
declarations: three files (exponent lattices, the Hodge expansion, cokernel killers) and four definitions of a
fourth. Theorems A to~D and every lemma below are proved from the Inputs of Section~\ref{s:setting}, which belong to
the background library and are stated and not proved here; the main one is the frame theorem for the Kisin module
of a crystalline lattice. Every statement was compared with its Lean text, argument by argument. The Lean text read
is \PinText; Appendix~A is the concordance of \ConcDecls{} declarations.}

\sect{The statements}{s:intro}
The objects are those of Settings \ref{set:rings} and~\ref{set:lattice}: a prime~$p$, a group~$\Gamma$ over the
absolute Galois group of~$\Qp$, a Kummer chart with ring \m{\fS\subset\Ainf}, coordinate~$z$ and Eisenstein
element~$\cE$, and lattice data~$L$ with their periods. The exponent module \m{\wExp{L}{i}{o}} is
Definition~\ref{d:exp}, the cost \m{\cost{\nu}{n}} is Definition~\ref{d:cost}.

\statement{arith}{\phantomsection\label{t:A}\textbf{Theorem A} (the depth step).\enspace
\emph{In Setting~\ref{set:rings}, let \m{\nu\ge0} be an integer and let~$n$ be an integer with \m{n\ne0}. Then
there are \m{g\in\Gamma'}, a unit~$w$ of~$\Ainf$ and elements \m{y,X\in\Ainf} such that
\[\begin{gathered}g\cdot\cE_\nu=w\,\cE_\nu,\qquad \cE_\nu\mid g\cdot s-s\ \text{ in $\Ainf$ for every }s\in\fS,\\
\cE_\nu\mid y\,X-p^{\cost{\nu}{n}},\qquad z\tau\cdot X=w^{n}-1 .\end{gathered}\]}}
\leanline{\lean{existsdepthstep}}

\statement{arith}{\phantomsection\label{t:B}\textbf{Theorem B} (standard exponents are Frobenius pullbacks).\enspace
\emph{In Setting~\ref{set:lattice}, let $x\in\WC\otimes L$, let~$i$ be a letter and let \m{l\ge0} be an integer.
Then \m{x\in\wExp{L}{i}{\std_l}} if and only if \m{\varphi^{-l}x\in\Minf(L)} and \m{x\in(\WC\otimes L)_i}.}}
\leanline{\lean{memwithExponentsstdExponentiff}}

\statement{arith}{\phantomsection\label{t:C}\textbf{Theorem C} (exponents add).\enspace
\emph{In Setting~\ref{set:lattice}, let $A$, $L$, $C$ be lattice data over~$\cO$ and let \m{b\colon A\times L\to C}
be an $\cO$-bilinear map of their lattices with \m{b(g a,g t)=g\,b(a,t)} for all \m{g\in\Gamma}, \m{a\in A},
\m{t\in L}. Let~$i$ be a letter, let $o_A$, $o$, $o_C$ be exponent functions, let \m{y\in\wExp{A}{i}{o_A}} and
\m{x\in\wExp{L}{i}{o}}, and suppose that
\[o_C(\nu,v)\ \le\ o(\nu,u')+o_A(\nu,u)\]
for every \m{\nu\ge0} and all \m{u\in U_A(i^{p^{-\nu}})}, \m{u'\in U_L(i^{p^{-\nu}})}, \m{v\in U_C(i^{p^{-\nu}})} with
\m{u+u'\le v}. Then $\wPair{b}{y}{x}\in\wExp{C}{i}{o_C}$.}}
\leanline{\lean{wittPairingmemwithExponents}}

\statement{arith}{\phantomsection\label{t:D}\textbf{Theorem D} (a cokernel killer over an unramified chart).\enspace
\emph{In Setting~\ref{set:lattice}, suppose that the Kummer chart is unramified and that \m{p\ge5}. Let $S$ be a
side over~$\cO$, let $i$ be a letter and let \m{\kappa,\rho\ge0} be integers such that
\m{u\le(p-3)(\kappa+\rho)} for every letter~$j$ and every \m{u\in U_{Y_S}(j)}. Then \m{\Kills{S}{i}{\kappa}{\rho}}:
the element \m{p^{\kappa}\cE^{\rho}} multiplies the part at~$i$ of \m{\Ainf\otimes C_S} into the image of
\m{\wExp{\extL{S}}{i}{\std_0}}.}}
\leanline{\lean{killsCokernelofisUnramified}}

Theorem~A is the arithmetic heart: the cost \m{v_p(n)+\fd+3\nu} is what one Galois step costs. Theorem~B is the
bridge between the exponents and Kisin's module, and makes the definition of the exponents usable. Theorem~C is
what lets Part~2 pair a class with a dual class. Theorem~D is the first source of cokernel killers; Part~3 has a
second, for tame types.

\sect{The setting}{s:setting}
\begin{conventions}\phantomsection\label{c:conventions}
Throughout, $p$ is a prime number. The declarations of this unit take the following ambient objects as implicit or
instance arguments; they are fixed by Settings \ref{set:rings} and~\ref{set:lattice} and named again in each
statement that needs them: the prime~$p$; the group~$\Gamma$ over~$\GQp$; the Kummer chart~$\fc$ for~$\Gamma$; the
coefficient ring~$\cO$ with its root of unity~$\omega$; lattice data $L$, $L'$, $A$, $C$ over~$\cO$; a side~$S$; a
letter~$i$; integers $l$, $\nu$, $n$, $\kappa$, $r$; exponent functions $o$, $o'$; elements $x$, $y$ of
\m{\WC\otimes L}; a period~$\mu$; an element~$g$ of~$\Gamma'$. Rings are commutative with~$1$. For an
integer~$n$, \m{v_p(n)} is the exponent of~$p$ in~$n$, with the convention \m{v_p(0)=0}. A unit~$w$ of a ring
has all integral powers \m{w^{n}}, \m{n\in\mathbb{Z}}. Tensor products without a subscript are over~$\Zp$. An
\kw{exponent function} is a function $o\colon\mathbb{N}\times\mathbb{Z}\to\mathbb{Z}$, written
\m{(\nu,u)\mapsto o(\nu,u)}; inequalities between exponent functions are meant pointwise. The Lean text calls
the artifact's period ring \emph{Bcris}; it is Colmez's ring $B_{\max}$, and this note writes $B_{\max}$.
\end{conventions}

\begin{setting}[period rings and a Kummer chart; from the background library]\label{set:rings}
\emph{Period rings.} $C$ is the completion of an algebraic closure~$\bar{\mathbb{Q}}_p$ of~$\Qp$, with ring of
integers~$\mathcal{O}_C$; $\GQp=\mathrm{Gal}(\bar{\mathbb{Q}}_p/\Qp)$. $\mathcal{O}_C^\flat$ is the tilt
of~$\mathcal{O}_C$, with fraction field~$C^\flat$; $\Ainf=W(\mathcal{O}_C^\flat)\subset\WC$ are the rings of Witt
vectors, with Frobenius automorphism~$\varphi$ and the action of~$\GQp$; $[\,\cdot\,]$ is the Teichmüller lift;
$\theta\colon\Ainf\to\mathcal{O}_C$ is Fontaine's map. Fix $\varepsilon\in\mathcal{O}_C^\flat$, a compatible
system of $p$-power roots of unity whose component of level one is a primitive $p$-th root of unity, and put
\[\tau=[\varepsilon^{1/p}]-1,\qquad \xi=\sum_{0\le a<p}[\varepsilon^{1/p}]^{a},\qquad\text{so that}\quad
\xi\,\tau=[\varepsilon]-1 ;\]
$\xi$ generates the kernel of~$\theta$. $B^+_{\mathrm{dR}}$ is the completion of \m{\Ainf[1/p]} along the kernel
of~$\theta$, a complete discrete valuation ring with residue field~$C$, and $\BdR$ is its fraction field; for
\m{n\in\mathbb{Z}}, \m{\Fil^{n}\BdR} is the $n$-th power of the maximal ideal of $B^+_{\mathrm{dR}}$, so that
\m{\Fil^{0}\BdR=B^+_{\mathrm{dR}}}. The group~$\GQp$ acts on~$\BdR$ and preserves each \m{\Fil^{n}\BdR}. For
\m{\nu\ge0}, the composite of $\varphi^{-\nu}\colon\Ainf\to\Ainf$ with \m{\Ainf\to\BdR} is again written
$\varphi^{-\nu}\colon\Ainf\to\BdR$, and $\theta_\nu=\theta\circ\varphi^{-\nu}\colon\Ainf\to\mathcal{O}_C$.

\emph{The group.} $\Gamma$ is a group with a homomorphism \m{j_\Gamma\colon\Gamma\to\GQp} whose image is open; it
acts on the rings above through~$j_\Gamma$, and \m{\chi\colon\Gamma\to\Zp^\times} is the cyclotomic character
composed with~$j_\Gamma$. A property holds \kw{near~$1$} if it holds for all \m{g\in\Gamma} with \m{j_\Gamma(g)} in
some neighbourhood of~$1$. The group \m{\mathbb{Z}\times\Gamma} acts on~$\WC$ by
$(a,g)\cdot c=\varphi^{a}(j_\Gamma(g)c)$; write $\varphi$ for \m{(1,1)} and $g$ for \m{(0,g)}.

\emph{The chart.} A \kw{chart} consists of an integer \m{h\ge1}, the ring \m{\WW=W(\mathbb{F}_{p^h})} of Witt
vectors of the field $\mathbb{F}_{p^h}=\{c\in\mathcal{O}_C^\flat: c^{p^h}=c\}$, a homomorphism of
$\WW$-algebras \m{\WW[[X]]\to\Ainf} that sends~$X$ to a Teichmüller lift~$z$ which is not a unit, a natural
number~$e_{\fc}$ and a unit~$\mathsf{A}$ of \m{\WW[[X]]} such that the image~$\cE$ of
\m{X^{e_{\fc}}\mathsf{A}+p} is~$\xi$ times a unit of~$\Ainf$. Its ring~$\fS$ is the image of \m{\WW[[X]]}; it
contains $z$ and~$\cE$, is stable under~$\varphi$, and $\cE_\nu=\varphi^{\nu}(\cE)\in\fS$ for \m{\nu\ge0}. The
integer $h$ is the \kw{residue degree}, $e_{\fc}$ the \kw{index}, and
\[\fd=0\ \text{ if }e_{\fc}=1,\qquad \fd=v_p(e_{\fc})+1\ \text{ otherwise}.\]
A \kw{Kummer chart}~$\fc$ for~$\Gamma$ is a chart together with a subgroup~$\Gamma'$ of~$\Gamma$, of index prime
to~$p$, that contains all elements near~$1$, such that \m{j_\Gamma(g)} fixes $\WW$ and~\m{\theta(z)} for every
\m{g\in\Gamma'}, every element of~$\GQp$ that fixes $\WW$ and \m{\theta(z)} lies in \m{j_\Gamma(\Gamma')}, and for
every integer \m{a\ge0} some \m{g\in\Gamma'} has \m{g\cdot z=[\varepsilon]^{a}z}. It is \kw{unramified} if
\m{\cE=z+p}, \m{e_{\fc}=1} and \m{\fd=0}. The field~$F$ is the subfield of~$\bar{\mathbb{Q}}_p$ fixed by
\m{j_\Gamma(\Gamma')}. The ring~$\fS_B$ is the $\fS$-algebra of the elements of \m{B^{+}=A_{\max}[1/p]} whose
multiplication commutes with every $\fS$-linear endomorphism of~$B^{+}$ that preserves the image of~$A_{\max}$; it
contains the image of~$\fS$. For \m{\nu\ge0}, $F_\nu$ is the subfield of~$\BdR$ generated over~$F$ by the products
of an element of~$\fS_B$ and an element of \m{\varphi^{-\nu}(\fS)}; so $\varphi^{-\nu}(\fS)\subseteq F_\nu$ and
\m{F_0\subseteq F_\nu}. Finally $\operatorname{co}_0\colon\Ainf\to\fS$ is a fixed $\fS$-linear map that is the
identity on~$\fS$, one of the coordinates \m{(\operatorname{co}_a)_a} of a fixed topological basis of~$\Ainf$
over~$\fS$ for the \m{(p,z)}-adic topology; an element of~$\Ainf$ all of whose coordinates vanish is~$0$.
\end{setting}

\begin{setting}[lattice data, periods and Kisin modules; from the background library]\label{set:lattice}
In Setting~\ref{set:rings}:

\emph{Letters and coefficients.} A \kw{letter} is an element of \m{\mathbb{F}_{p^h}^\times}; the letters are
permuted by \m{i\mapsto i^{p}}, and \m{i^{p^{-\nu}}} is the letter whose $p^{\nu}$-th power is~$i$. A
\kw{coefficient ring} is the ring of integers~$\cO$ of a finite extension of~$\Qp$ together with an element
\m{\omega\in\cO} with \m{\omega^{p^h-1}=1}.

\emph{Lattices.} A \kw{lattice} over~$\cO$ is a finite free $\cO$-module with an $\cO$-linear action
of~$\Gamma$. For a lattice~$T$, \m{\WC\otimes T} carries the action
$(a,g)(c\otimes t)=\varphi^{a}(j_\Gamma(g)c)\otimes gt$ of \m{\mathbb{Z}\times\Gamma}, and \m{\Ainf\otimes T}
is the $\Ainf$-submodule spanned by the \m{1\otimes t}. A $\Zp$-bilinear map \m{b\colon T\times T'\to T''} has
a $\WC$-bilinear extension to these modules. For a letter~$i$, \m{(\WC\otimes T)_i} is the $\Ainf$-submodule on
which \m{1\otimes\omega} acts as~$[i]$, and
\[\pr_i=\frac1{p^h-1}\sum_{0\le a<p^h-1}\bigl([i]^{-1}(1\otimes\omega)\bigr)^{a}\]
is the projector onto it. The lattice~$T$ is \kw{potentially crystalline} if the elements of
\m{B_{\max}\otimes T} fixed near~$1$ span it over~$B_{\max}$, and \kw{crystalline on~$\Gamma'$} if those fixed by
all of~$\Gamma'$ do. For a $\Zp$-algebra map \m{\gamma\colon\cO\to C} and \m{u\in\mathbb{Z}}, the
\kw{multiplicity of the Hodge--Tate jump~$u$ of~$T$ at~$\gamma$} is the dimension over~$\bar{\mathbb{Q}}_p$ of
the space of $\gamma$-semilinear maps \m{\lambda\colon T\to C} with
$\lambda(gt)=\chi(g)^{-u}\,g\lambda(t)$ for each~$t$ and all~$g$ near~$1$. \m{Z^{1}_{\mathrm{f}}(T)} is the
$\cO$-module of the maps \m{c\colon\Gamma\to T} for which some \m{y\in B_{\max}\otimes T} has
\m{1\otimes c(g)=g y-y} for all~$g$ (the crystalline cocycles).

\emph{Lattice data.} A \kw{lattice datum}~$L$ over~$\cO$ is a lattice, also written~$L$, that is crystalline
on~$\Gamma'$, together with a finite set \m{U_L(i)\subset\mathbb{Z}} for every letter~$i$, such that
\m{U_L(i)} contains every~$u$ that is a Hodge--Tate jump of~$L$ of non-zero multiplicity at some
\m{\gamma\colon\cO\to C} with \m{\gamma(\omega)=\theta([i])}. (The formal structure also records a jump set for
each~$\gamma$, contained in the corresponding~\m{U_L(i)}.) An \kw{equivariant map} \m{f\colon L\to L'} of lattice
data is a $\Gamma$-equivariant $\cO$-linear map of the lattices; \m{1\otimes f} is its extension
\m{\WC\otimes L\to\WC\otimes L'}. The \kw{jump bound} \m{h_L\ge0} is the least integer~$a$ such that all
Hodge--Tate jumps of~$L$, as a $\Zp$-representation of~$\Gamma'$, lie in \m{[-a,a]}.

\emph{Periods.} \m{\Per(L)} is the set of the $\Ainf$-linear maps \m{\mu\colon\WC\otimes L\to\BdR} with
$\mu(1\otimes gt)=g\cdot\mu(1\otimes t)$ for \m{t\in L}, \m{g\in\Gamma'}. For a letter~$i$ and
\m{u\in\mathbb{Z}}, \m{\Per_{i,u}(L)} is the set of the \m{\mu\in\Per(L)} with
$\mu(\pr_i(1\otimes t))\in\Fil^{-u}\BdR$ for all \m{t\in L} (the periods \kw{of type~$u$ at~$i$}). For
\m{\nu\ge0} the map \m{x\mapsto\mu(\varphi^{-\nu}x)} is semilinear along $\varphi^{-\nu}\colon\Ainf\to\BdR$.
The \kw{fraction span} of~$L$ is the set of the $x\in\WC\otimes L$ that a non-zero element of~$\Ainf$
multiplies into \m{\Ainf\otimes L}.

\emph{Kisin modules.} \m{\Minf(L)} is the $\Ainf$-module of the $x\in\WC\otimes L$ such that
\m{\tau^{a}x\in\Ainf\otimes L} for some \m{a\ge0} and $\mu(\varphi^{a+1}x)\in\Fil^{0}\BdR$ for all
\m{a\ge0} and \m{\mu\in\Per(L)}. \m{\Mlocinf(L)} is the $\Ainf$-module of the~$x$ that some finite product of
the~$\cE_\nu$ multiplies into \m{\Minf(L)}. For an $\Ainf$-submodule~$Q$ of \m{\WC\otimes L}, its \kw{points
over~$F_0$} are the \m{x\in Q} with \m{\mu(x)\in F_0} for all \m{\mu\in\Per(L)}; they form an $\fS$-module.
\m{\mathfrak{M}(L)} is the $\fS$-module of the points of \m{\Minf(L)} over~$F_0$ (the \kw{Kisin module}), and
\m{\Mloc(L)} that of the points of \m{\Mlocinf(L)} over~$F_0$. For \m{l\ge0}, \m{\std_l} is the exponent
function with \m{\std_l(\nu,u)=0} for \m{\nu<l} and \m{\std_l(\nu,u)=-u} for \m{\nu\ge l}.
\end{setting}

\begin{definition}[the cost of a step]\label{d:cost}
In Setting~\ref{set:rings}, for integers \m{\nu\ge0} and~$n$,
\[\cost{\nu}{n}\ =\ v_p(n)+\fd+3\nu .\]
\end{definition}
\leanline{\lean{cost}}

\begin{definition}[elements with exponents]\label{d:exp}
In Setting~\ref{set:lattice}, for a lattice datum~$L$, a letter~$i$ and an exponent function~$o$,
\m{\wExp{L}{i}{o}} is the set of the $x\in\Mlocinf(L)\cap(\WC\otimes L)_i$ such that, for every \m{\nu\ge0},
\[\mu(\varphi^{-\nu}x)\in\Fil^{\,o(\nu,u)}\BdR\qquad\text{for all }u\in U_L(i^{p^{-\nu}})\text{ and all }
\mu\in\Per_{i^{p^{-\nu}},u}(L).\]
It is an $\Ainf$-submodule of \m{\WC\otimes L}. For an integer \m{l\ge0},
\[\locLP{l}{i}(L)\ =\ \Mloc(L)\cap\wExp{L}{i}{\std_l},\]
an $\fS$-module.
\end{definition}
\leanline{\lean{withExponents}\sep \lean{locLetterPart}}

\begin{definition}[parts, invariants, operations on exponents]\label{d:ops}
In Setting~\ref{set:lattice}, let $L$, $A$, $C$ be lattice data.
\begin{enumerate}[label=\textup{(\roman*)},leftmargin=2.3em,itemsep=2pt]
\item For an $\Ainf$-submodule~$Q$ of \m{\WC\otimes L}, \m{\locKP{Q}} is the $\Zp$-module of the elements of
\m{\Mloc(L)} that lie in~$Q$, and \m{\invMod{Q}} is the $\Zp$-module of the $x\in\WC\otimes L$ with
\m{g x-x\in Q} for all \m{g\in\Gamma'}. (The module \m{\Mloc(L)} is a $\Zp$-module because $\fS$ is a
$\Zp$-algebra; the formal text records this as an instance.)
\item For exponent functions $o^{\mathrm{lo}}$, $o^{\mathrm{hi}}$, an integer \m{\nu\ge0} and
\m{u\in\mathbb{Z}},
\[\begin{gathered}\impE{o^{\mathrm{lo}}}{o^{\mathrm{hi}}}{\nu}{u}=o^{\mathrm{hi}}(\nu,u)\ \text{ if }o^{\mathrm{lo}}(\nu,u)>0,\\
\impE{o^{\mathrm{lo}}}{o^{\mathrm{hi}}}{\nu}{u}=\min\bigl(o^{\mathrm{hi}}(\nu,u),0\bigr)\ \text{ otherwise.}\end{gathered}\]
\item For integers \m{l\ge0} and~$n$, an exponent function~$o$, an integer \m{\nu\ge0} and
\m{u\in\mathbb{Z}},
\[\frobX{l}{n}{o}{\nu}{u}=n\ \text{ if }\nu<l,\qquad \frobX{l}{n}{o}{\nu}{u}=o(\nu-l,u)\ \text{ if }\nu\ge l .\]
\item For an $\cO$-bilinear map \m{b\colon A\times L\to C} of the lattices, $y\in\WC\otimes A$ and
$x\in\WC\otimes L$, $\wPair{b}{y}{x}\in\WC\otimes C$ is the value at \m{(y,x)} of the $\WC$-bilinear
extension of~$b$.
\end{enumerate}
\end{definition}
\leanline{\lean{locKisinPart}\sep \lean{invariantsMod}\sep \lean{improvedExponent}\sep \lean{frobExponent}\sep
\lean{wittPairinga}}

\begin{definition}[sides]\label{d:side}
In Setting~\ref{set:lattice}, a \kw{side} over~$\cO$ is a triple \m{S=(Y_S,C_S,\alpha_S)} of a lattice
datum~$Y_S$ over~$\cO$ with \m{0\in U_{Y_S}(i)} for every letter~$i$, a lattice~$C_S$ over~$\cO$, and an
$\cO$-linear map $\alpha_S\colon C_S\to Z^{1}_{\mathrm{f}}(Y_S)$. For a side~$S$:
\begin{enumerate}[label=\textup{(\roman*)},leftmargin=2.3em,itemsep=2pt]
\item \m{\extL{S}} is the lattice datum on \m{Y_S\times C_S} with the action
\m{g(y,c)=(gy+\alpha_S(c)(g),\,c)} and \m{U_{\extL{S}}(i)=U_{Y_S}(i)} for every letter~$i$: the extension of the
module~$C_S$, with trivial action, by~$Y_S$ along the cocycles~$\alpha_S$. It is crystalline on~$\Gamma'$, and its
non-zero jumps are jumps of~$Y_S$ (Input~\ref{i:ext}).
\item \m{\trivL{S}} is the lattice datum formed by the module~$C_S$ with the trivial action of~$\Gamma$ and
\m{U_{\trivL{S}}(i)=\{0\}} for every letter~$i$.
\item $\projC{S}\colon\extL{S}\to\trivL{S}$ is the equivariant map \m{(y,c)\mapsto c}.
\end{enumerate}
\end{definition}
\leanline{\lean{PairingSide}\sep \lean{extLattice}\sep \lean{trivLattice}\sep \lean{projTriv}}

\begin{definition}[windows]\label{d:win}
\begin{enumerate}[label=\textup{(\roman*)},leftmargin=2.3em,itemsep=2pt]
\item For integers \m{\eta,\beta,W\ge0}, a sign \m{\varsigma\in\{1,-1\}} and a finite set~$U$ of integers,
$\JumpWin{\eta}{\varsigma}{\beta}{W}{U}$ means that every \m{u\in U} satisfies \m{-\eta\le u\le\eta} or
\m{\beta\le\varsigma u\le\beta+W}.
\item For integers $n$, \m{\nu\ge0} and~$u$, \m{\raisedE{n}{\nu}{u}=n-u} if \m{\nu=0} and
\m{\raisedE{n}{\nu}{u}=-u} if \m{\nu>0}.
\item For integers \m{n\ge0} and~$c$, $\twoWin{n}{c}=[-n,n]\cup[c-n,c+n]$, a finite set of integers.
\item For integers \m{p,n\ge0} and~$\beta$, \m{\NoMult{p}{n}{\beta}} means that every
$a\in\twoWin{n}{\beta}\cup\twoWin{n}{-\beta}$ that is divisible by~$p$ is~$0$.
\end{enumerate}
\end{definition}
\leanline{\lean{IsInJumpWindow}\sep \lean{raisedAtE}\sep \lean{twoWindows}\sep \lean{NoMultipleNear}}

\begin{definition}[constants, the lifting ideal, cokernel killers]\label{d:lift}
In Setting~\ref{set:lattice}, let $S$ be a side over~$\cO$ and $i$ a letter.
\begin{enumerate}[label=\textup{(\roman*)},leftmargin=2.3em,itemsep=2pt]
\item For \m{c\in C_S}, $\constAt{S}{i}{c}=\pr_i(1\otimes c)\in\WC\otimes C_S$; the map
\m{c\mapsto\constAt{S}{i}{c}} is $\Zp$-linear.
\item For integers \m{r,l\ge0}, \m{\liftId{S}{r}{l}{i}} is the ideal of~$\Ainf$ generated by~$r$ and by the
\m{a\in\Ainf} with
\[a\cdot\pr_i\bigl(\Ainf\otimes C_S\bigr)\ \subseteq\ (1\otimes\projC{S})\bigl(\wExp{\extL{S}}{i}{\std_l}\bigr).\]
\item For integers \m{\kappa,\rho\ge0}, \m{\Kills{S}{i}{\kappa}{\rho}} means
$p^{\kappa}\cE^{\rho}\in\liftId{S}{0}{0}{i}$; in words, \m{p^{\kappa}\cE^{\rho}} \kw{kills the cokernel
of~$S$ at~$i$}.
\end{enumerate}
\end{definition}
\leanline{\lean{constantAt}\sep \lean{liftIdeal}\sep \lean{KillsCokernel}}

\inputbox{\textbf{Inputs from the background library.} The following results are taken for granted. Each is a
declaration of the background library, read in its Lean text; the number in brackets is the unit that owns it.
All are in Setting~\ref{set:lattice}, for a lattice datum~$L$.
\begin{enumerate}[label=\textup{I\arabic*.},ref=I\arabic*,leftmargin=2.6em,itemsep=2.5pt]
\item\label{i:div} \emph{Divisibility and filtration} [83]. For \m{r\in\Ainf} and \m{\nu\ge0}: \m{\cE_\nu\mid r}
if and only if \m{\theta_\nu(r)=0}, if and only if $\varphi^{-\nu}(r)\in\Fil^{1}\BdR$. The image
\m{\varphi^{-\nu}(\cE_\nu)} generates the maximal ideal of~$B^{+}_{\mathrm{dR}}$. If \m{\cE_\nu\nmid r}, then
\m{\varphi^{-\nu}(r)} is a unit of~$B^{+}_{\mathrm{dR}}$. Moreover \m{\cE_\nu\nmid\cE_{\nu'}} for
\m{\nu\ne\nu'}, \m{\cE_\nu\nmid z\tau}, \m{\cE_\nu\nmid p}, \m{\cE_\nu\mid\varphi^{l}(\tau)} for \m{\nu<l},
\m{z\tau\mid\varphi^{\nu}(z\tau)}, \m{\theta(z)\mid p} and \m{\tau\mid\xi-p}; $p$ divides neither $z$ nor~$\tau$,
so \m{z\tau} and the~$\cE_\nu$ are units of~$\WC$.
\item\label{i:gal} \emph{Galois action on the chart} [83]. For \m{g\in\GQp}, \m{g\cdot\cE} is~$\cE$ times a unit
of~$\Ainf$. If $g$ fixes~$\WW$ and \m{d\in\Ainf} divides \m{gz-z}, then $d$ divides \m{gs-s} for every
\m{s\in\fS}; if moreover $g$ fixes \m{\theta(z)}, it fixes \m{\theta(s)} for every \m{s\in\fS}. For
\m{g\in\Gamma'} and \m{s\in\fS}, \m{z\tau\cE} divides \m{gs-s}.
\item\label{i:der} \emph{A derivative} [83]. There is a polynomial~$P_0$ over~$\WW$ with \m{\cE\mid P_0(z)}
in~$\fS$ and \m{\theta(P_0'(z))\mid p^{\fd}}. The handout's catalogue relates it to the different of a totally ramified extension (J.-P.~Serre, \emph{Local Fields}, Ch.~III, §6; place not confirmed).
\item\label{i:lte} \emph{Lifting the exponent} [83]. Let \m{f\colon\Ainf\to R} be a surjective ring
homomorphism, \m{a\in\Ainf}, and \m{v\in\Ainf} a non-unit dividing \m{a-1} with
\m{p\,f(a-1)\,f(v)\mid f(a-1)^{p}}. Then $f(a^{n}-1)\mid p^{v_p(n)}f(a-1)$ for every integer \m{n\ge1}. Also:
for \m{n\ge1} there are \m{\alpha,\beta\in\Ainf} with
$\tau\cE\,(p^{v_p(n)}+\beta z)=\alpha\,([\varepsilon]^{n}-1)$. A divisibility form of the folklore lemma on lifting the
exponent.
\item\label{i:mod} \emph{Residues modulo~$\cE_\nu$} [83]. For \m{\nu\ge0}, every \m{s\in\fS} is congruent
modulo~$\cE_\nu$ to $\sum_{a<p^{\nu}e_{\fc}}w_a z^{a}$ with \m{w_a\in\WW}. Moreover \m{z^{e_{\fc}}} lies in the
ideal \m{(\cE,p)} of~$\fS$, and \m{\theta_\nu([\varepsilon])} is a primitive $p^{\nu}$-th root of unity.
\end{enumerate}}

\inputbox{\textbf{Inputs from the background library, continued.}
\begin{enumerate}[label=\textup{I\arabic*.},ref=I\arabic*,leftmargin=2.6em,itemsep=2.5pt,start=6]
\item\label{i:tate} \emph{Tate} [58]. Let $V\in\BdR$ and let $M$ be a $B^{+}_{\mathrm{dR}}$-submodule of~$\BdR$.
If \m{gV-V\in M} for all \m{g\in\Gamma'}, there is \m{c\in F} with \m{V-c\in M}. An element~$c$
of~$\bar{\mathbb{Q}}_p$ lies in \m{\Fil^{n}\BdR} if and only if \m{c=0} or \m{n\le0}. An element of~$\BdR$
fixed by~$\Gamma'$ lies in \m{\Fil^{0}\BdR\cap F_0}. The first statement is, in the second reader's comparison,
a weaker form of the degree-zero parts of Tate's two theorems (J.~Tate, \emph{$p$-Divisible Groups}, §3.3,
Theorems 1 and~2; P.~Colmez, \emph{Théorie d'Iwasawa des représentations de de Rham d'un corps local}, §III.1).
\item\label{i:val} \emph{Values in the fields~$F_\nu$} [83]. \emph{(a)} Let \m{f\colon\Ainf\to\fS} be
$\fS$-linear, \m{r_k\in\Ainf} and \m{v_k\in F_\nu}, finitely many, and $M$ a $B^{+}_{\mathrm{dR}}$-submodule
of~$\BdR$. If $\sum_k\varphi^{-\nu}(r_k)v_k\in M$, then $\sum_k\varphi^{-\nu}(f(r_k))v_k\in M$.
\emph{(b)} Let \m{V\in F_\nu\cap\Fil^{0}\BdR} and \m{g,t_1,\dots,t_m\in\Gamma'}. If
$\cE_\nu\mid g\sum_k t_k s-\sum_k t_k s$ for all \m{s\in\fS}, then
$g\sum_k t_kV-\sum_k t_kV\in\Fil^{1}\BdR$.
\item\label{i:frame} \emph{The frame theorem} [83]. \emph{(a)} If \m{v_1,\dots,v_m} are points of the fraction
span over~$F_0$ and \m{c_k\in\Ainf} with \m{\sum_kc_kv_k=0}, then \m{\sum_kf(c_k)v_k=0} for every $\fS$-linear
\m{f\colon\Ainf\to\fS}. \emph{(b)} $\Minf(L)=\Ainf\cdot\mathfrak{M}(L)$. \emph{(c)}
$\tau^{h_L}\Minf(L)\subseteq\Ainf\otimes L$ and $\tau^{h_L}(\Ainf\otimes L)\subseteq\Minf(L)$.
\emph{(d)} $\cE^{h_L}\varphi(x)\in\mathfrak{M}(L)$ for \m{x\in\mathfrak{M}(L)}. \emph{(e)} Let
\m{d\in\Ainf} be such that $d$ divides~$r$ whenever $d$ divides \m{\cE^{a}r} in~$B^{+}$. If \m{x\in\Minf(L)}
and \m{\cE^{2h_L}\mu(x)\in B^{+}d} for every crystalline period~$\mu$ (a period in \m{\Per(L)} whose values on
\m{1\otimes L} lie in the image of~$B_{\max}$), then \m{x\in d\,\Minf(L)}; and \m{\cE^{h_L}\mu(x)} lies in
the image of~$\fS_B$ for such~$\mu$ and \m{x\in\mathfrak{M}(L)}. The element \m{d=z\tau} has the divisibility
property, and \m{z\tau\mid gs-s} in~$B^{+}$ for \m{g\in\Gamma'}, \m{s\in\fS_B}. This is the artifact's form of
the existence and the finite height of the Breuil--Kisin module of a crystalline lattice; the file names it a
known result (Kisin lattice of a crystalline lattice). Reference to be confirmed.
\end{enumerate}}

\inputbox{\textbf{Inputs from the background library, concluded.}
\begin{enumerate}[label=\textup{I\arabic*.},ref=I\arabic*,leftmargin=2.6em,itemsep=2.5pt,start=9]
\item\label{i:forms} \emph{Periods} [82, 83]. \emph{(a)} For \m{\mu\in\Per(L)}, $x$ in the fraction span and
\m{g\in\Gamma'}: \m{\mu(gx)=g\cdot\mu(x)}. The fraction span contains \m{\Minf(L)} and is stable under
\m{\mathbb{Z}\times\Gamma}, as is \m{\Ainf\otimes L}; \m{\Minf(L)} is stable under~$\Gamma'$; \m{\WC\otimes L} has
no $\Ainf$-torsion. \emph{(b)} For an equivariant \m{f\colon L\to L'}: \m{\mu\circ(1\otimes f)\in\Per(L)} for
\m{\mu\in\Per(L')}, and \m{1\otimes f} maps \m{\Minf(L)} into \m{\Minf(L')}; \m{\mu\circ\pr_i\in\Per(L)} for
\m{\mu\in\Per(L)}. \emph{(c)} $\varphi^{\nu}x\in(\WC\otimes L)_i$ if and only if
$x\in(\WC\otimes L)_{i^{p^{-\nu}}}$; $\pr_i$ commutes with~$\Gamma'$ and with \m{1\otimes f}. \emph{(d)} If
\m{u\notin U_L(i)}, then $\Per_{i,u}(L)\subseteq\Per_{i,u-1}(L)$. \emph{(e)} Every $\Zp$-linear map
\m{L\to\BdR} is \m{t\mapsto\mu(1\otimes t)} for an $\Ainf$-linear \m{\mu\colon\WC\otimes L\to\BdR}; two such
maps that agree on \m{1\otimes L} agree on the fraction span.
\item\label{i:hodge} \emph{A Hodge basis} [87]. Let \m{(e_k)_k} be a $\Zp$-basis of~$L$. There are integers
$\kappa_k$ and an invertible matrix~$\mathsf{H}$ over~$\BdR$, the product of an invertible matrix
over~$B^{+}_{\mathrm{dR}}$ and of a diagonal matrix with entries generating \m{\Fil^{\kappa_k}\BdR}, such that
$\mathsf{R}(g)\cdot g(\mathsf{H})=\mathsf{H}$ for all \m{g\in\Gamma'}, where \m{\mathsf{R}(g)} is the matrix
of~$g$ on~$L$ in the basis. This is the Hodge--Tate decomposition of a potentially crystalline lattice; in the
second reader's comparison the passage from potential crystallinity to Hodge--Tate weights is classical for
field coefficients (J.-M.~Fontaine, \emph{Représentations $p$-adiques semi-stables}, Astérisque~223, Exposé~III,
Theorem~3.8(i), Remark~3.9(ii), §5.1.7) and the formulation for coefficient fibres is not located.
\item\label{i:ext} \emph{Extensions and twists} [86, 87]. For a side~$S$, the lattice \m{\extL{S}} is
crystalline on~$\Gamma'$, and for \m{u\ne0} the multiplicity of the jump~$u$ of \m{\extL{S}} at any~$\gamma$ is
at most that of~$Y_S$. If $\Gamma$ acts on a lattice through the $c$-th power of~$\chi$, its jumps are all
equal to~$-c$. Both rest on Tate's vanishing of the invariants of the non-trivial twists of~$C$.
\item\label{i:rad} \emph{A radical} [59]. If \m{a\in\Ainf} is not a unit and \m{d\in\Ainf} is not divisible
by~$p$, then $a$ lies in the radical of the ideal \m{(p,d)}. The ring~$\Ainf$ is $p$-adically complete, and
$p$ does not divide~$\tau$.
\end{enumerate}}

\sect{The depth steps}{s:depth}
Throughout this section the setting is Setting~\ref{set:rings}. By Input~\ref{i:div}, the ring homomorphism
\m{\theta_\nu} has kernel \m{\cE_\nu\Ainf}; it commutes with~$\GQp$, so this ideal is stable under~$\GQp$.

\begin{proof}[Proof of Theorem A]
\emph{The elements.} Choose \m{\gamma\in\Gamma'} with $\gamma z=[\varepsilon]^{p^\nu}z$. By Input~\ref{i:gal},
\m{\gamma\cE=w_0\cE} for a unit~$w_0$ of~$\Ainf$, and \m{z\tau\cE} divides
\m{\gamma\cE-\cE=(w_0-1)\cE}; as $\cE\ne0$ and $\Ainf$ is a domain, \m{z\tau\mid w_0-1}. For every integer
\m{n\ne0}, \m{w_0-1} divides \m{w_0^{n}-1} (for \m{n<0} write \m{w_0^{n}-1=-w_0^{n}(w_0^{-n}-1)}). Since
\m{z\tau\mid\varphi^\nu(z\tau)} (Input~\ref{i:div}), \m{z\tau} divides
\m{\varphi^\nu(w_0^{n}-1)=w^{n}-1}, where \m{w=\varphi^\nu(w_0)}. Let \m{g=\gamma} and let~$X$ be the
quotient, so that \m{z\tau X=w^{n}-1}. Then $g\cE_\nu=\varphi^\nu(\gamma\cE)=w\cE_\nu$, because the actions
of~$\varphi$ and of~$\GQp$ commute. Next $\theta_\nu([\varepsilon]^{p^\nu}-1)=\theta([\varepsilon]-1)=0$, so
$\cE_\nu$ divides $\gamma z-z=([\varepsilon]^{p^\nu}-1)z$, and therefore \m{\gamma s-s} for every
\m{s\in\fS}, by Input~\ref{i:gal} ($\gamma$ fixes~$\WW$).

\emph{Reduction.} It remains to find~$y$ with $\cE_\nu\mid yX-p^{\cost{\nu}{n}}$. It suffices that
\m{\theta_\nu(X)} divides \m{p^{\cost{\nu}{n}}} in~$\mathcal{O}_C$: then, $\theta$ being surjective, there is
$y_0$ with $\theta(y_0)\theta_\nu(X)=p^{\cost{\nu}{n}}$, and \m{y=\varphi^\nu(y_0)} has
$\theta_\nu(yX-p^{\cost{\nu}{n}})=0$. Now \m{\theta_\nu(z\tau)\ne0} (Input~\ref{i:div}) and
$\theta_\nu(z\tau)\,\theta_\nu(X)=\theta_\nu(w^{n}-1)=\theta(w_0^{n}-1)$. So it suffices to prove
\begin{equation}\label{e:goal}
\theta(w_0^{n}-1)\ \mid\ p^{\,v_p(n)+\fd+3\nu}\cdot\theta_\nu(z\tau)\qquad\text{in }\mathcal{O}_C .
\end{equation}

\emph{The term \m{v_p(n)}.} Write \m{w_0-1=z\tau d} with \m{d\in\Ainf}. The number \m{\theta(\tau)+1} is a
$p$-th root of unity, so the binomial theorem gives $\theta(\tau)^{p}=-p\,\theta(\tau)\,c$ for some
\m{c\in\mathcal{O}_C}. Hence
\[\theta(w_0-1)^{p}=-p\,c\,\theta(\tau)\,\theta(z)^{p}\theta(d)^{p}\qquad\text{is divisible by}\qquad
p\cdot\theta(w_0-1)\cdot\theta(z)=p\,\theta(\tau)\,\theta(z)^{2}\theta(d),\]
because \m{p\ge2}. So Input~\ref{i:lte} applies to \m{f=\theta}, \m{a=w_0} and the non-unit \m{v=z}, with
the exponent~$|n|$: \m{\theta(w_0^{|n|}-1)} divides \m{p^{v_p(n)}\theta(w_0-1)}, and \m{w_0^{n}-1} is
\m{w_0^{|n|}-1} times a unit.

\emph{The term \m{\fd+\nu}.} Let $P_0$ be the polynomial of Input~\ref{i:der}, with \m{P_0(z)=\cE G},
\m{G\in\fS}. Taylor's formula gives \m{r\in\Ainf} with
\[P_0(\gamma z)=P_0(z)+P_0'(z)\,(\gamma z-z)+r\,(\gamma z-z)^{2}.\]
The left side is $\gamma(P_0(z))=w_0\cE\,\gamma(G)$, as $\gamma$ fixes~$\WW$. On the right,
$\gamma z-z=([\varepsilon]^{p^\nu}-1)z$ and
$[\varepsilon]^{p^\nu}-1=([\varepsilon]-1)\,\sigma_\nu$ with
$\sigma_\nu=\sum_{a<p^\nu}[\varepsilon]^{a}$, while $[\varepsilon]-1=\xi\tau=\cE\tau\,\mathsf{u}$ for a
unit~$\mathsf{u}$. Dividing by~$\cE$:
\[w_0\,\gamma(G)-G=P_0'(z)\,\tau\mathsf{u}\,\sigma_\nu z+r\,\tau\mathsf{u}\,\sigma_\nu z\cdot
([\varepsilon]^{p^\nu}-1)z .\]
Apply~$\theta$. The last term vanishes; \m{\theta(\gamma G)=\theta(G)} by Input~\ref{i:gal}; and
\m{\theta(\sigma_\nu)=p^{\nu}}, since \m{\theta([\varepsilon])=1}. Therefore
$\theta(G)\,\theta(w_0-1)=\theta(P_0'(z))\,\theta(\mathsf{u})\,p^{\nu}\,\theta(z\tau)$, and
\m{\theta(w_0-1)} divides \m{p^{\fd+\nu}\,\theta(z\tau)}.

\emph{The term \m{2\nu}.} If \m{\nu=0} then \m{\theta_\nu=\theta}. If \m{\nu\ge1}, then \m{\theta(z)\mid p}
and \m{\theta(\tau)\mid p} (Input~\ref{i:div}: $\tau$ divides \m{\xi-p} and \m{\theta(\xi)=0}), so
\m{\theta(z\tau)} divides~$p^{2}$, hence \m{p^{2\nu}\theta_\nu(z\tau)}. In both cases
$\theta(z\tau)\mid p^{2\nu}\theta_\nu(z\tau)$. The three divisibilities give~\eqref{e:goal}.
\end{proof}

\begin{lemma}[the depth step at exponent~$0$]\label{l:zero}
In Setting~\ref{set:rings}, let \m{\nu\ge0} be an integer and put \m{m=p^{\cost{\nu}{0}}}. There are
\m{t_1,\dots,t_m\in\Gamma'} such that, for every \m{g\in\Gamma'} and every \m{s\in\fS},
\[\cE_\nu\ \mid\ g\cdot\sum_{k=1}^{m}t_k\cdot s-\sum_{k=1}^{m}t_k\cdot s\qquad\text{in }\Ainf .\]
\end{lemma}
\leanline{\lean{existsdepthstepzero}}
\begin{proof}
For \m{0\le k<m} choose \m{t_k\in\Gamma'} with \m{t_kz=[\varepsilon]^{k}z}. Divisibility by~$\cE_\nu$ is the
vanishing of~$\theta_\nu$, which commutes with~$\Gamma'$; so by Input~\ref{i:mod} it suffices to treat
\m{s=wz^{a}} with \m{w\in\WW}, and to show that
$\theta_\nu(g\sum_kt_ks)=\theta_\nu(\sum_kt_ks)$. Each~$t_k$ fixes~$w$, so
$\theta_\nu(t_ks)=\theta_\nu(w)\theta_\nu(z)^{a}\varrho^{ka}$, where $\varrho=\theta_\nu([\varepsilon])$ is a
primitive $p^\nu$-th root of unity. If \m{p^\nu\nmid a}, then \m{\varrho^{a}\ne1} and
\m{(\varrho^{a})^{m}=1}, because \m{p^\nu} divides \m{m=p^{\fd+3\nu}}; the geometric sum
\m{\sum_{k<m}\varrho^{ka}} vanishes, and both sides are~$0$. If \m{a=p^\nu a'}, then
\m{z^{a}=\varphi^{\nu}(z^{a'})} and, for every \m{\gamma\in\Gamma'},
$\theta_\nu(\gamma s)=\theta(\gamma\varphi^{-\nu}(w))\,\theta(\gamma z^{a'})=\theta_\nu(s)$, because
$\gamma$ fixes $\WW$ and~\m{\theta(z)}. So $\theta_\nu(\sum_kt_ks)=m\,\theta_\nu(s)$, and the same
computation with \m{gt_k} in place of~$t_k$ gives the same value.
\end{proof}

\sect{Exponent modules: first properties}{s:first}
Throughout, the setting is Setting~\ref{set:lattice}; $L$ and~$L'$ are lattice data over~$\cO$.

\begin{lemma}[the localized modules]\label{l:loc}
Let $x\in\WC\otimes L$, let \m{\nu,l\ge0} and $n$ be integers, let \m{g\in\Gamma'} and let
\m{f\colon L\to L'} be an equivariant map.
\begin{enumerate}[label=\textup{(\alph*)},leftmargin=2.3em,itemsep=2pt]
\item \m{\cE_\nu x\in\Mlocinf(L)} if and only if \m{x\in\Mlocinf(L)}.
\item \m{\cE_\nu^{\,n}x\in\Mloc(L)} if and only if \m{x\in\Mloc(L)}.
\item If \m{x\in\Mlocinf(L)}, then $(1\otimes f)(x)\in\Mlocinf(L')$; if \m{x\in\Mloc(L)}, then
$(1\otimes f)(x)\in\Mloc(L')$.
\item If \m{x\in\Mlocinf(L)}, then \m{g x\in\Mlocinf(L)}.
\item \m{\Mlocinf(L)} is contained in the fraction span of~$L$.
\item If \m{x\in\Mlocinf(L)} and \m{\mu\in\Per(L)}, then
$\mu(\varphi^{-\nu}(gx-x))=g\cdot\mu(\varphi^{-\nu}x)-\mu(\varphi^{-\nu}x)$.
\item If \m{x\in\Mloc(L)}, then \m{\varphi^{l}x\in\Mloc(L)}.
\item The $\Ainf$-span of \m{\Mloc(L)} is \m{\Mlocinf(L)}.
\item If \m{x\in\Mloc(L)} and \m{\mu\in\Per(L)}, then \m{\mu(\varphi^{-\nu}x)\in F_\nu}.
\end{enumerate}
\end{lemma}
\leanline{\lean{frobEsmulmemkisinLocAinfiff}\sep \lean{frobEUnitzpowsmulmemkisinLoc}\sep
\lean{wittMapmemkisinLocAinf}\sep \lean{wittMapmemkisinLoc}\sep \lean{galsmulmemkisinLocAinf}\sep
\lean{kisinLocAinflefracSpan}\sep \lean{valueAtsmulsub}\sep \lean{phipowmemkisinLoc}\sep \lean{spankisinLoc}\sep
\lean{valueAtmemvalueField}\sep \lean{valueAtzero} (the identity \m{\mu(\varphi^{0}x)=\mu(x)})}
\begin{proof}
Several proofs go by \emph{induction on the localization}: a property that holds on \m{\Minf(L)} and passes
from \m{\cE_\nu x} to~$x$, for every~$\nu$, holds on \m{\Mlocinf(L)}.

(a) holds by the definition of \m{\Mlocinf(L)} as a saturation. (b) The units~$c$ of~$\WC$ with
\m{c\,\Mloc(L)=\Mloc(L)} form a group, so it suffices to show that \m{\cE_\nu x\in\Mloc(L)} if and only if
\m{x\in\Mloc(L)}. For \m{\Mlocinf} this is~(a); and \m{\mu(\cE_\nu x)=\cE_\nu\mu(x)} lies in~$F_0$ exactly
when \m{\mu(x)} does, because the image of~$\cE_\nu$ in~$\BdR$ is a non-zero element of~$F_0$.

(c) \m{1\otimes f} is $\WC$-linear and maps \m{\Minf(L)} into \m{\Minf(L')}, hence \m{\Mlocinf(L)} into
\m{\Mlocinf(L')}; and \m{\mu\circ(1\otimes f)\in\Per(L)} for \m{\mu\in\Per(L')} (Input~\ref{i:forms}), which
gives the values in~$F_0$. (d) By induction on the localization: \m{\Minf(L)} is stable under~$\Gamma'$, and
\m{g(\cE_\nu x)=g(\cE_\nu)\,gx} with \m{g(\cE_\nu)} equal to~$\cE_\nu$ times a unit of~$\Ainf$, so (a) applies.
(e) By the same induction from \m{\Minf(L)}, which lies in the fraction span. (f) $g$ commutes
with~$\varphi^{-\nu}$, the fraction span is stable under \m{\mathbb{Z}\times\Gamma}, and
\m{\mu(gx')=g\mu(x')} on it (Input~\ref{i:forms}).

(g) It suffices to treat \m{l=1}. By induction on the localization, applied to the points over~$F_0$: for
\m{x\in\mathfrak{M}(L)}, $\cE^{h_L}\varphi(x)\in\mathfrak{M}(L)$ by the frame theorem
(Input~\ref{i:frame}(d)), so \m{\varphi(x)\in\Mloc(L)} by~(b); and
$\varphi(\cE_\nu x)=\cE_{\nu+1}\varphi(x)$, so (b) passes the property from \m{\cE_\nu x} to~$x$.

(h) By induction on the localization: $\Minf(L)=\Ainf\cdot\mathfrak{M}(L)$ (Input~\ref{i:frame}(b)), and if
\m{\cE_\nu x=\sum_ka_kv_k} with \m{v_k\in\Mloc(L)}, then \m{x=\sum_ka_k\,\cE_\nu^{-1}v_k} with
\m{\cE_\nu^{-1}v_k\in\Mloc(L)} by~(b).

(i) By (e) there is \m{r\ne0} in~$\Ainf$ with $r\varphi^{-\nu}x\in\Ainf\otimes L$; then
$\tau^{h_L}r\varphi^{-\nu}x\in\Minf(L)$ (Input~\ref{i:frame}(c)), and by~(h) it is \m{\sum_ka_kv_k} with
\m{a_k\in\Ainf}, \m{v_k\in\Mloc(L)}. Apply~$\varphi^{\nu}$:
\[-\varphi^{\nu}(\tau^{h_L}r)\,x+\sum_k\varphi^{\nu}(a_k)\,\varphi^{\nu}v_k=0 ,\]
a relation with coefficients in~$\Ainf$ between $x$ and the \m{\varphi^{\nu}v_k}, which lie in \m{\Mloc(L)}
by~(g). By Input~\ref{i:frame}(a) it persists when a coordinate~\m{\operatorname{co}_a} is applied to the
coefficients; choose~$a$ with $c=\operatorname{co}_a(\varphi^{\nu}(\tau^{h_L}r))\ne0$. Applying
\m{\mu\circ\varphi^{-\nu}} gives
$\varphi^{-\nu}(c)\,\mu(\varphi^{-\nu}x)=\sum_k\varphi^{-\nu}(\operatorname{co}_a(\varphi^\nu a_k))\,\mu(v_k)$.
Here \m{\mu(v_k)\in F_0\subseteq F_\nu} and $\varphi^{-\nu}(\fS)\subseteq F_\nu$, and $F_\nu$ is a field.
\end{proof}

\begin{lemma}[the exponent modules]\label{l:exp}
Let $x\in\WC\otimes L$, let $i$ be a letter and let $o$, $o'$ be exponent functions.
\begin{enumerate}[label=\textup{(\alph*)},leftmargin=2.3em,itemsep=2pt]
\item Let \m{x\in\wExp{L}{i}{o}}. Suppose that, for every \m{\nu\ge0}, every \m{u\in U_L(i^{p^{-\nu}})} with
\m{o(\nu,u)<o'(\nu,u)} and every \m{\mu\in\Per_{i^{p^{-\nu}},u}(L)}, one has
$\mu(\varphi^{-\nu}x)\in\Fil^{o'(\nu,u)}\BdR$. Then \m{x\in\wExp{L}{i}{o'}}.
\item Let \m{x\in\wExp{L}{i}{o}}. If \m{o'(\nu,u)\le o(\nu,u)} for all \m{\nu\ge0} and
\m{u\in U_L(i^{p^{-\nu}})}, then \m{x\in\wExp{L}{i}{o'}}.
\item Let \m{\nu\ge0}, \m{u\in\mathbb{Z}} and \m{\mu\in\Per_{i^{p^{-\nu}},u}(L)}. If \m{x\in\Ainf\otimes L} and
\m{x\in(\WC\otimes L)_i}, then $\mu(\varphi^{-\nu}x)\in\Fil^{-u}\BdR$.
\item Let \m{f\colon L\to L'} be an equivariant map with \m{U_{L'}(j)\subseteq U_L(j)} for every letter~$j$. If
\m{x\in\wExp{L}{i}{o}}, then $(1\otimes f)(x)\in\wExp{L'}{i}{o}$.
\item Let \m{r\in\Ainf} with \m{\cE_\nu\nmid r} for all \m{\nu\ge0}, and let \m{x\in\Mlocinf(L)}. If
\m{rx\in\wExp{L}{i}{o}}, then \m{x\in\wExp{L}{i}{o}}.
\item Let \m{\nu_0\ge0} and $n$ be integers and \m{x\in\wExp{L}{i}{o}}. If
\m{o'(\nu,u)\le o(\nu,u)+n} for \m{\nu=\nu_0} and \m{o'(\nu,u)\le o(\nu,u)} for \m{\nu\ne\nu_0}, for all
\m{u\in U_L(i^{p^{-\nu}})}, then $\cE_{\nu_0}^{\,n}x\in\wExp{L}{i}{o'}$.
\item If \m{x\in\wExp{L}{i}{o}} and \m{g\in\Gamma'}, then \m{gx-x\in\wExp{L}{i}{o}}.
\end{enumerate}
\end{lemma}
\leanline{\lean{memwithExponentsofmemhodgeFil}\sep \lean{memwithExponentsofle}\sep
\lean{valueAtmemhodgeFilofmemainfLattice}\sep \lean{mapmemwithExponents}\sep \lean{memwithExponentsofsmul}\sep
\lean{frobEUnitzpowmemwithExponents}\sep \lean{smulsubmemwithExponents}}
\begin{proof}
(a) and (b) hold because the filtration decreases. (c) The element \m{\varphi^{-\nu}x} lies in
\m{\Ainf\otimes L} and in the part at~\m{i^{p^{-\nu}}} (Input~\ref{i:forms}), so it is an $\Ainf$-combination
of elements \m{\pr_{i^{p^{-\nu}}}(1\otimes t)}, on which $\mu$ takes values in the
$B^{+}_{\mathrm{dR}}$-module \m{\Fil^{-u}\BdR}. (d) Membership in \m{\Mlocinf} is Lemma~\ref{l:loc}(c);
\m{1\otimes f} commutes with \m{1\otimes\omega} and with~$\varphi$; and for \m{\mu\in\Per_{j,u}(L')} with
\m{u\in U_{L'}(j)\subseteq U_L(j)} the period \m{\mu\circ(1\otimes f)} lies in \m{\Per_{j,u}(L)}, because
\m{1\otimes f} commutes with~$\pr_j$. (e) $r$ is a non-zero-divisor on \m{\WC\otimes L} and commutes with
\m{1\otimes\omega}, so $x$ lies in the part at~$i$; and
$\mu(\varphi^{-\nu}(rx))=\varphi^{-\nu}(r)\mu(\varphi^{-\nu}x)$ with \m{\varphi^{-\nu}(r)} a unit
of~$B^{+}_{\mathrm{dR}}$ (Input~\ref{i:div}). (f) \m{\cE_{\nu_0}^{\,n}x} lies in \m{\Mlocinf(L)} by
Lemma~\ref{l:loc}(a) and in the part at~$i$, and
$\mu(\varphi^{-\nu}(\cE_{\nu_0}^{\,n}x))=\varphi^{-\nu}(\cE_{\nu_0})^{n}\mu(\varphi^{-\nu}x)$; the factor
\m{\varphi^{-\nu}(\cE_{\nu_0})} is a uniformizer of~$B^{+}_{\mathrm{dR}}$ for \m{\nu=\nu_0} and a unit
otherwise (Input~\ref{i:div}). (g) By Lemma~\ref{l:loc}(d) and~(f): \m{gx-x} lies in \m{\Mlocinf(L)}, in the
part at~$i$, and its values are $g\mu(\varphi^{-\nu}x)-\mu(\varphi^{-\nu}x)$, in the Galois-stable module
\m{\Fil^{o(\nu,u)}\BdR}.
\end{proof}

\begin{lemma}[division by \m{z\tau}]\label{l:ztau}
Let \m{g\in\Gamma'} and \m{x\in\Mloc(L)}.
\begin{enumerate}[label=\textup{(\alph*)},leftmargin=2.3em,itemsep=2pt]
\item $(z\tau)^{-1}(gx-x)\in\Mlocinf(L)$.
\item Let $i$ be a letter and $o$ an exponent function. If \m{gx-x\in\wExp{L}{i}{o}}, then
$(z\tau)^{-1}(gx-x)\in\wExp{L}{i}{o}$.
\end{enumerate}
\end{lemma}
\leanline{\lean{smulsubmemzTaukisinLocAinf}\sep \lean{smulsubmemzTauUnitsmulofdeep}}
\begin{proof}
(a) By induction on the localization, for the points over~$F_0$. Let \m{x\in\mathfrak{M}(L)}. Then
\m{gx-x\in\Minf(L)}, and for a crystalline period~$\mu$ we have \m{\mu(gx-x)=g\mu(x)-\mu(x)} and
\m{\cE^{h_L}\mu(x)=s} in the image of~$\fS_B$. Since \m{z\tau} divides \m{gs-s} and \m{g\cE-\cE} in~$B^{+}$,
the element \m{\cE^{2h_L}(g\mu(x)-\mu(x))} is a multiple of~\m{z\tau}; so \m{gx-x\in z\tau\,\Minf(L)} by
Input~\ref{i:frame}(e). For the step, let \m{\cE_\nu x} satisfy the claim, with \m{x\in\Mloc(L)}. Then
\[g(\cE_\nu x)-\cE_\nu x=(g\cE_\nu-\cE_\nu)\,gx+\cE_\nu(gx-x),\]
and \m{z\tau} divides \m{g\cE_\nu-\cE_\nu} in~$\Ainf$ (Input~\ref{i:gal}) while \m{gx\in\Mlocinf(L)}; so
$(z\tau)^{-1}\cE_\nu(gx-x)\in\Mlocinf(L)$, and Lemma~\ref{l:loc}(a) removes~$\cE_\nu$. (b) follows from (a)
and Lemma~\ref{l:exp}(e) with \m{r=z\tau} (Input~\ref{i:div}).
\end{proof}

\begin{lemma}[values at a period of any type]\label{l:any}
Let \m{\nu\ge0}, let $i$ be a letter, $o$ an exponent function, \m{u\in\mathbb{Z}}, \m{x\in\wExp{L}{i}{o}} and
\m{\mu\in\Per_{i^{p^{-\nu}},u}(L)}. Let $M$ be a $B^{+}_{\mathrm{dR}}$-submodule of~$\BdR$ with
\m{\Fil^{o(\nu,v)}\BdR\subseteq M} for every \m{v\in U_L(i^{p^{-\nu}})} with \m{v\le u}. Then
\m{\mu(\varphi^{-\nu}x)\in M}.
\end{lemma}
\leanline{\lean{valueAtmemofhodgeFille}\sep the integer~$u$ need not lie in \m{U_L(i^{p^{-\nu}})}}
\begin{proof}
Put \m{j=i^{p^{-\nu}}} and suppose \m{\mu(\varphi^{-\nu}x)\notin M}. We show by descending induction that
\m{\mu\in\Per_{j,v}(L)} for every \m{v\le u}. If this holds for~$v$ and \m{v\in U_L(j)}, then
$\mu(\varphi^{-\nu}x)\in\Fil^{o(\nu,v)}\BdR\subseteq M$, which is excluded; so \m{v\notin U_L(j)}, and
\m{\mu\in\Per_{j,v-1}(L)} by Input~\ref{i:forms}(d). Now choose \m{r\ne0} in~$\Ainf$ with
\m{rx\in\Ainf\otimes L} (Lemma~\ref{l:loc}(e)). By Lemma~\ref{l:exp}(c),
$\varphi^{-\nu}(r)\,\mu(\varphi^{-\nu}x)\in\Fil^{-v}\BdR$ for every \m{v\le u}, so it is~$0$; as \m{\varphi^{-\nu}(r)\ne0},
\m{\mu(\varphi^{-\nu}x)=0\in M}: a contradiction.
\end{proof}

\sect{The retraction}{s:ret}
\begin{lemma}[a retraction onto the module over~$\fS$]\label{l:ret}
In Setting~\ref{set:lattice}, for a lattice datum~$L$ there is an $\fS$-linear map
\m{r\colon\Mlocinf(L)\to\Mloc(L)} with $r(a x)=\operatorname{co}_0(a)\,x$ for all \m{a\in\Ainf} and
\m{x\in\Mloc(L)}.
\end{lemma}
\leanline{\lean{existscoordRetract}}
\begin{proof}
By Lemma~\ref{l:loc}(h), \m{\Mlocinf(L)} is the image of \m{\Ainf\otimes_{\fS}\Mloc(L)}. The $\fS$-linear map
$a\otimes x\mapsto\operatorname{co}_0(a)x$ on the tensor product factors through this image if and only if
every relation \m{\sum_ka_kx_k=0} with \m{x_k\in\Mloc(L)} implies $\sum_k\operatorname{co}_0(a_k)x_k=0$; this
is Input~\ref{i:frame}(a), the elements of \m{\Mloc(L)} being points of the fraction span over~$F_0$.
\end{proof}

\begin{definition}[the retraction]\label{d:ret}
In Setting~\ref{set:lattice}, for a lattice datum~$L$, $\ret_L\colon\Mlocinf(L)\to\Mloc(L)$ is a chosen map
with the property of Lemma~\ref{l:ret}.
\end{definition}
\leanline{\lean{coordRetract}}

\begin{lemma}[properties of the retraction]\label{l:retprop}
In Setting~\ref{set:lattice}, let $L$, $L'$ be lattice data.
\begin{enumerate}[label=\textup{(\alph*)},leftmargin=2.3em,itemsep=2pt]
\item If \m{v_1,\dots,v_m\in\Mloc(L)}, \m{a_1,\dots,a_m\in\Ainf} and $y\in\Mlocinf(L)$ with
\m{\sum_ka_kv_k=y}, then $\ret_L(y)=\sum_k\operatorname{co}_0(a_k)v_k$. If \m{x\in\Mloc(L)}, then
\m{\ret_L(x)=x}.
\item Let \m{f\colon L\to L'} be an equivariant map, $y\in\Mlocinf(L)$ and $y'\in\Mlocinf(L')$ with
$(1\otimes f)(y)=y'$. Then $(1\otimes f)(\ret_L(y))=\ret_{L'}(y')$.
\item Let \m{n\in\mathbb{Z}}, \m{l\ge0}, \m{\mu\in\Per(L)}, $y\in\Mlocinf(L)$ and \m{v\in F_l}. If
$\mu(\varphi^{-l}y)-v\in\Fil^{n}\BdR$, then $\mu(\varphi^{-l}\ret_L(y))-v\in\Fil^{n}\BdR$.
\item Let $i$ be a letter, $o$ an exponent function and $y\in\Mlocinf(L)$. If \m{y\in\wExp{L}{i}{o}}, then
\m{\ret_L(y)\in\wExp{L}{i}{o}}.
\end{enumerate}
\end{lemma}
\leanline{\lean{coordRetractsum}\sep \lean{coordRetractofmemkisinLoc}\sep \lean{mapcoordRetract}\sep
\lean{valueAtcoordRetractsubmemhodgeFil}\sep \lean{coordRetractmemwithExponents}}
\begin{proof}
(a) is the defining property, with \m{\operatorname{co}_0(1)=1}. (b) Write \m{y=\sum_ka_kv_k} with
\m{v_k\in\Mloc(L)} (Lemma~\ref{l:loc}(h)). Then \m{y'=\sum_ka_k(1\otimes f)(v_k)} with
\m{(1\otimes f)(v_k)\in\Mloc(L')}, and both sides equal $\sum_k\operatorname{co}_0(a_k)(1\otimes f)(v_k)$
by~(a). (c) With $y$ written as in~(b),
\[\mu(\varphi^{-l}y)-v=\sum_k\varphi^{-l}(a_k)\,\mu(\varphi^{-l}v_k)+\varphi^{-l}(1)\cdot(-v),\]
and the values \m{\mu(\varphi^{-l}v_k)} lie in~$F_l$ by Lemma~\ref{l:loc}(i). Input~\ref{i:val}(a) with
\m{f=\operatorname{co}_0} replaces each coefficient~$a_k$ by \m{\operatorname{co}_0(a_k)} and keeps~$1$.
(d) $\ret_L(y)\in\Mloc(L)\subseteq\Mlocinf(L)$. Multiplication by~$\omega$ is an equivariant map
\m{L\to L}, so by~(b) and the $\fS$-linearity of~$\ret_L$,
$(1\otimes\omega)\ret_L(y)=\ret_L((1\otimes\omega)y)=\ret_L([i]y)=[i]\ret_L(y)$, as \m{[i]\in\WW}. The
values are in the right steps of the filtration by~(c) with \m{v=0}.
\end{proof}

\sect{Deepening}{s:deep}
\begin{lemma}[small facts]\label{l:small}
In Setting~\ref{set:lattice}:
\begin{enumerate}[label=\textup{(\alph*)},leftmargin=2.3em,itemsep=2pt]
\item If the Kummer chart is unramified and $n$ is an integer such that \m{p\mid n} implies \m{n=0}, then
\m{\cost{\nu}{n}\le3\nu} for every \m{\nu\ge0}.
\item Every element~$c$ of~$F$, seen in~$\BdR$, satisfies \m{g\cdot c=c} for all \m{g\in\Gamma'}.
\end{enumerate}
\end{lemma}
\leanline{\lean{costleofisUnramified}\sep \lean{galMonoidsmullocalFieldOf}}
\begin{proof}
(a) \m{v_p(n)=0} and \m{\fd=0}. (b) $F$ is the field fixed by \m{j_\Gamma(\Gamma')}.
\end{proof}

\begin{lemma}[deep elements have better exponents]\label{l:deep}
In Setting~\ref{set:lattice}, let $L$ be a lattice datum, $i$ a letter, $o^{\mathrm{lo}}$ and
$o^{\mathrm{hi}}$ exponent functions, and $x\in\wExp{L}{i}{o^{\mathrm{lo}}}$. If
$x\in\invMod{\wExp{L}{i}{o^{\mathrm{hi}}}}$, then
$x\in\wExp{L}{i}{\impE{o^{\mathrm{lo}}}{o^{\mathrm{hi}}}{\cdot}{\cdot}}$.
\end{lemma}
\leanline{\lean{memwithExponentsofdeep}}
\begin{proof}
Fix \m{\nu\ge0}, \m{u\in U_L(i^{p^{-\nu}})} and \m{\mu\in\Per_{i^{p^{-\nu}},u}(L)}, and write
\m{a=o^{\mathrm{lo}}(\nu,u)}, \m{b=o^{\mathrm{hi}}(\nu,u)} and \m{V=\mu(\varphi^{-\nu}x)}. Then
\m{V\in\Fil^{a}\BdR}, and $gV-V=\mu(\varphi^{-\nu}(gx-x))\in\Fil^{b}\BdR$ for all \m{g\in\Gamma'}
(Lemma~\ref{l:loc}(f)). By Input~\ref{i:tate} there is \m{c\in F} with \m{V-c\in\Fil^{b}\BdR}. Then
\m{c\in\Fil^{\min(a,b)}\BdR}, so \m{c=0} or \m{\min(a,b)\le0}. If \m{a>0}: either \m{c=0}, or \m{b\le0} and
$c\in\Fil^{0}\BdR\subseteq\Fil^{b}\BdR$; in both cases \m{V\in\Fil^{b}\BdR}. If \m{a\le0}: \m{V-c} and~$c$
lie in \m{\Fil^{\min(b,0)}\BdR}.
\end{proof}

\begin{proposition}[one step deeper, at an explicit cost]\label{p:succ}
In Setting~\ref{set:lattice}, let $L$ be a lattice datum, \m{\nu\ge0} an integer and $n$ an integer. There are
an integer \m{m\ge0}, elements \m{t_1,\dots,t_m\in\Gamma'} and \m{y\in\Ainf} such that for every
\m{x\in\Mloc(L)} and every \m{\mu\in\Per(L)}: if there is \m{c\in F} with
$\mu(\varphi^{-\nu}x)-c\in\Fil^{n}\BdR$, then there is \m{c'\in F} with
\[\mu\Bigl(\varphi^{-\nu}\Bigl(p^{\cost{\nu}{n}}x+\sum_{k=1}^{m}y\,(z\tau)^{-1}(t_kx-x)\Bigr)\Bigr)-c'
\ \in\ \Fil^{n+1}\BdR .\]
\end{proposition}
\leanline{\lean{existsreachesdepthsucc}}
\begin{proof}
\emph{The case \m{n\ne0}.} Take \m{g,w,y_0,X} from Theorem~A for \m{(\nu,n)}, and \m{m=1}, \m{t_1=g},
\m{y=-y_0}. Let \m{V=\mu(\varphi^{-\nu}x)} and \m{V-c\in\Fil^{n}\BdR}. Put
$\pi_\nu=\varphi^{-\nu}(\cE_\nu)$, a uniformizer of~$B^{+}_{\mathrm{dR}}$ in~$F_\nu$, and
\m{V-c=\pi_\nu^{\,n}V_0}; then \m{V_0\in F_\nu\cap\Fil^{0}\BdR}, by Lemma~\ref{l:loc}(i). With
\m{\hat w=\varphi^{-\nu}(w)} we have \m{g\pi_\nu=\hat w\pi_\nu} and \m{gc=c}, so
\[gV-V=\pi_\nu^{\,n}\bigl(\hat w^{n}\,gV_0-V_0\bigr)
=\pi_\nu^{\,n}\bigl((\hat w^{n}-1)V_0+\hat w^{n}(gV_0-V_0)\bigr).\]
Here \m{gV_0-V_0\in\Fil^{1}\BdR} by Input~\ref{i:val}(b), applied to the single element~$1$ and to the
divisibility \m{\cE_\nu\mid gs-s} of Theorem~A; and $\hat w^{n}-1=\varphi^{-\nu}(z\tau)\,\varphi^{-\nu}(X)$.
As \m{\varphi^{-\nu}(z\tau)} is a unit of~$B^{+}_{\mathrm{dR}}$, Lemma~\ref{l:loc}(f) gives
\[\mu\bigl(\varphi^{-\nu}((z\tau)^{-1}(gx-x))\bigr)\ \equiv\ \pi_\nu^{\,n}\varphi^{-\nu}(X)V_0
\pmod{\Fil^{n+1}\BdR}.\]
Multiply by \m{\varphi^{-\nu}(y_0)} and use
$\varphi^{-\nu}(y_0X)\equiv p^{\cost{\nu}{n}}$ modulo \m{\Fil^{1}\BdR}: the value of~$\mu$ at
$\varphi^{-\nu}(y_0(z\tau)^{-1}(gx-x))$ is congruent to \m{p^{\cost{\nu}{n}}(V-c)}. So the element of the
statement has value congruent to $p^{\cost{\nu}{n}}V-p^{\cost{\nu}{n}}(V-c)=p^{\cost{\nu}{n}}c$, and
\m{c'=p^{\cost{\nu}{n}}c} serves.

\emph{The case \m{n=0}.} Take the \m{m=p^{\cost{\nu}{0}}} elements~$t_k$ of Lemma~\ref{l:zero} and
\m{y=z\tau}. The element of the statement is \m{mx+\sum_k(t_kx-x)=\sum_kt_kx}, with value
\m{\sum_kt_kV}, where $V=\mu(\varphi^{-\nu}x)\in F_\nu\cap\Fil^{0}\BdR$. By Lemma~\ref{l:zero} and
Input~\ref{i:val}(b), $g\sum_kt_kV-\sum_kt_kV\in\Fil^{1}\BdR$ for all \m{g\in\Gamma'}; Input~\ref{i:tate}
gives~$c'$.
\end{proof}

\sect{The Hodge expansion and the proof of Theorem B}{s:hodge}
Throughout, the setting is Setting~\ref{set:lattice}. The space \m{\operatorname{Fun}(L)} of all $\Ainf$-linear
maps \m{\WC\otimes L\to\BdR} is a vector space over~$\BdR$.

\begin{lemma}[testing a value on a Hodge basis]\label{l:basis}
Let $L$ be a lattice datum, $x$ an element of the fraction span of~$L$, $i$ a letter, \m{n\in\mathbb{Z}} and
\m{\mu\in\operatorname{Fun}(L)}. Suppose that $D(\mu)\cdot\mu'(x)\in\Fil^{n}\BdR$ for every
\m{\kappa\in\mathbb{Z}}, every \m{\mu'\in\Per_{i,\kappa}(L)} and every $\BdR$-linear map
$D\colon\operatorname{Fun}(L)\to\BdR$ with the two properties
\begin{enumerate}[label=\textup{(D\arabic*)},leftmargin=2.9em,itemsep=2pt]
\item if $\lambda\in\operatorname{Fun}(L)$, \m{a\in\mathbb{Z}} and $\lambda(1\otimes t)\in\Fil^{a}\BdR$ for
all \m{t\in L}, then $D(\lambda)\in\Fil^{a+\kappa}\BdR$;
\item if $\lambda,\lambda'\in\operatorname{Fun}(L)$, \m{g\in\Gamma'} and
$\lambda'(1\otimes gt)=g\cdot\lambda(1\otimes t)$ for all \m{t\in L}, then \m{D(\lambda')=g\cdot D(\lambda)}.
\end{enumerate}
Then \m{\mu(x)\in\Fil^{n}\BdR}.
\end{lemma}
\leanline{\lean{memhodgeFilofhodgebasis}}
\begin{proof}
Take a $\Zp$-basis \m{(e_k)_k} of~$L$ and $\kappa_k$, $\mathsf{H}$ as in Input~\ref{i:hodge}. By
Input~\ref{i:forms}(e) there are \m{\mu_k\in\operatorname{Fun}(L)} with
$\mu_k(1\otimes e_a)=(\mathsf{H}^{-1})_{ka}$. For $\lambda\in\operatorname{Fun}(L)$ put
$D_k(\lambda)=\sum_a\lambda(1\otimes e_a)\mathsf{H}_{ak}$. Then \m{\lambda} and \m{\sum_kD_k(\lambda)\mu_k}
agree on the basis, hence on the fraction span, so
$\mu(x)=\sum_kD_k(\mu)\,\mu_k(x)$. It suffices to check that each triple \m{(\kappa_k,\mu_k,D_k)} is
admissible. The identity $\mathsf{R}(g)\,g(\mathsf{H})=\mathsf{H}$ says exactly that $\mu_k$ is
$\Gamma'$-equivariant on \m{1\otimes L} and that $D_k$ has~(D2). The entries of the $k$-th row
of~$\mathsf{H}^{-1}$ lie in \m{\Fil^{-\kappa_k}\BdR}; as $\pr_i(1\otimes t)\in\Ainf\otimes L$, this gives
\m{\mu_k\in\Per_{i,\kappa_k}(L)}. The entries of the $k$-th column of~$\mathsf{H}$ lie in
\m{\Fil^{\kappa_k}\BdR}, which gives~(D1).
\end{proof}

\begin{lemma}[Frobenius shifts the exponents]\label{l:shift}
Let $L$ be a lattice datum, $i$ a letter, \m{l\ge0} an integer, $o$ an exponent function and
$x\in\WC\otimes L$. If \m{x\in\wExp{L}{i^{p^{-l}}}{o}}, then
$\varphi^{l}x\in\wExp{L}{i}{\frobX{l}{0}{o}{\cdot}{\cdot}}$.
\end{lemma}
\leanline{\lean{phipowmemwithExponents}}
\begin{proof}
\m{\varphi^{l}x} lies in \m{\Mlocinf(L)} by Lemma~\ref{l:loc}(g) and~(h), and in the part at~$i$. Let
\m{\mu\in\Per_{i^{p^{-\nu}},u}(L)}. For \m{\nu\ge l},
$\mu(\varphi^{-\nu}\varphi^{l}x)=\mu(\varphi^{-(\nu-l)}x)\in\Fil^{o(\nu-l,u)}\BdR$, the letter being
\m{(i^{p^{-l}})^{p^{-(\nu-l)}}}. For \m{\nu<l} we need $\mu(\varphi^{a+1}x)\in\Fil^{0}\BdR$ with
\m{a=l-\nu-1\ge0}: this holds on \m{\Minf(L)} by definition, and it passes from \m{\cE_{\nu'}x} to~$x$
because $\mu(\varphi^{a+1}(\cE_{\nu'}x))=\cE_{\nu'+a+1}\,\mu(\varphi^{a+1}x)$ and the image
of~\m{\cE_{\nu'+a+1}} in~$\BdR$ is a unit of~$B^{+}_{\mathrm{dR}}$ (Input~\ref{i:div}, with index~$0$).
\end{proof}

\begin{proof}[Proof of Theorem B]
Write \m{x=\varphi^{l}x_0} and \m{j=i^{p^{-l}}}; then \m{x\in(\WC\otimes L)_i} if and only if
\m{x_0\in(\WC\otimes L)_j}.

\emph{From the Kisin module to the exponents.} Let \m{x_0\in\Minf(L)} lie in the part at~$j$. Then
\m{\tau^{a}x_0\in\Ainf\otimes L} for some~$a$, and for \m{\mu\in\Per_{j^{p^{-\nu}},u}(L)} Lemma~\ref{l:exp}(c)
gives $\varphi^{-\nu}(\tau)^{a}\mu(\varphi^{-\nu}x_0)\in\Fil^{-u}\BdR$, with \m{\varphi^{-\nu}(\tau)} a unit
of~$B^{+}_{\mathrm{dR}}$. So \m{x_0\in\wExp{L}{j}{\std_0}}, and Lemma~\ref{l:shift} gives
$x\in\wExp{L}{i}{\frobX{l}{0}{\std_0}{\cdot}{\cdot}}=\wExp{L}{i}{\std_l}$.

\emph{From the exponents to the Kisin module.} Let \m{x\in\wExp{L}{i}{\std_l}}. We must show that $x_0$ is
killed into \m{\Ainf\otimes L} by a power of~$\tau$, and that $\mu(\varphi^{a+1}x_0)\in\Fil^{0}\BdR$ for
\m{a\ge0}, \m{\mu\in\Per(L)}.

\emph{First.} We show that every \m{x'\in\wExp{L}{i}{\std_l}} is multiplied into \m{\Ainf\otimes L} by a power
of~\m{\varphi^{l}(\tau)}; applying~$\varphi^{-l}$ then gives the claim for~$x_0$. This goes by induction on the
localization. On \m{\Minf(L)} it holds because $\tau$ divides \m{\varphi^{l}(\tau)}. Suppose it holds for
\m{\cE_\nu x'}, where \m{x'\in\wExp{L}{i}{\std_l}}: so \m{c\,\cE_\nu x'\in\Ainf\otimes L} with~$c$ a power of
\m{\varphi^{l}(\tau)}. If \m{\nu<l}, then \m{\cE_\nu} divides \m{\varphi^{l}(\tau)} (Input~\ref{i:div}), and a
higher power works for~$x'$. If \m{\nu\ge l}, write $c\,\cE_\nu x'=\sum_kb_k(1\otimes e_k)$ in a
$\Zp$-basis of~$L$ and let $\lambda_k\in\operatorname{Fun}(L)$ extend the $k$-th coordinate. Then
\[\varphi^{-\nu}(b_k)=\lambda_k\bigl(\varphi^{-\nu}(c\,\cE_\nu x')\bigr)
=\varphi^{-\nu}(c)\,\varphi^{-\nu}(\cE_\nu)\,\lambda_k(\varphi^{-\nu}x').\]
We claim $\lambda_k(\varphi^{-\nu}x')\in\Fil^{0}\BdR$. By Lemma~\ref{l:basis}, at the letter
\m{i^{p^{-\nu}}} and for the element \m{\varphi^{-\nu}x'} of the fraction span, it suffices that
$D(\lambda_k)\mu'(\varphi^{-\nu}x')\in\Fil^{0}\BdR$ for admissible \m{(\kappa,\mu',D)}: by~(D1),
$D(\lambda_k)\in\Fil^{\kappa}\BdR$, and by Lemma~\ref{l:any} with \m{M=\Fil^{-\kappa}\BdR},
$\mu'(\varphi^{-\nu}x')\in\Fil^{-\kappa}\BdR$, since \m{\std_l(\nu,v)=-v\ge-\kappa} for \m{v\le\kappa}. So
$\varphi^{-\nu}(b_k)\in\Fil^{1}\BdR$, that is \m{\cE_\nu\mid b_k}, and \m{c\,x'\in\Ainf\otimes L} because
\m{\WC\otimes L} has no $\Ainf$-torsion.

\emph{Second.} Let \m{\mu\in\Per(L)} and \m{a\ge0}; apply Lemma~\ref{l:basis} to~$\mu$, the element
\m{\varphi^{a+1}x_0} and the letter \m{j^{p^{a+1}}}, with \m{n=0}. For admissible \m{(\kappa,\mu',D)}:
\m{D(\mu)} is fixed by~$\Gamma'$ by~(D2), because $\mu$ is a period, so \m{D(\mu)\in\Fil^{0}\BdR}
(Input~\ref{i:tate}). And $\varphi^{a+1}x\in\wExp{L}{i^{p^{a+1}}}{\frobX{a+1}{0}{\std_l}{\cdot}{\cdot}}$ by
Lemma~\ref{l:shift}, with \m{\frobX{a+1}{0}{\std_l}{l}{v}=0} for every~$v$; so
$\mu'(\varphi^{a+1}x_0)=\mu'(\varphi^{-l}\varphi^{a+1}x)\in\Fil^{0}\BdR$ by Lemma~\ref{l:any} with
\m{M=\Fil^{0}\BdR}.
\end{proof}

\begin{lemma}[projections to a letter]\label{l:proj}
Let $L$ be a lattice datum, $x\in\WC\otimes L$, $i$ a letter and \m{l\ge0}.
\begin{enumerate}[label=\textup{(\alph*)},leftmargin=2.3em,itemsep=2pt]
\item If \m{\varphi^{-l}x\in\Minf(L)}, then \m{\pr_i(x)\in\wExp{L}{i}{\std_l}}.
\item If $x$ is a point over~$F_0$ of \m{\varphi^{l}\Minf(L)}, that is \m{\varphi^{-l}x\in\Minf(L)} and
\m{\mu(x)\in F_0} for all \m{\mu\in\Per(L)}, then \m{\pr_i(x)\in\locLP{l}{i}(L)}.
\end{enumerate}
\end{lemma}
\leanline{\lean{letterProjmemwithExponents}\sep \lean{letterProjmemofvalueFieldPoints}}
\begin{proof}
(a) $\pr_i$ is a combination, with coefficients in~$\Ainf$, of powers of \m{1\otimes\omega}, and
multiplication by~$\omega$ is an equivariant map \m{L\to L}; so $\varphi^{-l}\pr_i(x)\in\Minf(L)$
(Input~\ref{i:forms}(b)), and Theorem~B applies. (b) By (a), and because \m{\mu\circ\pr_i\in\Per(L)}.
\end{proof}

\sect{Exponents add: the proof of Theorem C}{s:add}
\begin{lemma}[pairings]\label{l:pair}
In Setting~\ref{set:lattice}, let $A$, $L$, $C$ be lattice data and \m{b\colon A\times L\to C} an
$\cO$-bilinear map.
\begin{enumerate}[label=\textup{(\alph*)},leftmargin=2.3em,itemsep=2pt]
\item For a letter~$i$, $x\in\WC\otimes L$ and $y\in\WC\otimes A$:
$\wPair{b}{\pr_i(y)}{x}=\pr_i\bigl(\wPair{b}{y}{x}\bigr)$.
\item Let $T$ be a multiplicative subset of~$\Ainf$ and let $N_A$, $N_L$, $N_C$ be $\Ainf$-submodules of
\m{\WC\otimes A}, \m{\WC\otimes L}, \m{\WC\otimes C}. If some element of~$T$ multiplies $y$ into~$N_A$, some
element of~$T$ multiplies $x$ into~$N_L$, and \m{\wPair{b}{y'}{x'}\in N_C} for all \m{y'\in N_A},
\m{x'\in N_L}, then some element of~$T$ multiplies \m{\wPair{b}{y}{x}} into~$N_C$.
\end{enumerate}
\end{lemma}
\leanline{\lean{wittPairingletterProjleft}\sep \lean{wittPairingmemsaturateBy}}
\begin{proof}
(a) $b$ is $\cO$-bilinear, so its extension carries \m{1\otimes\omega} in the first variable to
\m{1\otimes\omega} on~$C$. (b) By bilinearity.
\end{proof}

\begin{proof}[Proof of Theorem C]
\emph{A criterion.} Let \m{\mu\in\Per(C)}, \m{m\in\mathbb{Z}\times\Gamma}, a letter~$j$, \m{u\in\mathbb{Z}} and
\m{n\in\mathbb{Z}}; let $x'$ lie in the fraction span of~$L$ and \m{y'\in\WC\otimes A}. We claim that
$\mu(m\cdot\wPair{b}{y'}{x'})\in\Fil^{n}\BdR$ provided that
\begin{equation}\label{e:crit}
\lambda(my')\cdot\mu'(mx')\in\Fil^{n}\BdR
\end{equation}
for all \m{\kappa\in\mathbb{Z}}, all \m{\mu'\in\Per_{j,\kappa}(L)} and all \m{\lambda\in\Per(A)} such that
\m{\lambda\in\Per_{j,u-\kappa}(A)} if \m{\mu\in\Per_{j,u}(C)}. Indeed, the equivariance of~$b$ gives
$m\cdot\wPair{b}{y'}{x'}=\wPair{b}{my'}{mx'}$. Apply Lemma~\ref{l:basis} to the lattice datum~$L$, the
element \m{mx'} and the functional \m{\mu\circ\wPair{b}{my'}{\cdot}}. For an admissible triple
\m{(\kappa,\mu',D)} put $\lambda(a)=D\bigl(\mu\circ\wPair{b}{a}{\cdot}\bigr)$ for \m{a\in\WC\otimes A}. This
is an $\Ainf$-linear map; it is a period of~$A$ by~(D2) and the equivariance of $b$ and~$\mu$; and if
\m{\mu\in\Per_{j,u}(C)}, then by Lemma~\ref{l:pair}(a) the functional
$\mu\circ\wPair{b}{\pr_j(1\otimes a)}{\cdot}$ takes values in \m{\Fil^{-u}\BdR} on \m{1\otimes L}, so
$\lambda(\pr_j(1\otimes a))\in\Fil^{-u+\kappa}\BdR$ by~(D1): $\lambda$ has type \m{u-\kappa} at~$j$.
So~\eqref{e:crit} is what Lemma~\ref{l:basis} asks.

\emph{The localized module.} By Lemma~\ref{l:pair}(b) for the multiplicative set generated by the~$\cE_\nu$,
then for the powers of~$\tau$, it suffices to see that \m{\wPair{b}{y'}{x'}\in\Minf(C)} for
\m{y'\in\Minf(A)}, \m{x'\in\Minf(L)}. The pairing maps $(\Ainf\otimes A)\times(\Ainf\otimes L)$ into
\m{\Ainf\otimes C}. And for \m{\mu\in\Per(C)}, \m{a\ge0}, the criterion with \m{m=\varphi^{a+1}} and
\m{n=0} applies: \m{\lambda(\varphi^{a+1}y')} and \m{\mu'(\varphi^{a+1}x')} lie in \m{\Fil^{0}\BdR} by the
definition of the Kisin modules.

\emph{The letter.} \m{y=\pr_i(y)}, so $\wPair{b}{y}{x}=\pr_i(\wPair{b}{y}{x})$ by Lemma~\ref{l:pair}(a).

\emph{The exponents.} Let \m{\nu\ge0}, \m{j=i^{p^{-\nu}}}, \m{v\in U_C(j)} and \m{\mu\in\Per_{j,v}(C)}. By the
criterion with \m{m=\varphi^{-\nu}}, \m{u=v} and \m{n=o_C(\nu,v)}, it suffices that
$\lambda(\varphi^{-\nu}y)\,\mu'(\varphi^{-\nu}x)\in\Fil^{o_C(\nu,v)}\BdR$ for \m{\mu'\in\Per_{j,\kappa}(L)}
and \m{\lambda\in\Per_{j,v-\kappa}(A)}. Apply Lemma~\ref{l:any} to~$y$ and~$\lambda$, with the module
$M_1=\{a\in\BdR: a\,\mu'(\varphi^{-\nu}x)\in\Fil^{o_C(\nu,v)}\BdR\}$: we must show
$\Fil^{o_A(\nu,u)}\BdR\subseteq M_1$ for \m{u\in U_A(j)}, \m{u\le v-\kappa}. Fix such~$u$ and
\m{a\in\Fil^{o_A(\nu,u)}\BdR}, and apply Lemma~\ref{l:any} to~$x$ and~$\mu'$, with
$M_2=\{c\in\BdR: a\,c\in\Fil^{o_C(\nu,v)}\BdR\}$: for \m{u'\in U_L(j)} with \m{u'\le\kappa} we have
\m{u+u'\le v}, so $o_C(\nu,v)\le o(\nu,u')+o_A(\nu,u)$ and
$a\cdot\Fil^{o(\nu,u')}\BdR\subseteq\Fil^{o_C(\nu,v)}\BdR$.
\end{proof}

\sect{The lifting ideal and cokernel killers}{s:kill}
Throughout, the setting is Setting~\ref{set:lattice} and $S$ is a side over~$\cO$.

\begin{lemma}[counting in two windows]\label{l:count}
Let \m{n,p,d,m,H\ge0} be integers and $c$ an integer. If \m{n\le p\,m} and \m{|c|+n\le H}, then
\[\sum_{a\in\twoWin{n}{c}}\bigl(v_p(a)+d\bigr)\ \le\ (4m+2)\,\lfloor\log_pH\rfloor+(4n+2)\,d .\]
\end{lemma}
\leanline{\lean{sumtwoWindowsle}\sep here $p$ is any natural number, \m{v_p(0)=0}, and \m{\lfloor\log_pH\rfloor} is read as~$0$ when \m{p\le1} or \m{H=0}}
\begin{proof}
The set \m{\twoWin{n}{c}} has at most \m{4n+2} elements, which accounts for the terms~$d$. A non-zero
\m{a\in\twoWin{n}{c}} has \m{|a|\le H}, so \m{v_p(a)\le\lfloor\log_pH\rfloor}; and \m{v_p(a)\ne0} only if
\m{a=pa'}, where $a'$ lies in \m{\twoWin{m}{\lfloor c/p\rfloor}}, a set with at most \m{4m+2} elements.
\end{proof}

\begin{lemma}[cyclotomic twists]\label{l:chi}
Let $L$ be a lattice datum, \m{c\ge0} an integer, $i$ a letter and \m{x\in\wExp{L}{i}{\std_0}}. If
\m{g t=\chi(g)^{c}t} for all \m{g\in\Gamma} and \m{t\in L}, then \m{\tau^{c}x\in\Ainf\otimes L}.
\end{lemma}
\leanline{\lean{taupowsmulmemainfLatticeofchi}}
\begin{proof}
By Theorem~B with \m{l=0}, \m{x\in\Minf(L)}; by the frame theorem \m{\tau^{h_L}x\in\Ainf\otimes L}; and
\m{h_L\le c} because all jumps of~$L$ are~$-c$ (Input~\ref{i:ext}).
\end{proof}

\begin{lemma}[the lifting ideal: formal properties]\label{l:ideal}
Let \m{r,l,n,c,a\ge0} be integers, $i$ a letter and \m{s\in\Ainf}.
\begin{enumerate}[label=\textup{(\alph*)},leftmargin=2.3em,itemsep=2pt]
\item \m{s\in\liftId{S}{r}{l}{i}} if and only if there is $t\in\Ainf$ such that for every $c'\in C_S$
some $x\in\wExp{\extL{S}}{i}{\std_l}$ has $(1\otimes\projC{S})(x)=(s-rt)\cdot\constAt{S}{i}{c'}$.
\item $p^{\kappa}\in\liftId{S}{p^{\kappa}}{l}{i}$ for every integer \m{\kappa\ge0}.
\item $\varphi^{l}\,\constAt{S}{i}{c'}=\constAt{S}{i^{p^{l}}}{c'}$ for $c'\in C_S$.
\item If \m{s\in\liftId{S}{r}{0}{i}} and \m{u\le n} for every \m{u\in U_{Y_S}(i)}, then
$s\,\cE^{n}\in\liftId{S}{r}{1}{i}$.
\item If $p^{c}\in\liftId{S}{p^{c+1}}{l}{i}$, then $\liftId{S}{p^{c+1}}{l}{i}\subseteq\liftId{S}{0}{l}{i}$.
\item $\constAt{S}{i}{c'}\in\locLP{l}{i}(\trivL{S})$ for $c'\in C_S$.
\item $\varphi^{l}(s)\in\liftId{S}{r}{l}{i^{p^{l}}}$ if and only if \m{s\in\liftId{S}{r}{0}{i}}.
\item $p^{c}z^{pa}\in\liftId{S}{r}{1}{i^{p}}$ if and only if $p^{c}z^{a}\in\liftId{S}{r}{0}{i}$.
\end{enumerate}
\end{lemma}
\leanline{\lean{memliftIdealiff}\sep \lean{ppowmemliftIdealself}\sep \lean{phipowsmulconstantAt}\sep
\lean{mulEpowmemliftIdealone}\sep \lean{liftIdealleofppowmem}\sep \lean{constantAtmemlocLetterPart}\sep
\lean{phipowmemliftIdealiff}\sep \lean{ppowmulzpowmemliftIdealiff}}
\begin{proof}
(a) \m{\Ainf\otimes C_S} is spanned over~$\Ainf$ by the \m{1\otimes c'}, and the image of
\m{\wExp{\extL{S}}{i}{\std_l}} is an $\Ainf$-module. (b) It is the generator. (c) $\varphi$ maps the
Teichmüller lift of~$i$ to that of~$i^{p}$, and fixes \m{1\otimes\omega} and \m{1\otimes c'}.

(d) Let $x$ lift \m{(s-rt)\constAt{S}{i}{c'}} as in~(a), with \m{l=0}. By Lemma~\ref{l:exp}(f) with
\m{\nu_0=0}, \m{\cE^{n}x} has the exponents \m{n-u\ge0=\std_1(0,u)} at \m{\nu=0} and \m{-u=\std_1(\nu,u)} at
\m{\nu\ge1}, for \m{u\in U_{Y_S}(i)}; it lifts $(s\cE^{n}-rt\cE^{n})\constAt{S}{i}{c'}$.

(e) Put \m{J=\liftId{S}{0}{l}{i}}; then $\liftId{S}{p^{c+1}}{l}{i}=(p^{c+1})+J$. The hypothesis says
$(p^{c})\subseteq p\cdot(p^{c})+J$. As $\Ainf$ is $p$-adically complete, $p$ lies in its Jacobson radical,
and Nakayama's lemma gives \m{(p^{c})\subseteq J}.

(f) The element \m{1\otimes c'} lies in \m{\Ainf\otimes C_S}, is fixed by~$\varphi$, and every period takes a
$\Gamma'$-fixed value on it, which lies in \m{\Fil^{0}\BdR\cap F_0} (Input~\ref{i:tate}); so it is a point
over~$F_0$ of \m{\varphi^{l}\Minf(\trivL{S})}, and Lemma~\ref{l:proj}(b) applies.

(g) By (a), Theorem~B and~(c): $x\in\wExp{\extL{S}}{i^{p^{l}}}{\std_l}$ if and only if
$\varphi^{-l}x\in\wExp{\extL{S}}{i}{\std_0}$, and $\varphi^{l}$ is a bijection that commutes with
\m{1\otimes\projC{S}} and fixes the integer~$r$. (h) is (g) for \m{l=1}, since \m{\varphi(z)=z^{p}}.
\end{proof}

\begin{lemma}[lifts over~$\fS$]\label{l:liftS}
Let \m{r,l\ge0} be integers, $i$ a letter and \m{f\in\fS} with \m{f\in\liftId{S}{r}{l}{i}}. Then there is
\m{t\in\fS} such that for every \m{y\in\locLP{l}{i}(\trivL{S})} some \m{x\in\locLP{l}{i}(\extL{S})} has
$(1\otimes\projC{S})(x)=(f-rt)\,y$.
\end{lemma}
\leanline{\lean{existsliftofmemliftIdeal}}
\begin{proof}
Write \m{f=rt_0+s} with \m{t_0\in\Ainf} and $s$ in the colon ideal of Definition~\ref{d:lift}(ii), and put
\m{t=\operatorname{co}_0(t_0)}. Let \m{y\in\locLP{l}{i}(\trivL{S})}. As \m{U_{\trivL{S}}(j)=\{0\}}, $y$ has
the exponents~$\std_0$, so \m{y\in\Ainf\otimes C_S} by Lemma~\ref{l:chi} with \m{c=0}, and
\m{y=\pr_i(y)}. Hence \m{sy=(1\otimes\projC{S})(x_1)} with $x_1\in\wExp{\extL{S}}{i}{\std_l}$. Put
\m{x=\ret_{\extL{S}}(x_1)}; it lies in \m{\locLP{l}{i}(\extL{S})} by Lemma~\ref{l:retprop}(d), and by
Lemma~\ref{l:retprop}(b) and~(a),
$(1\otimes\projC{S})(x)=\ret_{\trivL{S}}(sy)=\operatorname{co}_0(s)\,y=(f-rt)\,y$.
\end{proof}

\begin{lemma}[differentiating in the lifting ideal]\label{l:der}
Let \m{r\ge0} be an integer, $i$ a letter, \m{f\in\fS}, \m{s,\alpha,\beta\in\Ainf} and \m{g\in\Gamma'}. If
\m{f\in\liftId{S}{r}{1}{i}} and $z\tau\,\cE\,(s+\beta f)=\alpha\,(g f-f)$, then
\m{s\in\liftId{S}{r}{1}{i}}.
\end{lemma}
\leanline{\lean{memliftIdealofderivation}}
\begin{proof}
Take \m{t\in\fS} from Lemma~\ref{l:liftS} with \m{l=1}, and \m{v\in\Ainf} with \m{gt-t=z\tau\cE v}
(Input~\ref{i:gal}). Let $c'\in C_S$, \m{y=\constAt{S}{i}{c'}}, which lies in \m{\locLP{1}{i}(\trivL{S})}
by Lemma~\ref{l:ideal}(f) and is fixed by~$g$, and let \m{x\in\locLP{1}{i}(\extL{S})} with
\m{(1\otimes\projC{S})(x)=(f-rt)y}. Then
$(1\otimes\projC{S})(gx-x)=\bigl((gf-f)-r\,z\tau\cE v\bigr)y$, and
\[x'=\alpha\,\cE^{-1}(z\tau)^{-1}(gx-x)-\beta x\qquad\text{has}\qquad
(1\otimes\projC{S})(x')=\bigl(s-r(\alpha v-\beta t)\bigr)\,y .\]
By Lemma~\ref{l:ideal}(a) it remains to see that $x'\in\wExp{\extL{S}}{i}{\std_1}$. The element \m{gx-x}
has the exponents~$\std_1$ (Lemma~\ref{l:exp}(g)). At \m{\nu=0} it has more: for a period~$\mu$ of type~$u$
at~$i$, \m{V=\mu(x)} lies in \m{F_0\cap\Fil^{0}\BdR}, and \m{gV-V\in\Fil^{1}\BdR} by Input~\ref{i:val}(b),
since \m{\cE\mid gs'-s'} for \m{s'\in\fS}. So $gx-x\in\wExp{\extL{S}}{i}{\frobX{1}{1}{\std_0}{\cdot}{\cdot}}$
by Lemma~\ref{l:exp}(a), and so does \m{(z\tau)^{-1}(gx-x)} by Lemma~\ref{l:ztau}(b). Lemma~\ref{l:exp}(f)
with \m{\nu_0=0}, \m{n=-1} then gives the exponents~$\std_1$ for \m{\cE^{-1}(z\tau)^{-1}(gx-x)}.
\end{proof}

\begin{lemma}[from level one to level zero, and removing~$\cE$]\label{l:down}
Let \m{r,a,n,c,\kappa,b\ge0} be integers and $i$ a letter.
\begin{enumerate}[label=\textup{(\alph*)},leftmargin=2.3em,itemsep=2pt]
\item If \m{pa\le n}, $p^{c}z^{n}\in\liftId{S}{r}{1}{i^{p}}$ and
\m{c+v_p(n!)\le\kappa+v_p((pa)!)}, then $p^{\kappa}z^{a}\in\liftId{S}{r}{0}{i}$.
\item If $p^{c}z^{a}\cE^{n}\in\liftId{S}{r}{1}{i}$, \m{c+\fd\,n+v_p(n!)\le b} and \m{a+n\le\kappa}, then
$p^{b}z^{\kappa}\in\liftId{S}{r}{1}{i}$.
\end{enumerate}
\end{lemma}
\leanline{\lean{memliftIdealzeroofmemone}\sep \lean{memliftIdealofmulEpow}}
\begin{proof}
(a) By induction on \m{n\ge pa}. If \m{n=pa}, then \m{c\le\kappa} and Lemma~\ref{l:ideal}(h) applies. For
the step from $n$ to \m{n+1}, let $f=p^{c}z^{n+1}\in\liftId{S}{r}{1}{i^{p}}$ with
$c+v_p((n+1)!)\le\kappa+v_p((pa)!)$. Choose \m{g\in\Gamma'} with \m{gz=[\varepsilon]z}, and
\m{\alpha,\beta} with $\tau\cE(p^{v_p(n+1)}+\beta z)=\alpha([\varepsilon]^{n+1}-1)$
(Input~\ref{i:lte}). Then $gf-f=p^{c}([\varepsilon]^{n+1}-1)z^{n+1}$ and
\[z\tau\,\cE\,\bigl(p^{c+v_p(n+1)}z^{n}+\beta f\bigr)=\alpha\,(gf-f).\]
By Lemma~\ref{l:der}, $p^{c+v_p(n+1)}z^{n}\in\liftId{S}{r}{1}{i^{p}}$, and the induction hypothesis applies
to it, because $c+v_p(n+1)+v_p(n!)=c+v_p((n+1)!)$.

(b) By induction on~$n$; for \m{n=0} the element \m{p^{b}z^{\kappa}} is a multiple of \m{p^{c}z^{a}}. For the
step, let $f=p^{c}z^{a}\cE^{n+1}\in\liftId{S}{r}{1}{i}$. Take \m{g,w,y,X} from Theorem~A for
\m{(0,n+1)}, and write \m{yX-p^{\cost{0}{n+1}}=\cE b_0}, \m{gz=z(1+\tau\cE v)} (Input~\ref{i:gal}) and
$(1+\tau\cE v)^{a}=1+\tau\cE vt'$. Then
\[gf-f=p^{c}z^{a}\cE^{n+1}\bigl((1+\tau\cE vt')w^{n+1}-1\bigr)
=p^{c}z^{a}\cE^{n+1}\,\tau\,\bigl(zX+\cE vt'w^{n+1}\bigr),\]
and a direct computation gives
\[\begin{gathered}zy\cdot(gf-f)=z\tau\,\cE\,\bigl(s_1+\beta_1f\bigr),\\
s_1=p^{c+\cost{0}{n+1}}z^{a+1}\cE^{n},\qquad \beta_1=zb_0+yvt'w^{n+1}.\end{gathered}\]
By Lemma~\ref{l:der}, \m{s_1\in\liftId{S}{r}{1}{i}}. The induction hypothesis applies to it, with \m{a+1} in
place of~$a$ and \m{c'=c+\cost{0}{n+1}} in place of~$c$: indeed \m{a+1+n\le\kappa} and
\[c'+\fd n+v_p(n!)=c+v_p(n+1)+\fd+\fd n+v_p(n!)=c+\fd(n+1)+v_p((n+1)!)\le b .\]
\end{proof}

\begin{lemma}[a power of the coordinate, modulo~$p^{\kappa}$]\label{l:zpow}
Let \m{n\ge0} be an integer with \m{u\le n} for every letter~$j$ and every \m{u\in U_{Y_S}(j)}, and let
\m{\kappa\ge0} be an integer. Then there is an integer \m{a\ge0} with
$z^{a}\in\liftId{S}{p^{\kappa}}{0}{j}$ for every letter~$j$ and \m{a\,p\le e_{\fc}\,(n+\kappa)+a}.
\end{lemma}
\leanline{\lean{existszpowmemliftIdeal}}
\begin{proof}
Let $I$ be the intersection of the ideals \m{\liftId{S}{p^{\kappa}}{0}{j}} over the letters~$j$. It
contains~$p^{\kappa}$. It contains \m{\tau^{h}} with \m{h=h_{\extL{S}}}: for $c'\in C_S$,
$\tau^{h}(1\otimes(0,c'))\in\Minf(\extL{S})$ by the frame theorem, its projection to the letter~$j$ has the
exponents~$\std_0$ by Lemma~\ref{l:proj}(a), and it maps to \m{\tau^{h}\constAt{S}{j}{c'}}. Since $z$ is not
a unit and \m{p\nmid\tau}, $z$ lies in the radical of \m{(p,\tau)} (Input~\ref{i:rad}), hence in the radical
of~$I$. Let $a$ be the least integer with \m{z^{a}\in I}, and put \m{N=e_{\fc}(n+\kappa)+a} and
\m{a'=\lfloor N/p\rfloor}. It suffices to show \m{z^{a'}\in I}: then \m{a\le a'}, that is \m{ap\le N}.

Fix a letter~$i$. By Lemma~\ref{l:ideal}(d), $z^{a}\cE^{n}\in\liftId{S}{p^{\kappa}}{1}{i^{p}}$, and
\m{p^{\kappa}z^{a}} lies there too. Since \m{z^{e_{\fc}}\in(\cE,p)} (Input~\ref{i:mod}),
$z^{e_{\fc}(n+\kappa)}\in(\cE,p)^{n+\kappa}\subseteq(\cE^{n},p^{\kappa})$, and so
$z^{N}\in\liftId{S}{p^{\kappa}}{1}{i^{p}}$. Now Lemma~\ref{l:down}(a), with \m{c=0}, exponent~$0$ of~$p$ in
the conclusion, and \m{pa'\le N}, gives $z^{a'}\in\liftId{S}{p^{\kappa}}{0}{i}$: its numerical condition
\m{v_p(N!)\le v_p((pa')!)} holds because no multiple of~$p$ lies in \m{(pa',N]}.
\end{proof}

\begin{proof}[Proof of Theorem D]
Here \m{\cE=z+p}, \m{e_{\fc}=1} and \m{\fd=0}. Put \m{q=\kappa+\rho}. By Lemma~\ref{l:zpow} for
\m{n=(p-3)q} and the modulus \m{p^{q+1}} there is~$b$ with $z^{b}\in\liftId{S}{p^{q+1}}{0}{j}$ for all
letters~$j$ and \m{(p-1)\,b\le(p-2)\,q+1}. We claim that, for every integer \m{a\ge0},
\begin{equation}\label{e:mono}
p^{\max(q-a,0)}z^{a}\in\liftId{S}{p^{q+1}}{0}{i}.
\end{equation}
Put \m{a_1=\max(a,b)}. Then $z^{a_1}\in\liftId{S}{p^{q+1}}{0}{i}$, so
$z^{pa_1}\in\liftId{S}{p^{q+1}}{1}{i^{p}}$ by Lemma~\ref{l:ideal}(h), and Lemma~\ref{l:down}(a) with
\m{c=0} and \m{n=pa_1} gives~\eqref{e:mono} as soon as
$v_p((pa_1)!)\le\max(q-a,0)+v_p((pa)!)$. For \m{a\ge b} this is clear. For \m{a<b}, since
\m{v_p((pa)!)=a+v_p(a!)}, it follows from \m{b+v_p(b!)\le q}; this is clear for \m{b=0}, and \m{(p-1)v_p(b!)\le b-1} for
\m{b\ge1}, so
\[(p-1)^{2}\bigl(b+v_p(b!)\bigr)\le(p-1)(pb-1)\le p(p-2)q+1<(p-1)^{2}(q+1).\]
For \m{a=0}, \eqref{e:mono} reads $p^{q}\in\liftId{S}{p^{q+1}}{0}{i}$, so
$\liftId{S}{p^{q+1}}{0}{i}\subseteq\liftId{S}{0}{0}{i}$ by Lemma~\ref{l:ideal}(e). Finally
\[p^{\kappa}\cE^{\rho}=\sum_{a=0}^{\rho}\binom{\rho}{a}\,p^{\,q-a}z^{a},\]
and each term lies in \m{\liftId{S}{p^{q+1}}{0}{i}} by~\eqref{e:mono}.
\end{proof}

\sect{Where the results are used}{s:used}
Parts 2 and~3 of the notes on this unit (\emph{Cup products of crystalline classes vanish up to a power of~$p$:
the walk along Frobenius orbits}, and the note on few low weights and tame types) use everything above. Proposition~\ref{p:succ} and Lemma~\ref{l:deep} are
iterated in Part~2 over finite sets of steps \m{(\nu,n)}, the sum of whose costs is the price of a deep point;
Theorem~C pairs a point of a side with a point of the dual side; Theorem~D feeds the walk
along Frobenius orbits that bounds cup products of crystalline classes; and Lemma~\ref{l:count} and the lifting
ideal with its differentiation lemma return in Part~3, in the estimates for cut windows and in the cokernel killer
for tame types. Outside this unit nothing of this note is
used directly.

\sect{Notes on the formal text}{s:notes}
\emph{Arguments as the text gives them.} The arguments of the declarations were read from the Lean text and from
the census of the workspace; they were not printed by Lean. The census reads a bound variable of the definition of
the exponent module as an explicit argument (an index~$\nu$): the definition has the two explicit arguments
$i$ and~$o$ and the implicit~$L$, as every use shows, and the note writes all three. In the definition of the
constants the census counts the element~$c$ to which the linear map is applied as an argument, and the note
writes it so. In the pairing \m{\wPair{b}{y}{x}} the census reads the arrows of the bilinear map as two
arguments, which are $y$ and~$x$. The declarations of the three files take the prime, the group and the chart as
instance arguments through the types of their other arguments, also where the census does not list them.

\emph{What the formal statements contain beyond the natural ones.} A lattice datum records a jump set at every
embedding of~$\cO$, of which only the union over a letter, \m{U_L(i)}, is used here. A side carries an action of~$\Gamma$ on its module~$C_S$ that is never used: the extension and the trivial datum give~$C_S$ the trivial action. The sets \m{U_L(i)} are
bounds: they contain the jumps and may be larger; this is why Lemma~\ref{l:any} allows a type~$u$ outside the
set. The localized modules invert all~$\cE_\nu$, \m{\nu\ge0}, at once, and the exponent at each of them is a
separate condition. In Definition~\ref{d:win}(iv) and Lemma~\ref{l:count} the letter~$p$ is an arbitrary natural
number. Theorem~D assumes \m{p\ge5}; we did not examine whether the formal proof needs this, and the estimate
given above uses only \m{p\ge3}. Theorem~A is stated for all \m{n\ne0}, negative ones included; the cost depends
only on \m{v_p(n)}.

\emph{Assemblies and small results.} Lemma~\ref{l:loc}, Lemma~\ref{l:exp}, Lemma~\ref{l:retprop},
Lemma~\ref{l:small}, Lemma~\ref{l:pair} and Lemma~\ref{l:ideal} collect several declarations each; the Lean
names are listed under them. The identity \m{\mu(\varphi^{0}x)=\mu(x)} is a declaration of its own in the formal
text and is not stated here.

\emph{The background.} The Inputs are stated as their Lean texts read; the frame theorem (Input~\ref{i:frame})
has further clauses that this note does not use (a spanning statement up to a power of~$p$, used in Part~2). The
period ring is Colmez's $B_{\max}$ under the name \emph{Bcris}. The second reader of the workspace compared the
$p$-adic Hodge inputs with the literature: for Input~\ref{i:tate} she finds a weaker form of Tate's theorems,
for Input~\ref{i:hodge} a classical statement for field coefficients whose formulation here is not located. We
did not confirm a place in the literature for the other Inputs.

\subhead{How this note was written.} From the reader's dossier of the unit: the statements of all members and
the structured proofs of the declarations treated here. The three Lean files and the first part of the fourth
were read in full, together with the Lean definitions of the background that the Settings restate (the chart, the
Kummer chart, lattice data, periods, the Kisin modules, the frame theorem). The proof of Theorem~A was read line
by line in the Lean text, because the reader had marked its bookkeeping as least certain; the account above of
where the three parts of the cost come from is ours, and it agrees with the reader's steps. The proofs of
Theorem~B, of Lemma~\ref{l:basis} and of Theorem~C were also rebuilt from the Lean terms. The arithmetic at the
end of the proof of Theorem~D is ours and replaces an automated step of the formal proof. The program check of
the unit's dossier covered \UnitSteps{} leaf steps of \UnitProofs{} proofs.

\head{Appendix A. Concordance with the Lean artifact}
{\fontsize{8.6}{11.6}\selectfont\renewcommand{\leanfont}{\ttfamily\fontsize{7.7}{10}\selectfont\color{leancol}}
\setlength{\LTleft}{0pt}\setlength{\LTright}{\fill}\setlength{\tabcolsep}{0pt}
\begin{longtable}{@{}p{30mm}p{56mm}p{43mm}r@{}}
\textbf{In this note} & \textbf{Lean declaration} & \textbf{File} & \textbf{Lines}\\\midrule\endhead
\ConcordanceRows
\end{longtable}}

Two statements as the artifact has them:
\quoted{existsdepthstep}
\quoted{memwithExponentsstdExponentiff}

\emph{Legend.} The formal names of the objects and the letters of this note: the chart~\texttt{k} is~$\fc$, its
ring \texttt{k.S} is~$\fS$, \texttt{k.E} is~$\cE$ and \texttt{k.frobE a} is~$\cE_\nu$; \texttt{k.f} is~$h$,
\texttt{k.e} is~$e_{\fc}$ and \texttt{differentExp} is~$\fd$; \texttt{galMonoid} is~$\Gamma'$ acting through
\m{(0,g)}; \texttt{zTau} is~\m{z\tau}; \texttt{hodgeFil n} is \m{\Fil^{n}\BdR}; \texttt{valueField a}
is~$F_\nu$ and \texttt{localFieldOf} is~$F$; \texttt{coord₀} is \m{\operatorname{co}_0}; \texttt{forms} and
\texttt{formsOfType} are \m{\Per(L)} and \m{\Per_{i,u}(L)}; \texttt{valueAt a μ x} is
\m{\mu(\varphi^{-\nu}x)}; \texttt{kisinAinf}, \texttt{kisinLocAinf}, \texttt{kisinLoc}, \texttt{kisinModule} are
\m{\Minf(L)}, \m{\Mlocinf(L)}, \m{\Mloc(L)}, \m{\mathfrak{M}(L)}; \texttt{stepBack i a} is
\m{i^{p^{-\nu}}}; \texttt{stdExponent l} is~$\std_l$; \texttt{ζ} of the coefficients is~$\omega$.

\head{Notation}
\begingroup\fontsize{8.6}{10.9}\selectfont\raggedright
\setlength{\columnsep}{9mm}\setlength{\parskip}{0pt}
\begin{multicols}{2}
\nt{$\ell$}{the prime of Chapter~VIII of the handout; it is~$p$ here}
\nt{$p$}{a prime number; in Definition~\ref{d:win}(iv) and Lemma~\ref{l:count} any natural number}
\nt{$\Qp$, $\Zp$, $\bar{\mathbb{Q}}_p$, $C$, $\mathcal{O}_C$}{the $p$-adic numbers and integers, an algebraic closure, its completion and the integers of the latter}
\nt{$\GQp$}{the absolute Galois group of~$\Qp$}
\nt{$\mathcal{O}_C^\flat$, $C^\flat$, $\Ainf$, $\WC$}{the tilt, its fraction field, their rings of Witt vectors}
\nt{$\varphi$, $\theta$, $\theta_\nu$, $[\,\cdot\,]$}{Frobenius, Fontaine's map, \m{\theta\circ\varphi^{-\nu}}, the Teichmüller lift}
\nt{$\varepsilon$, $\tau$, $\xi$}{a system of $p$-power roots of unity in the tilt, \m{[\varepsilon^{1/p}]-1}, a generator of \m{\ker\theta}}
\nt{$B^{+}_{\mathrm{dR}}$, $\BdR$, $\Fil^{n}\BdR$}{the de Rham period rings and the filtration}
\nt{$B_{\max}$, $A_{\max}$, $B^{+}$}{Colmez's period rings; \m{B^{+}=A_{\max}[1/p]}}
\nt{$\Gamma$, $j_\Gamma$, $\chi$}{a group over~$\GQp$, its map, the cyclotomic character}
\nt{$\fc$, $\Gamma'$}{the Kummer chart and its subgroup}
\nt{$h$, $\WW$, $\mathbb{F}_{p^h}$, $X$, $\mathsf{A}$}{the residue degree, \m{W(\mathbb{F}_{p^h})}, the field, the variable and the unit series of the chart}
\nt{$\fS$, $z$, $\cE$, $\cE_\nu$}{the ring of the chart, its coordinate, the Eisenstein element, \m{\varphi^\nu(\cE)}}
\nt{$e_{\fc}$, $\fd$}{the index of the chart; $0$ or \m{v_p(e_{\fc})+1}}
\nt{$F$, $F_\nu$, $\fS_B$}{the field fixed by~$\Gamma'$, the value fields, the ring of Setting~\ref{set:rings}}
\nt{$\operatorname{co}_0$, $\operatorname{co}_a$}{coordinates \m{\Ainf\to\fS}}
\nt{$\cO$, $\omega$}{the coefficient ring and its root of unity}
\nt{$i$, $j$}{letters, elements of \m{\mathbb{F}_{p^h}^\times}}
\nt{$L$, $L'$, $A$, $C$, $T$, $T'$, $T''$}{lattice data and lattices}
\nt{$U_L(i)$, $h_L$}{the jump set at a letter, the jump bound}
\nt{$\pr_i$, \m{(\WC\otimes L)_i}}{the projector to the letter~$i$ and its image}
\nt{$\Per(L)$, $\Per_{i,u}(L)$, $\operatorname{Fun}(L)$}{periods, periods of type~$u$ at~$i$, all $\Ainf$-linear functionals}
\nt{$\mu$, $\mu'$, $\mu_k$, $\lambda$, $\lambda'$, $\lambda_k$}{periods and functionals; in Setting~\ref{set:lattice} $\lambda$ is a semilinear map}
\nt{$\Minf(L)$, $\Mlocinf(L)$, $\Mloc(L)$, $\mathfrak{M}(L)$}{Kisin modules over~$\Ainf$, localized, over~$\fS$ localized, over~$\fS$}
\nt{$\std_l$}{the standard exponents of level~$l$}
\nt{$o$, $o'$, $o_A$, $o_C$, $o^{\mathrm{lo}}$, $o^{\mathrm{hi}}$}{exponent functions}
\nt{\m{\cost{\nu}{n}}}{the cost of a step}
\nt{\m{\wExp{L}{i}{o}}, \m{\locLP{l}{i}(L)}}{elements with exponents; those in \m{\Mloc(L)} with standard exponents}
\nt{\m{\locKP{Q}}, \m{\invMod{Q}}}{the elements of \m{\Mloc(L)} in~$Q$; the invariants modulo~$Q$}
\nt{$\impE{o^{\mathrm{lo}}}{o^{\mathrm{hi}}}{\nu}{u}$, \m{\frobX{l}{n}{o}{\nu}{u}}}{improved and shifted exponents}
\nt{\m{\wPair{b}{y}{x}}}{the extension of a bilinear map}
\nt{$\ret_L$}{the retraction}
\nt{$S$, $Y_S$, $C_S$, $\alpha_S$}{a side and its parts}
\nt{\m{\extL{S}}, \m{\trivL{S}}, \m{\projC{S}}}{the extension, the trivial datum, the projection}
\nt{\m{Z^{1}_{\mathrm{f}}(T)}}{the crystalline cocycles}
\nt{$\JumpWin{\eta}{\varsigma}{\beta}{W}{U}$}{the set~$U$ lies in a low and a far window}
\nt{\m{\raisedE{n}{\nu}{u}}, \m{\twoWin{n}{c}}, \m{\NoMult{p}{n}{\beta}}}{exponents raised at~$\cE$; two windows; no multiple of~$p$ near}
\nt{\m{\constAt{S}{i}{c}}, \m{\liftId{S}{r}{l}{i}}, \m{\Kills{S}{i}{\kappa}{\rho}}}{constants, the lifting ideal, a cokernel killer}
\nt{$\nu$, $\nu_0$, $\nu'$, $l$}{indices of Frobenius translates and levels}
\nt{$u$, $u'$, $v$, $\kappa$, $\kappa_k$}{jumps and types; $\kappa$ is also an exponent of~$p$, $v$ also an element of~$\Ainf$ or of a field}
\nt{$n$, $m$, $a$, $a'$, $a_1$, $b$, $c$, $d$, $k$, $q$, $r$, $H$, $N$, $W$, $\eta$, $\beta$, $\rho$}{integers; $b$ is also a bilinear map, $c$ also an element, $r$ also a map or an element}
\nt{$\varsigma$}{a sign}
\nt{$g$, $\gamma$, $t_k$}{elements of~$\Gamma'$ or of~$\GQp$; $\gamma$ is also a map \m{\cO\to C}}
\nt{$w$, $w_0$, $\hat w$, $\mathsf{u}$}{units}
\nt{$x$, $x'$, $x_0$, $x_1$, $y$, $y'$, $y_0$, $X$, $s$, $s'$, $t$, $t_0$, $t'$, $f$, $G$, $P_0$, $\alpha$, $\beta$, $s_1$, $\beta_1$, $a_k$, $b_k$, $b_0$, $c_k$, $v_k$, $w_a$, $e_k$}{elements of rings and modules; $f$ is also a map, $t$ also an element of a lattice}
\nt{$V$, $V_0$, $\pi_\nu$, $\varrho$, $\sigma_\nu$}{elements of~$\BdR$, a uniformizer, a root of unity, a sum}
\nt{$D$, $D_k$, $\mathsf{H}$, $\mathsf{R}(g)$}{coordinates in a Hodge basis, the Hodge matrix, the matrix of~$g$}
\nt{$M$, $M_1$, $M_2$, $Q$, $N_A$, $N_L$, $N_C$, $I$, $J$}{submodules and ideals}
\nt{$c'$}{an element of~$C_S$ or of~$F$; in the proof of Lemma~\ref{l:down} an integer}
\end{multicols}
\endgroup

\end{document}
